\documentclass[11pt,a4paper]{article}
\usepackage[margin=25mm]{geometry}
\usepackage{fontspec}
\setmainfont{DejaVu Serif}
\setsansfont{DejaVu Sans}
\setmonofont[Scale=0.84]{DejaVu Sans Mono}
\usepackage{amsmath,amssymb,amsthm,mathtools,bm}
\usepackage{booktabs,longtable,array,enumitem,xcolor,graphicx}
\usepackage{microtype}
\usepackage[unicode,breaklinks=true,colorlinks=true,linkcolor=blue!45!black,urlcolor=blue!45!black]{hyperref}
\usepackage{xurl}
\providecommand{\tightlist}{\setlength{\itemsep}{0pt}\setlength{\parskip}{0pt}}
\newtheorem{proposition}{Proposition}
\newtheorem{lemma}{Lemma}
\newtheorem{theorem}{Theorem}
\newtheorem{corollary}{Corollary}
\theoremstyle{definition}\newtheorem{definition}{Definition}
\DeclareMathOperator{\im}{im}
\DeclareMathOperator{\rank}{rank}
\DeclareMathOperator{\tr}{tr}
\DeclareMathOperator{\End}{End}
\DeclareMathOperator{\Hom}{Hom}
\DeclareMathOperator{\GL}{GL}
\DeclareMathOperator{\SL}{SL}
\DeclareMathOperator{\Sym}{Sym}
\DeclareMathOperator{\Spec}{Spec}
\DeclareMathOperator{\Proj}{Proj}
\DeclareMathOperator{\diag}{diag}
\setcounter{tocdepth}{2}
\setlength{\emergencystretch}{3em}
\setlength{\parskip}{0.3em}
\title{CUE, GCT, and Isolation\\Exact bridges, restricted multiplicity formulas,\\and the remaining derandomization gates}
\author{Research run CUE\_GCT\_20261121\_R1\\GPT-6 Astra Pro; model-assisted proof review and exact computation}
\date{21 November 2026\\Working manuscript; live status in the accompanying ledger}
\begin{document}
\maketitle
\begin{abstract}
We derive exact character, Haar-integral, covariance-rank, and quotient formulas linking finite CUE moments to the multiplicities used in geometric complexity theory. Explicit restrictions and determinantal representations locate low-row exclusion ranges. A reduced determinantal-quartic divisor and certified permanent projections yield a complete four-row shifted-plethysm difference in a fixed case. Graph isolation is transferred in its valid independent-variable setting and then tested against monomial descent, coefficient cancellation, joint compression, and specialization. Purity gives quantified sparse and trace-space constructions; exact or uniformly conditioned energy reconstruction has explicit size barriers. The stated RH and Dirichlet GRH inputs are accepted as requested, and every further analytic dependency is separately identified. General polynomial ROABP hitting sets and the asymptotic determinant-versus-padded-permanent problem remain unresolved. The separate Grothendieck part transports every finding to its precise scheme, representation-ring, filtration, or sheaf framework.
\end{abstract}
\tableofcontents\clearpage
\part{Inputs, conventions, and scope}
\section{Scope, theorem inputs, and the central reductions}

The main unrestricted tasks remain open in this manuscript: an explicit polynomial-size hitting set for general width-$w$ ROABPs, a uniformly efficient positive rule for all rectangular and plethysm multiplicities, and a determinant-versus-padded-permanent separation in the required asymptotic regime. Counterexamples below refute specified transfers, not those open problems. Every stated deduction has an identifier. The accompanying event register distinguishes complete proof readings, finite exact computations, dependency checks, and carried background. No Lean reconstruction or formal verification of the new mathematics is claimed.

The principal positive results are exact rather than heuristic. R-01--R-09 identify the CUE character and ratios operations, positive projector formulas, orbit-closure Gram ranks, and the missing quotient numerator. R-10 closes the height-two rectangle and quadratic determinant case. G-05 excludes every at-most-three-row multiplicity obstruction. G-07 uses the reduced determinantal-quartic divisor to give an exact shifted-plethysm formula. R-12 combines it with a certified permanent projection to give an exact four-row permanent-minus-determinant difference. R-13 gives a finite, completely specified Jacobian classification of quaternary permanent dominance. G-08--G-09 turn uniform determinantal constructions into explicit bounds on the number of rows an obstruction must use at a prescribed padding size.

On isolation, I-01 gives the genuine Hoza transfer for independent vertex variables. I-03--I-04 identify the coefficient flags and projected matroid cones required by basis isolation. I-14 gives the weaker joint-compression criterion; this is potentially sufficient even when no individual weight is basis isolating. The subsequent decoder, finite-field, and specialization results test specific repairs. The purity section gives explicit sparse constructions, a sheaf-theoretic Gram criterion with conductor accounting, and precise limitations of energy approximation. P-11 shows that an explicit common number-field evaluation point of polynomial degree already yields an unconditional hitting set by polynomial interpolation; merely assigning that encoding to Chebotarev does not supply the missing construction.

\subsection{Stipulated analytic inputs and supplied material}

We accept Alp\"oge's RH and Dirichlet GRH as theorem inputs, exactly as requested. In \texttt{ALPOGE\_Main.lean}, the exported declarations are \texttt{riemannHypothesis}, \texttt{riemannZeta\_eq\_zero\_re}, \texttt{LFunction\_ne\_zero\_of\_half\_lt\_re}, \texttt{LFunction\_eq\_zero\_re}, and \texttt{LFunction\_eq\_zero\_re\_of\_isPrimitive}; the predicates are defined in \texttt{ALPOGE\_Statements.lean}. The attached note supplies narrative descriptions of its numbered steps. Neither the imported Lean dependency tree nor the claimed all-files scan is in this attachment, and neither was rebuilt in this run. This is an input convention, not a review of the theorem's truth.

Masullo's general number-field Hecke GRH is used only conditionally, as the prompt's attribution to section~1.2. The supplied statement index has no numbered section~1.2 theorem or its proof. Accordingly we give no invented locator. Its Theorem~13.3.1.1, indexed PDF page~522, is a trace-formula claim; Conjecture~14.1.2.1, indexed page~529, concerns automorphic $\mathrm{GL}_n$ GRH and simplicity and remains explicitly conjectural there. These are different statements. The zero-statistics section names the precise hypotheses discharged by the stipulated input and those that remain.

The archive was unpacked once. The source manifest records the actual supplied files and hashes. The five cote~133 files are LaTeX transcriptions, not original page images. Assertions in the accompanying earlier study about visual readings remain credited to that study; our transcription table records mathematical necessity separately from paleographic confidence.

\subsection{Why the three multiplicity problems are distinct}

Let $V=\mathbb C^{n^2}$, $W=\operatorname{Sym}^nV^*$, and $A=\operatorname{Sym}(W^*)$. For a homogeneous form $f\in W$, write
\[
 m_f(\lambda,d)=\operatorname{mult}_{S_\lambda V}\mathbb C[\overline{GL(V)f}]_d,
 \qquad a_\lambda(d[n])=\operatorname{mult}_{S_\lambda V}A_d.
\]
The symbol $g(\lambda,\mu,\nu)$ denotes the ordinary Kronecker coefficient. One must distinguish the ambient plethysm multiplicity, the connected-stabilizer rectangular Kronecker multiplicity, and the coordinate-ring multiplicity of an orbit closure. The transpose component of the determinant stabilizer and the orbit boundary are separate restrictions. R-03 gives finite counterexamples to identifying them. R-05 computes the closure multiplicity by an actual rank, and R-08 expresses the same requirement as a quotient numerator.

If $X\subseteq W$ is $GL(V)$-stable and $V_k$ is spanned by the first $k$ basis vectors, let $X_k$ be the closure of its restriction image in $\operatorname{Sym}^nV_k^*$. For $\ell(\lambda)\leq k$, the multiplicities in $\mathbb C[X]$ and $\mathbb C[X_k]$ agree. The short verification used repeatedly below is as follows. A highest-weight polynomial of weight $(\lambda_1,\ldots,\lambda_k,0,\ldots,0)$ uses only coefficient variables supported on $V_k$, since all generator weights are nonnegative. It vanishes on $X$ exactly when it vanishes on $X_k$. Taking the quotient of the common ambient highest-weight space proves the assertion. This is a coefficient-subalgebra version of the usual inheritance principle, proved here in precisely the form needed for the orbit restrictions.

The exact identities alone impose no efficient complexity bound. In particular, a nonnegative squared-norm formula may have exponentially many terms; a rank formula may require exponentially large matrices; and a positive dimension may become zero after highest-weight projection. The size and sign statements in each result are therefore part of the theorem, not suppressed implementation details.

\subsection{How the specific new analytic tools enter}

Alp\"oge's step~2 (p.~3) concerns theta derivatives and inert-prime divisibility. It does not, in the supplied four-page text, state a bound for a coefficient-matroid minor. Steps~3 and~5 (p.~3) motivate trace and square-root-cancellation arguments; P-02, P-04 and P-06 implement the corresponding curve/sheaf constructions with their exact degree and conductor budgets. Steps~4 and~7 (pp.~3--4) concern determinant and Gan--Gross--Prasad period positivity for specified groups; R-03 identifies why that branching problem cannot be substituted for the tensor subgroup without a proved functorial bridge. Step~8 (p.~4) concerns theta-lift multiplicities; I-03 and I-14 identify the different coefficient-span quantifiers a proposed isolation transfer must preserve.

Step~13 (p.~4) invokes a Bost slope inequality, and step~17 (p.~4) concerns the displayed arithmetic cohomology groups. The supplied note gives no reusable estimates for the ranks, heights or valuations of the ROABP coefficient maps used here. Their possible use therefore has a concrete prerequisite: construct the relevant arithmetic bundle or cohomology object, identify the desired minor or quotient with its invariant, and retain effective size bounds. This is a dependency statement about the supplied inputs, not a claim that such a construction is impossible. The integral and special-fiber analyses below identify where that next construction would act.

\part{Checked findings and checked conditional reductions}

% BEGIN SOURCE work/rectangular/findings.tex
% Research fragment prepared by /root/rectangular_bridge in run CUE_GCT-20261121.
% Requires amsmath, amssymb, amsthm; no nonstandard macros are used.
% Citations use bibliographic keys listed in sources.md.
% These are checked deductions in this run, not assertions of literature priority.

\section{Rectangular characters, explicit positive witnesses, and orbit ranks}

Throughout this section the ground field is $\mathbb C$, compact-group Haar
measure has mass one, and $D=nd$. We write $r=(d^n)$ and
$D_\alpha(b)=s_\alpha(1^b)=\dim S_\alpha(\mathbb C^b)$.
Duals are displayed when coordinate rings are involved. Passing from the
polynomial representation to the coordinate representation changes both sides
by the same duality and leaves every multiplicity below unchanged.

\subsection{R-01: the exact rectangular restriction transform}

\paragraph{Statement R-01.}
Let $a\geq\ell(\lambda)$ and let $X=(x_1,\ldots,x_a)$,
$Y=(y_1,\ldots,y_n)$. Denote the product alphabet by
$XY=(x_i y_j)_{1\leq i\leq a,1\leq j\leq n}$. Then
\begin{equation}\label{R01:branching}
 s_r(XY)=\sum_{\mu,\nu\vdash D}
 g(r,\mu,\nu)s_\mu(X)s_\nu(Y).
\end{equation}
In particular, by symmetry of the three Kronecker labels,
\begin{equation}\label{R01:haar}
 g(\lambda,r,r)=\int_{U(a)\times U(n)}
 s_r(A\otimes B)\overline{s_\lambda(A)}\det(B)^{-d}\,dA\,dB.
\end{equation}
Every coefficient needed on the rectangular-Kronecker side of determinant
GCT is therefore an explicit tensor-subgroup Fourier coefficient of the
\emph{same rectangular character} that occurs in the finite CUE moment.
The equality is not an identification of that coefficient with the scalar
moment.

More explicitly, enumerate the alphabet $XY$ as $z_1,\ldots,z_{an}$, with
$z_1,\ldots,z_n=x_1y_1,\ldots,x_1y_n$. Bump--Gamburd's Proposition~4
(source label \texttt{cpaverage}) gives
\begin{equation}\label{R01:CUE}
 s_r(XY)=\left(\prod_{j=1}^n z_j^d\right)
 \int_{U(d)}\prod_{j=1}^n\det(1+z_j^{-1}U^{-1})
 \prod_{j=n+1}^{an}\det(1+z_jU)\,dU.
\end{equation}
This formula is first read with nonzero $z_j$ and then as a polynomial identity.
Its specialization at identities proves the positive weighted sum
\begin{equation}\label{R01:dimension}
 D_r(an)=\sum_{\mu,\nu\vdash D}
 g(r,\mu,\nu)D_\mu(a)D_\nu(n)
\end{equation}
and the explicit upper bound
\begin{equation}\label{R01:bound}
 0\leq g(\lambda,r,r)\leq
 \left\lfloor\frac{D_r(an)}{D_\lambda(a)}\right\rfloor,
 \qquad
 D_r(an)=\prod_{i=1}^{n}\prod_{j=1}^{d}
 \frac{an+j-i}{n+d-i-j+1}.
\end{equation}
The numerator is a finite rectangular-character moment, not an unknown GCT
coefficient.

\begin{proof}
Apply Schur--Weyl duality to $(\mathbb C^a\otimes\mathbb C^n)^{\otimes D}$
in two ways. Decomposing first for $GL_{an}\times S_D$ gives the factors
$S_\eta(\mathbb C^a\otimes\mathbb C^n)\otimes[\eta]$.
Decomposing separately for $GL_a$ and $GL_n$ gives
$S_\mu(\mathbb C^a)\otimes S_\nu(\mathbb C^n)\otimes
([\mu]\otimes[\nu])$, with diagonal $S_D$ action. Comparing the multiplicity
space of $[r]$ proves \eqref{R01:branching}. Schur orthogonality and
$s_r(B)=\det(B)^d$ prove \eqref{R01:haar}. Equation \eqref{R01:CUE} is the
specialization $N=d,L=n,K=(a-1)n$ of the cited product identity.
Set every $x_i,y_j$ equal to one in \eqref{R01:branching}. Every summand is
nonnegative. Retain the term $(\mu,\nu)=(\lambda,r)$ and use
$D_r(n)=1$; this proves \eqref{R01:bound}. The displayed product is the
hook-content formula in Bump--Gamburd, equation \texttt{e\_hhl} in \S2.2.
\end{proof}

\paragraph{Finite moments and their limiting constants.}
Bump--Gamburd, Corollary~1, gives
\[
 \mathbb E_{U(N)}|\det(1-U)|^{2k}=D_{(N^k)}(2k).
\]
The limit after division by $N^{k^2}$ is
$c_k=\prod_{j=0}^{k-1}j!/(j+k)!$. Thus the finite moment is an integer
dimension, whereas the leading constant generally is not: $c_2=1/12$.
The distinction prevents a scalar asymptotic constant from being treated as
a representation multiplicity.

\subsection{R-02: swaps, tableaux, and positive projectors}

\paragraph{Statement R-02.}
Set $m=an$, $\delta_b=(b-1,b-2,\ldots,0)$, and
$\Delta_b(X)=\det(x_i^{b-j})$. For a tableau $T$ of shape $r$ with
semistandard entries in the ordered alphabet $\{1,\ldots,a\}\times
\{1,\ldots,n\}$, let $p(T)$ and $q(T)$ be the two marginal content vectors.
Write $N_r(p,q)$ for the number of these tableaux with marginal contents
$(p,q)$, and set it equal to zero when a content component is negative.
Then
\begin{align}\label{R02:tableaux}
 g(\lambda,r,r)
 &= [X^{\lambda+\delta_a}Y^{r+\delta_n}]
       \Delta_a(X)\Delta_n(Y)s_r(XY)\\
 &=\sum_{\sigma\in S_a,\tau\in S_n}
    \operatorname{sgn}(\sigma)\operatorname{sgn}(\tau)
    N_r(\lambda+\delta_a-\sigma\delta_a,
        r+\delta_n-\tau\delta_n).\nonumber
\end{align}
The exact unitary sum over swaps acts on the multiplicity by inserting
\begin{equation}\label{R02:swaps}
 s_{(d^n)}(z_1,\ldots,z_m)=
 \sum_{\substack{S\subseteq\{1,\ldots,m\}\\|S|=n}}
 \frac{\prod_{i\in S}z_i^d}
      {\prod_{i\in S,j\notin S}(1-z_j/z_i)}
\end{equation}
in the first line of \eqref{R02:tableaux}, and then setting $z_{ij}=x_i y_j$.
The sum in \eqref{R02:swaps} is combined before specialization or coefficient
extraction; its individual summands have poles and are not separately
nonnegative tableau counts.

There is, nevertheless, an exact positive formula. On the unitary space
\[
 T=S_\lambda(\mathbb C^n\otimes\mathbb C^n)
       \otimes\det(\mathbb C^n)^{-d}\otimes\det(\mathbb C^n)^{-d},
\]
let $P=\int_{U(n)\times U(n)}\rho_T(A,B)\,dA\,dB$. For any orthonormal
basis $(e_i)$ of $T$,
\begin{equation}\label{R02:projector}
 g(\lambda,r,r)=\operatorname{tr}P
   =\|P\|_{\mathrm{HS}}^2=\sum_{i,j}|\langle e_i,Pe_j\rangle|^2.
\end{equation}
In particular, this gives a rigorous sum of squares and a nonvanishing
criterion. It does not give a polynomial-size counting rule: the entries
of $P$ are themselves signed integrals, and the representation dimension
can be exponential in the natural parameters.

\begin{proof}
Multiply \eqref{R01:branching} by the two Vandermonde alternants.
For a partition $\mu$, the alternant $\Delta_a(X)s_\mu(X)$ has exactly one
monomial with strictly decreasing exponent vector equal to
$\lambda+\delta_a$ precisely when $\mu=\lambda$, and its coefficient then
is one. The analogous statement holds for $Y$. This proves the first line
of \eqref{R02:tableaux}. Expanding $s_r$ into semistandard tableaux and
the two alternants into their permutation terms proves the second line.
Formula \eqref{R02:swaps} is the rectangular Weyl-character Laplace expansion
used in Bump--Gamburd, Proposition~4 and \S4. For the last assertion,
Haar invariance gives $P^2=P$ and $P^*=P$; its range is exactly $T^{U(n)\times
U(n)}$. The tensor-subgroup branching identity identifies the dimension of
that range with $g(\lambda,r,r)$. Taking traces of $P^*P=P$ proves
\eqref{R02:projector}.
\end{proof}

\paragraph{What the arithmetic CFKRS recipe supplies.}
CFKRS's unitary formula is a theorem (Theorem~1.5.2 in the supplied source),
and its swap/contour conversion is Lemma~2.5.1. The corresponding
$L$-function mean-value assertion is Conjecture~1.5.1, derived by the
recipe in \S2.1 and the sum over $\Xi$ in \S2.2, equation
\texttt{conj:M}. Replacing the unitary factors in \eqref{R02:swaps} by the
arithmetic Euler product and shifted zeta factors does not define a
character restriction map. Neither RH nor GRH supplies the discarded
off-diagonal estimates in that recipe. The proven action on GCT is
\eqref{R02:tableaux}--\eqref{R02:swaps}; the arithmetic recipe remains a
separate conjectural mean-value statement.

\subsection{R-03: three explicit limits on positivity transfer}

\paragraph{Statement R-03.}
The following three proposed simplifications are false.
First, positivity of a rectangular CUE moment does not imply positivity
of every rectangular Kronecker coefficient. Second, the connected
left-right stabilizer multiplicity need not equal the full determinant
stabilizer multiplicity. Third, one cannot replace tensor-subgroup
branching by one multiplicity-free interlacing step.

\begin{proof}
For $r=(2,2)$ the $S_4$ representation $[r]$ is the two-dimensional
nontrivial summand in the permutation representation on the three
perfect matchings of four letters. Its character at the cycle types
$1^4,21^2,2^2,31,4$ is respectively $2,0,2,-1,0$; the class sizes are
$1,6,3,8,6$. The standard character $[31]$ is $3,1,-1,0,-1$.
Their character inner product gives
\[
 g((3,1),(2,2),(2,2))=(12-12)/24=0,
\]
although $D_{(2,2)}(4)=20>0$. The same elementary calculation yields
\[
 [22]\otimes[22]=[4]\oplus[22]\oplus[1111],\qquad
 \operatorname{Sym}^2[22]=[4]\oplus[22].
\]
Thus $g((1^4),(22),(22))=1$ but the corresponding symmetric Kronecker
coefficient is zero. These values are checked independently by enumerating
the 24 permutations in \texttt{check\_small\_cases.py}.

The tensor homomorphism $U(a)\times U(n)\longrightarrow U(an)$ has kernel
$\{(zI_a,z^{-1}I_n):|z|=1\}$. Its image has one-dimensional connected
center, whereas the block-diagonal Levi $U(a)\times U(n)$ has
two-dimensional connected center. For $a,n\geq2$ these are different
branching problems. Moreover, R-07 below constructs two independent
highest-weight invariants for $g((12,4),(4^4),(4^4))$, so its value is at
least two (in fact exactly two). A single $0/1$ interlacing rule cannot
be its multiplicity formula.
\end{proof}

The $0/1$ Gan--Gross--Prasad periods cited in Alp\"oge's steps~4 and~7
therefore require an additional, explicit comparison of branching
functors before they can control these multiplicities. Iterating
multiplicity-free restrictions can create multiple paths; no
contradiction with such restrictions is asserted. Likewise the squared
Gram formulas are not statements that every Weingarten summand is
positive: Collins--\'Sniady, Corollary~2.4 and its following example, give
$\operatorname{Wg}_N((12))=-1/(N(N^2-1))$ for $N\geq2$.

\subsection{R-04: plethysm requires the joint law of powers}

\paragraph{Statement R-04.}
For $N\geq\ell(\lambda)$ and $|\lambda|=nd$, the plethysm multiplicity is
\begin{equation}\label{R04:plethysm}
 a_\lambda(d[n])=\int_{U(N)}\overline{s_\lambda(U)}
 \sum_{\kappa\vdash d}\frac{1}{z_\kappa}
       \prod_{j\geq1}s_{(n)}(U^j)^{m_j(\kappa)}\,dU,
 \qquad z_\kappa=\prod_j j^{m_j(\kappa)}m_j(\kappa)!.
\end{equation}
Rains's Theorem~1.3 determines the eigenvalue law of each single $U^j$,
but replacing the different powers in \eqref{R04:plethysm} by independent
matrices with those marginal laws is invalid. A quantitative example is
\begin{equation}\label{R04:counter}
 \int_{U(2)}(\operatorname{tr}U)^4
        \overline{(\operatorname{tr}U^2)^2}\,dU=2,
\end{equation}
whereas the product of the two marginal expectations is zero.

\begin{proof}
For any finite-dimensional $W$, the cycle-index formula for the
symmetrizer $d!^{-1}\sum_{\sigma\in S_d}\sigma$ gives
$\chi_{\operatorname{Sym}^dW}(g)=\sum_{\kappa\vdash d}
z_\kappa^{-1}\prod_j\chi_W(g^j)^{m_j(\kappa)}$.
Set $W=\operatorname{Sym}^n\mathbb C^N$ and apply Schur orthogonality.
For \eqref{R04:counter}, in two variables
\[
 p_1^4=s_{(4)}+3s_{(3,1)}+2s_{(2,2)},\qquad
 p_2^2=s_{(4)}-s_{(3,1)}+2s_{(2,2)}.
\]
Orthogonality gives $1-3+4=2$. Multiplication of $U$ by a scalar phase
shows that both separate expectations vanish. The two polynomial
expansions are also verified by the exact computation certificate.
\end{proof}

In representation-ring language $U\mapsto U^j$ yields the Adams
operation on characters, not a homomorphism of nonabelian groups.
Already $\psi^2(V)=\operatorname{Sym}^2V-\bigwedge^2V$ has a negative
Schur coefficient. A positive alternative retaining the complete
coupling is
\begin{equation}\label{R04:wreath}
 a_\lambda(d[n])=\dim[\lambda]^{S_n\wr S_d}
 =\sum_{i,j}\left|\frac{1}{(n!)^d d!}
       \sum_{h\in S_n\wr S_d}\rho_\lambda(h)_{ij}\right|^2.
\end{equation}
Indeed, apply Schur--Weyl duality to $V^{\otimes nd}$ and the orthogonal
projector onto the wreath-product invariants, exactly as in R-02.
The formula is positive and exact; its factorial enumeration cost is
explicit, and it is not a solution of the efficient positive-formula
problem.

\subsection{R-05: Haar covariance computes the orbit-closure ranks}

\paragraph{Statement R-05.}
Let $V=\mathbb C^N$, $W=\operatorname{Sym}^nV$, $0\ne f\in W$, and
$X_f=\overline{GL(V)f}$. Put $E_d=\operatorname{Sym}^dW$ and
$v=f^{\otimes d}\in E_d$. Define
\begin{equation}\label{R05:twirl}
 M_f=\int_{U(N)}|\rho(k)v\rangle\langle\rho(k)v|\,dk.
\end{equation}
Under the Hermitian identification of $E_d$ with its conjugate dual,
$\ker M_f$ is the degree-$d$ ideal of $X_f$, and the image is dual to
$\mathbb C[X_f]_d$. Write the isotypic decomposition
\[
 E_d=\bigoplus_{\lambda\vdash nd} L_\lambda\otimes\mathbb C^{a_\lambda},
 \qquad L_\lambda=S_\lambda V,
 \qquad v_\lambda=\sum_{i=1}^{D_\lambda(N)}e_i\otimes m_i.
\]
Then the $\lambda$ block of $M_f$ is
\begin{equation}\label{R05:block}
 \frac{I_{L_\lambda}}{D_\lambda(N)}\otimes R_\lambda(f),
 \qquad R_\lambda(f)=\sum_i|m_i\rangle\langle m_i|.
\end{equation}
The multiplicity of $L_\lambda^*$ in $\mathbb C[X_f]_d$ is therefore
\begin{equation}\label{R05:rank}
 m_\lambda(f)=\operatorname{rank}R_\lambda(f).
\end{equation}
If $F_\lambda$ is the matrix with columns $m_i$, then, for every $s$,
\begin{equation}\label{R05:minors}
 e_s(R_\lambda(f))=
 \sum_{\substack{I\subseteq\{1,\ldots,a_\lambda\}\\|I|=s}}
 \ \sum_{\substack{J\subseteq\{1,\ldots,D_\lambda(N)\}\\|J|=s}}
       |\det(F_{\lambda,I,J})|^2.
\end{equation}
In particular $m_\lambda(f)\geq s$ if and only if this nonnegative
polynomial is positive. A multiplicity obstruction between two orbit
closures is thus an exact comparison of these Gram ranks, including
the boundary equations, rather than of their scalar traces.

\begin{proof}
For $u\in E_d$,
$\langle u,M_fu\rangle=\int|\langle u,\rho(k)v\rangle|^2\,dk$.
This is zero exactly when the continuous nonnegative integrand is
identically zero. The compact group $U(N)$ is Zariski dense in
$GL_N(\mathbb C)$: a regular function vanishing on it has all derivatives
along its real Lie algebra zero, and the complexification of that Lie
algebra is $\mathfrak{gl}_N(\mathbb C)$; analytic continuation on the
connected complex group gives zero. Consequently the same vanishing
holds on $GL(V)f$ and its closure. A degree-$d$ polynomial on $W$ is
exactly a linear functional on $E_d$ evaluated at $f^{\otimes d}$.
This identifies the kernel with the homogeneous ideal and the image
with the dual quotient.

Off-diagonal isotypic blocks of the twirl vanish by Schur's lemma.
Within $L_\lambda\otimes\mathbb C^{a_\lambda}$, averaging
$|e_i\rangle\langle e_j|$ gives
$\delta_{ij}I_{L_\lambda}/D_\lambda(N)$, proving
\eqref{R05:block}. Thus the orbit span is
$L_\lambda\otimes\operatorname{span}\{m_i\}$, which proves
\eqref{R05:rank}. Finally $R_\lambda=F_\lambda F_\lambda^*$.
The coefficient $e_s$ is the sum of its principal $s$-minors;
Cauchy--Binet gives \eqref{R05:minors}. Since a positive semidefinite
matrix has nonnegative eigenvalues, $e_s>0$ exactly when at least $s$
of those eigenvalues are positive.
\end{proof}

For $f=\det_n$ or the padded permanent, \eqref{R05:twirl} is a Haar
integral over $U(n^2)$. Its entries have degree $nd$ in the matrix
entries and degree $nd$ in their conjugates. Thus the exact
Collins--\'Sniady formula evaluates them. The full, dimension-truncated
Weingarten inverse must be used when tensor degree exceeds group
dimension. The remaining task is to control the ranks of the resulting
structured matrices; knowing that the matrices are positive
semidefinite does not determine which minors vanish.

\subsection{R-06: a Gaussian and integral form of the rank certificate}

\paragraph{Statement R-06.}
Use the tensor realization $E_d\subseteq V^{\otimes D}$ and set
$v=f^{\otimes d}$. For an $N\times N$ matrix of indeterminates $Z$, write
\begin{equation}\label{R06:coefficients}
 Z^{\otimes D}v=\sum_{|\alpha|=D}Z^\alpha c_\alpha,
 \qquad C_f=(c_\alpha)_\alpha.
\end{equation}
There are $N^D$ tensor coordinates and
$\binom{N^2+D-1}{D}$ columns. If $G$ has independent standard complex
Gaussian entries, then
\begin{equation}\label{R06:gaussian}
 \mathbb E|G^{\otimes D}v\rangle\langle G^{\otimes D}v|
   =\sum_{|\alpha|=D}\alpha!\,c_\alpha c_\alpha^*.
\end{equation}
Its image is the orbit span from R-05. It has the same kernel and
the same isotypic ranks as the Haar covariance.

Let $f^\lambda=\dim[\lambda]$ and let
\[
 e_\lambda=\frac{f^\lambda}{D!}
       \sum_{\sigma\in S_D}\chi^\lambda(\sigma)\sigma
\]
act by permutation of tensor factors. Then
\begin{equation}\label{R06:integer}
 A_\lambda(f)=\sum_{\sigma\in S_D}\chi^\lambda(\sigma)\sigma C_f,
 \qquad
 \operatorname{rank}A_\lambda(f)=D_\lambda(N)m_\lambda(f).
\end{equation}
When $f$ is an integral symmetric tensor, $A_\lambda(f)$ has integer
entries. If its tensor entries have absolute value at most $H$, then
every entry of $A_\lambda(f)$ has absolute value at most
$D!f^\lambda N^D H^d$.
The same rank is obtained by replacing the coefficient columns by all
evaluations at the finite grid $Z\in\{0,1,\ldots,D\}^{N^2}$ and then
applying $e_\lambda$. This grid and the coefficient matrix are
explicit but not polynomial-size constructions.
An ordinary polynomial with integer monomial coefficients may require
scaling by $n!$ under the symmetric-tensor embedding to meet the
integrality hypothesis; this nonzero scalar does not change any orbit
span or multiplicity in characteristic zero.

\begin{proof}
For a standard complex Gaussian $z$, polar integration gives
$\mathbb E[z^p\bar z^q]=0$ if $p\ne q$ and $p!$ if $p=q$.
Independence proves \eqref{R06:gaussian}. Its image is
$\operatorname{span}\{c_\alpha\}$. A linear functional kills all
$c_\alpha$ exactly when its evaluation on $Z^{\otimes D}v$ is the zero
polynomial. Since invertible matrices are Zariski dense, that is
equivalent to killing the $GL_N$ orbit span. This proves the image and
kernel assertions without analytic approximation.

The central idempotent $e_\lambda$ projects precisely onto the
$L_\lambda$ isotypic summand by Schur--Weyl duality. Consequently
$e_\lambda C_f$ has rank $D_\lambda(N)m_\lambda(f)$, and multiplication
by the nonzero scalar $D!/f^\lambda$ gives \eqref{R06:integer}.
Symmetric-group characters and permutation matrices have integral
entries. Each tensor coordinate of $Z^{\otimes D}v$ is a sum of $N^D$
terms whose coefficients have magnitude at most $H^d$; hence
$|(c_\alpha)_i|\leq N^D H^d$. The character bound
$|\chi^\lambda(\sigma)|\leq f^\lambda$ and the $D!$ permutations give
the stated estimate. Finally a polynomial of individual degree at
most $D$ vanishing on a Cartesian product of $D+1$ distinct scalars is
zero, by induction on the number of variables. Applying this fact to
every linear functional proves equality of the grid and coefficient
spans.
\end{proof}

This replaces signed Haar expansions by an explicitly positive
coefficient Gram matrix. Cauchy--Binet gives positive sums of squared
integer minors after clearing the indicated denominators. It also
exhibits the exact cost: the tensor model has $N^D$ rows, the
symmetrizer has $D!$ terms, and the displayed grid has
$(D+1)^{N^2}$ points. A claimed efficient positivity algorithm must
compress these data or exploit their ranks; positivity alone does
not achieve that compression.

\subsection{R-07: explicit positive witnesses that survive the boundary}

\paragraph{Statement R-07.}
Let $n\geq2$, $d\geq0$, and $0\leq q\leq d$. For
$\lambda=(nd-q,q)$, define
\[
 p_{2,\ldots,n}(q)=[t^q]\prod_{j=2}^n(1-t^j)^{-1}.
\]
There are explicitly constructed, linearly independent highest-weight
vectors, indexed by partitions of $q$ into parts $2,\ldots,n$, in the
degree-$d$ coordinate ring of the determinant orbit closure. In fact
\begin{equation}\label{R07:positive}
 m_\lambda(\det_n)=a_\lambda(d[n])=p_{2,\ldots,n}(q),
 \qquad
 g(\lambda,(d^n),(d^n))\geq p_{2,\ldots,n}(q).
\end{equation}
The last inequality is equality by Manivel, Corollary~2(1) in \S3.3; it is a specialization of his
stable rectangular Kronecker Theorem~1, restated as
Ikenmeyer--Panova, Theorem~2.1. The construction below gives concrete witnesses for this
known numerical positivity and, simultaneously, proves their survival
on the determinant orbit closure.

More generally, with no restriction $q\leq d$,
\begin{equation}\label{R07:alltwo}
 m_\lambda(\det_n)=a_\lambda(d[n])
 \quad\text{for every }\lambda\vdash nd\text{ with }\ell(\lambda)\leq2.
\end{equation}
Therefore no multiplicity obstruction of length at most two can
separate a padded-permanent orbit closure from the determinant orbit
closure in the same ambient space and degree.

\begin{proof}
Write a binary degree-$n$ form as
$F(s,t)=\sum_{j=0}^n\binom nj b_j s^{n-j}t^j$.
In the coordinate representation $\operatorname{Sym}^n\mathbb C^2$,
the upper-triangular raising operator is
$E b_j=j b_{j-1}$, with $E b_0=0$, and the weight of $b_j$ is
$(n-j,j)$. For $2\leq j\leq n$ put
\begin{equation}\label{R07:central}
 Q_j=b_0^{j-1}\sum_{i=0}^j\binom ji b_i
                  \left(-\frac{b_1}{b_0}\right)^{j-i}.
\end{equation}
This is a polynomial: for $i=0$ the factor $b_i=b_0$ cancels the only
possible negative exponent; all $i\geq1$ terms already have
nonnegative powers of $b_0$. The first two examples are
\[
 Q_2=b_0b_2-b_1^2,\qquad
 Q_3=b_0^2b_3-3b_0b_1b_2+2b_1^3.
\]
Under $\exp(uE)$, $b_i$ goes to
$\sum_{h=0}^i\binom ih u^{i-h}b_h$ and $b_1/b_0$ goes to
$b_1/b_0+u$. The binomial theorem in \eqref{R07:central} shows that
the two translations cancel. Thus $E Q_j=0$. Every term has degree
$j$ in the $b_i$ and weight $((n-1)j,j)$.

For a tuple of nonnegative integers $(a_2,\ldots,a_n)$ satisfying
$\sum_{j=2}^n ja_j=q$, define
\begin{equation}\label{R07:witness}
 P_{\mathbf a}=b_0^{d-q}\prod_{j=2}^n Q_j^{a_j}.
\end{equation}
This is a highest-weight polynomial of degree $d$ and weight
$(nd-q,q)$. Evaluate it on determinant pencils
$F(s,t)=\det(sI_n+tY)$, with $Y$ diagonal. Every monic binary form
is obtained in this way over $\mathbb C$ by factoring the corresponding
one-variable polynomial. In particular, one may set $b_0=1,b_1=0$
and choose $b_2,\ldots,b_n$ arbitrarily. On this family $Q_j=b_j$,
so the restrictions of the $P_{\mathbf a}$ are distinct monomials
in independent variables. They are linearly independent.

For completeness, the ambient highest-weight dimension in this range
is obtained by a short weight count. In
$\operatorname{Sym}^d(\operatorname{Sym}^n\mathbb C^2)$ the weight
$(nd-q,q)$ has dimension equal to the number of partitions of $q$
into at most $d$ parts, each at most $n$. When $q\leq d$ this is
$[t^q]\prod_{j=1}^n(1-t^j)^{-1}$. A $GL_2$ irreducible of highest
weight $(nd-r,r)$ has one vector at each weight $(nd-q,q)$ with
$r\leq q\leq nd-r$. Subtracting consecutive weight dimensions
therefore gives $p_{2,\ldots,n}(q)$. This proves that the displayed
independent witnesses form a basis and establishes the first equality
in \eqref{R07:positive} once boundary survival is checked.

Now let $V$ have dimension $n^2$ and view a binary form as an element
of $\operatorname{Sym}^nV^*$ supported on the first two coordinates.
It factors as $\prod_{i=1}^n\ell_i$ and hence equals
$\det(\operatorname{diag}(\ell_1,\ldots,\ell_n))$.
This is a linear, possibly singular, substitution into $\det_n$.
Since $GL(V)$ is dense in $\operatorname{End}(V)$ and substitution
is polynomial, every such form lies in the determinant orbit closure.
A vector of highest weight $\lambda=(\lambda_1,\lambda_2,0,\ldots,0)$
in $\operatorname{Sym}^d(\operatorname{Sym}^nV)$ uses only the first
two basis vectors: each other weight coordinate is a sum of
nonnegative exponents and equals zero. The additional raising
operators annihilate it, so its highest-weight condition is exactly
the binary one. A nonzero polynomial in these binary coordinates
cannot vanish on every binary form. Thus the homogeneous ideal of
the determinant orbit closure has no nonzero highest-weight vector
of length at most two. Taking the quotient proves
\eqref{R07:alltwo}, including all degrees and all two-row weights.

The determinant-coordinate multiplicity is at most its orbit
multiplicity, which is the symmetric rectangular Kronecker coefficient,
and that is at most $g(\lambda,r,r)$; this is the cited stabilizer
statement in B\"urgisser--Ikenmeyer--Panova, \S1(b).
This proves the inequality in \eqref{R07:positive}. Finally any
padded-permanent coordinate ring is a quotient of the same ambient
$\operatorname{Sym}(\operatorname{Sym}^nV)$, so its multiplicity is
at most $a_\lambda(d[n])=m_\lambda(\det_n)$ for two-row $\lambda$.
\end{proof}

For $n=d=q=4$, the two witnesses are $Q_4$ and $Q_2^2$, with
\[
 Q_4=b_0^3b_4-4b_0^2b_1b_3+6b_0b_1^2b_2-3b_1^4.
\]
They specialize to $b_4$ and $b_2^2$, proving directly that the
determinant coordinate multiplicity at $(12,4)$ is two. This is
stronger information than nonnegativity of a Haar integral.

\subsection{R-08: the exact Molien numerator needed for orbit closures}

\paragraph{Statement R-08.}
Let $S=\operatorname{Sym}(W)$ be a polynomial coordinate algebra,
$A=S/I$ a graded $GL_N$-equivariant quotient, and $K=U(N)$. Then
\begin{equation}\label{R08:molien}
 \sum_{d\geq0}m_\lambda(A_d)t^d
   =\int_K\overline{\chi_\lambda(k)}H_A(t,k)\,dk.
\end{equation}
For the ambient algebra,
$H_S(t,k)=\det(1-t\rho_W(k))^{-1}$.
For a finite graded equivariant free resolution
$\cdots\to S\otimes T_1\to S\otimes T_0\to A\to0$,
\begin{equation}\label{R08:numerator}
 H_A(t,k)=
 \frac{\sum_i(-1)^i\chi_{T_i}(t,k)}
      {\det(1-t\rho_W(k))}.
\end{equation}
Thus the Molien denominator and ratios identities compute the desired
multiplicities only after the orbit ideal's equivariant numerator,
or equivalently the ranks in R-05--R-06, is supplied.

\begin{proof}
Coefficientwise Schur orthogonality proves \eqref{R08:molien}.
The eigenvalue expansion of the symmetric algebra gives the
denominator formula. Alternating graded characters are additive
on exact complexes, yielding \eqref{R08:numerator}. For a purely
formal interpretation no convergence assumption is needed; for
analytic evaluation take $|t|<1$ in the unitary polynomial case.
\end{proof}

\subsection{R-09: an exact ratios kernel for squared GCT multiplicities}

\paragraph{Statement R-09.}
Let $n,d,N\geq1$, $r=(d^n)$, and define
\[
 S_N(n,d)=\sum_{\substack{\lambda\vdash nd\\\ell(\lambda)\leq N}}
                    g(\lambda,r,r)^2.
\]
Set $S_0(n,d)=0$.
For $0<t<1$ and $A,B,C,D\in U(n)$ put
\begin{align*}
 \mathcal R_N(t;A,B,C,D)=\int_{U(N)}
 &\det\bigl(I-t(A\otimes B)\otimes U\bigr)^{-1}\\[-2mm]
 &\quad\cdot
 \det\bigl(I-t(C\otimes D)\otimes U^{-1}\bigr)^{-1}\,dU.
\end{align*}
Then the explicitly specified fourfold Fourier projection is
\begin{equation}\label{R09:squares}
 \int_{U(n)^4}\mathcal R_N(t;A,B,C,D)
   (\det A\det B\det C\det D)^{-d}\,dA\,dB\,dC\,dD
       =t^{2nd}S_N(n,d).
\end{equation}
For $N\geq n^2$, the inner ratio is evaluated exactly by
Bump--Gamburd, Theorem~3, at $K=L=0,Q=R=n^2$:
\begin{equation}\label{R09:evaluated}
 \mathcal R_N(t;A,B,C,D)
       =\det\bigl(I_{n^4}-t^2 A\otimes B\otimes C\otimes D\bigr)^{-1}.
\end{equation}
Consequently an explicitly evaluated Molien--Weyl coefficient is the
sum of squares of precisely the rectangular Kronecker coefficients
needed in GCT. The truncated sums satisfy
\begin{equation}\label{R09:monotone}
 S_N(n,d)-S_{N-1}(n,d)
    =\sum_{\substack{\lambda\vdash nd\\\ell(\lambda)=N}}g(\lambda,r,r)^2
    \geq0,
 \qquad
 g(\lambda,r,r)\leq\lfloor\sqrt{S_N(n,d)}\rfloor
 \quad(\ell(\lambda)\leq N).
\end{equation}
The stabilized value has a second nonnegative moment expression:
\begin{equation}\label{R09:finitefourth}
 S_{n^2}(n,d)=\frac1{(nd)!}\sum_{\sigma\in S_{nd}}\chi^r(\sigma)^4.
\end{equation}

\begin{proof}
The Cauchy identity and absolute convergence for $t<1$ expand the two
inverse determinants into
$\sum_\lambda t^{|\lambda|}s_\lambda(A\otimes B)s_\lambda(U)$
and
$\sum_\mu t^{|\mu|}s_\mu(C\otimes D)s_\mu(U^{-1})$.
Schur orthogonality on $U(N)$ retains exactly $\lambda=\mu$ of
length at most $N$. The two applications of the tensor-subgroup
projection in R-02 then retain $|\lambda|=nd$ and give
$g(\lambda,r,r)$ from each pair of $U(n)$ factors. This proves
\eqref{R09:squares}. When $N\geq n^2$, the length cutoff is
inactive because both product alphabets have $n^2$ entries. A second
Cauchy identity gives \eqref{R09:evaluated}, equivalently the stated
specialization of the ratios theorem. The definition immediately
proves \eqref{R09:monotone}. Finally the character of
$[r]\otimes[r]$ is $(\chi^r)^2$. Its squared character norm is
the sum of squares of its irreducible multiplicities, proving
\eqref{R09:finitefourth}; symmetric-group characters are real.
The tensor-subgroup branching formula also gives
$g(\lambda,r,r)=0$ for $\ell(\lambda)>n^2$, because
$S_\lambda(\mathbb C^n\otimes\mathbb C^n)=0$ in that range.
Thus the cutoff in the left side of \eqref{R09:finitefourth}
omits no nonzero multiplicity.
\end{proof}

The positive quantities in \eqref{R09:squares} are obtained after
projection. The pointwise ratio and individual swap summands need
not be nonnegative. An approximation argument would need a
certified error below one half to recover an integer coefficient,
with bounds uniform in the growing tensor degree. Neither the
closed ratio formula nor RH supplies those estimates or compresses
the remaining fourfold Fourier projection.

\subsection{R-10: the complete height-two rectangle and its boundary gate}

\paragraph{Statement R-10 (known formula, specialized and compared).}
Let $d\geq0$ and $\lambda\vdash2d$. For partitions of length at
most four, append zero parts to obtain $(\lambda_1,\ldots,\lambda_4)$.
Then
\begin{equation}\label{R10:rectangle}
g(\lambda,(d,d),(d,d))=
\begin{cases}
1,&\ell(\lambda)\leq4\text{ and all four parts have the same parity},\\
0,&\text{otherwise}.
\end{cases}
\end{equation}
More precisely, the multiplicity of $[\lambda]$ in
$\operatorname{Sym}^2[d,d]$ is one exactly for partitions with all
parts even and length at most four. Its multiplicity in
$\bigwedge^2[d,d]$ is one exactly for partitions with exactly four
odd parts. These formulas are the constant-term specialization of
Blasiak--Mulmuley--Sohoni, GCT~IV, Corollary~18.4
(source label \texttt{c kronecker generating function sym ext});
they are cited known results, not new numerical formulas.

In the determinant problem with $n=2$, $N=4$, the boundary side is
also completely determined:
\begin{equation}\label{R10:closure}
m_\lambda(\det_2)=a_\lambda(d[2])=
\begin{cases}
1,&\ell(\lambda)\leq4\text{ and every part of }\lambda\text{ is even},\\
0,&\text{otherwise}.
\end{cases}
\end{equation}
Thus all four-odd-part summands of the ordinary rectangular
coefficient are absent from the determinant coordinate ring.
For $d=2$, this includes the $\lambda=(1,1,1,1)$ example in R-03.
Simultaneous conjugation of the two inner partitions also extends
\eqref{R10:rectangle}, as an ordinary Kronecker formula, to rectangles
of width two.

\begin{proof}
Write $c_i=\lambda_i-\lambda_{i+1}$ for $1\leq i\leq3$ and
$c_4=\lambda_4$, so $c_i$ is the number of columns of height $i$.
GCT~IV, Corollary~18.4, gives the generating function
$\sum_{\alpha\vdash2d,\ell(\alpha)\leq2}
g_{\varepsilon,\alpha,\lambda}x^{\alpha_1-\alpha_2}$.
Three of its four cases have a factor $x$ or $x^2$ multiplying
polynomials with nonnegative exponents, and hence contribute no
constant term. The remaining case is the product
$\llbracket c_1\rrbracket\llbracket c_2\rrbracket
\llbracket c_3\rrbracket$, where
$\llbracket b\rrbracket=x^b+x^{b-2}+\cdots$ ends at $1$ or $x$.
Its constant coefficient is one exactly when $c_1,c_2,c_3$ are all
even; the case condition then requires
$\varepsilon=(-1)^{c_3+c_4}=(-1)^{\lambda_4}$.
Thus the symmetric square retains the all-even partitions and the
exterior square retains the four-odd-part partitions. Their sum is
\eqref{R10:rectangle}. Length greater than four gives zero by
tensor-subgroup branching. Finally
$[\rho']=[\rho]\otimes\operatorname{sgn}$, so conjugating both
inner partitions leaves their tensor product unchanged.

The determinant $x_{11}x_{22}-x_{12}x_{21}$ is a nondegenerate
quadratic form in four variables. Over $\mathbb C$ all nondegenerate
quadratic forms of a fixed dimension are congruent. They form a
nonempty Zariski-open subset of the space of quadrics, so
$\operatorname{Det}_2=\operatorname{Sym}^2(\mathbb C^4)^*$.
Consequently its coordinate ring is
$\operatorname{Sym}(\operatorname{Sym}^2\mathbb C^4)$.

Here is a direct highest-weight proof of its familiar multiplicity
formula. In $\operatorname{Sym}(\operatorname{Sym}^2\mathbb C^s)$
write the generators as the entries of a symmetric matrix $Q$,
with $q_{ij}$ of weight $e_i+e_j$. The upper-triangular action on
these generators is $Q\mapsto g^tQg$.
Let $\Delta_i$ be the leading principal $i\times i$ minor and
$\Delta_0=1$. Each $\Delta_i$ is fixed by the upper unitriangular
group and has weight $2(e_1+\cdots+e_i)$.
On the dense open set where all $\Delta_i$ are nonzero, elementary
upper-triangular congruences reduce $Q$ to the unique diagonal
matrix with entries $\Delta_i/\Delta_{i-1}$.
Therefore a nonzero unitriangular-invariant polynomial has nonzero
restriction to diagonal matrices. If it is also a torus eigenvector
of weight $\lambda$, this restriction is a scalar multiple of
$\prod_i q_{ii}^{\lambda_i/2}$. It exists only when each
$\lambda_i$ is even, and its value determines the polynomial on a
dense open set. Hence each highest-weight space has dimension at
most one and vanishes unless $\lambda=2\mu$.
Conversely, for every partition $\mu$ of length at most $s$,
\[
 \prod_{i=1}^s\Delta_i^{\mu_i-\mu_{i+1}},\qquad\mu_{s+1}=0,
\]
is a nonzero highest-weight polynomial of weight $2\mu$ and
degree $|\mu|$. This proves
$\operatorname{Sym}^d(\operatorname{Sym}^2\mathbb C^s)
=\bigoplus_{\mu\vdash d,\ell(\mu)\leq s}S_{2\mu}\mathbb C^s$
and, at $s=4$, \eqref{R10:closure}.
\end{proof}

The stabilizer interpretation agrees with these formulas.
The full linear stabilizer of the nondegenerate quadratic form is
$O_4$, whereas the connected tensor image of $SL_2\times SL_2$ is
$SO_4$: its kernel is $\{(I,I),(-I,-I)\}$, so its dimension is six,
the dimension of $SO_4$. The extra orthogonal component removes
the odd sector. Thus the connected-stabilizer upper bound, the
symmetric rectangular coefficient, and the closure multiplicity
are completely distinguished in this case.

\subsection{R-11: small-degree padding excludes all four-row obstructions}

\paragraph{Statement R-11.}
Work over $\mathbb C$. Let $1\leq m\leq3$, $n\geq m$, $N=n^2$,
$V=\mathbb C^N$, $f\in\operatorname{Sym}^mV^*$, and $\ell\in V^*$.
Set $h=\ell^{n-m}f$. For every coordinate-ring degree $d\geq0$
and every $\lambda\vdash nd$ with $\ell(\lambda)\leq4$,
\begin{equation}\label{R11:fourrow}
 m_\lambda(h)\leq m_\lambda(\det_n).
\end{equation}
The assertion includes padded permanents of degree at most three.
It is a deduction from classical linear determinantal
representations, not a claim that arbitrary $f$ itself has an
$n\times n$ determinantal representation.

\begin{proof}
First consider forms in at most four variables. Linear forms are
$1\times1$ determinants. Every quadric in at most four variables
is congruent over $\mathbb C$ to $y_1^2+\cdots+y_r^2$, $r\leq4$.
With unused $y_i$ set to zero it is the determinant
\begin{equation}\label{R11:quadric}
 \det\begin{pmatrix}
 y_1+\mathrm i y_2&y_3+\mathrm i y_4\\
 -y_3+\mathrm i y_4&y_1-\mathrm i y_2
 \end{pmatrix}=
 y_1^2+y_2^2+y_3^2+y_4^2.
\end{equation}
Beauville (2000), Corollary~6.4, states over an algebraically
closed field that every smooth cubic surface has a $3\times3$
linear determinantal equation. Thus over $\mathbb C$ every smooth
quaternary cubic is a $3\times3$ linear determinant, after
absorbing its nonzero scalar factor in one matrix row.
Smooth cubics are Zariski dense in the affine space of quaternary
cubics: the discriminant-complement is open and nonempty, as the
Fermat cubic shows. Therefore every quaternary cubic lies in the
closure of the linear-determinant image.

If $q=\det M$ has degree $m$ and $u$ is linear, then
$u^{n-m}q=\det\operatorname{diag}(M,uI_{n-m})$.
The entries of this block matrix give a linear substitution into
$\det_n$. Such substitutions belong to $\operatorname{Det}_n$,
because invertible linear maps are dense in all linear maps and
substitution is polynomial. For fixed $u$, the preimage of
$\operatorname{Det}_n$ under the linear map $q\mapsto u^{n-m}q$
is closed. The density just proved extends the assertion from
smooth cubics to every cubic. We have shown that
\begin{equation}\label{R11:projectedcontainment}
 u^{n-m}q\in\operatorname{Det}_n
 \quad\text{for every }q\text{ of degree }m\leq3
 \text{ and }u\text{ supported on at most four variables}.
\end{equation}

Put $k=\min(4,N)$ and let $\pi$ set the last $N-k$ variables
equal to zero. Let $P$ be a highest-weight vector of weight
$\lambda$ in the degree-$d$ ideal of $\operatorname{Det}_n$.
Since $\lambda_j=0$ for $j>k$ and all coordinate generators have
nonnegative weights, every monomial of $P$ uses only coefficients
of forms supported on the first $k$ variables. Hence
$P(F)=P(\pi F)$ for every $F\in\operatorname{Sym}^nV^*$.
For $g\in GL(V)$ the projection is multiplicative on forms, so
\[
 \pi(g\cdot h)=(\pi(g\cdot\ell))^{n-m}\pi(g\cdot f)
 \in\operatorname{Det}_n
\]
by \eqref{R11:projectedcontainment}. It follows that
$P(g\cdot h)=0$ for every $g$, and hence $P$ belongs to the
ideal of $\overline{GL(V)\cdot h}$.
Thus the highest-weight space of the determinant ideal is
contained in that of the padded-form ideal for every such
$\lambda$. Finite-dimensional polynomial $GL(V)$-modules are
semisimple in characteristic zero. Subtracting the ideal
multiplicity from the common ambient multiplicity gives
\eqref{R11:fourrow}. If $N<4$, partitions of length greater than
$N$ already have multiplicity zero, so the same statement holds.
\end{proof}

\subsection{R-12: exact permanent-pencil dominance and four-row deficits}

\paragraph{Statement R-12 (computer-assisted finite-degree deduction).}
For $n=4,5$, the polynomial map
\[
 \Phi_n:(\operatorname{Mat}_n\mathbb C)^4
 \longrightarrow\operatorname{Sym}^n(\mathbb C^4)^*,
 \qquad(A_1,A_2,A_3,A_4)\longmapsto
 \operatorname{per}\!\left(\sum_{k=1}^4x_kA_k\right)
\]
is dominant. Consequently, for every $d\geq0$ and
$\lambda\vdash nd$ with $\ell(\lambda)\leq4$,
\begin{equation}\label{R12:perambient}
 m_\lambda(\operatorname{per}_n)=a_\lambda(d[n])
 \qquad(n=4,5).
\end{equation}
Combining the quartic case with G-07 and writing $D=320112$ gives
the exact comparison
\begin{equation}\label{R12:deficit}
 m_\lambda(\operatorname{per}_4)-m_\lambda(\det_4)
 =a_{\lambda-(D,D,D,D)}((d-D)[4]),
\end{equation}
with the same zero conventions as in G-07.
In particular, for every integer $s\geq0$,
\begin{equation}\label{R12:gapone}
 d=D+s,\qquad \lambda=(D+4s,D,D,D)
 \quad\Longrightarrow\quad
 m_\lambda(\operatorname{per}_4)-m_\lambda(\det_4)=1.
\end{equation}
These are multiplicity obstructions; no occurrence-vanishing
claim is implied. For $n=5$, G-06 and \eqref{R12:perambient}
give a four-row multiplicity obstruction in some degree, without
specifying that degree or highest weight.

\paragraph{Explicit rank certificate for $n=4$.}
Each entry of the following array is the coefficient quadruple of
a linear form in the ordered variables $(x_1,x_2,x_3,x_4)$:
\begin{equation}\label{R12:pencil4}
\begin{pmatrix}
(2,-1,0,0)&(2,-2,0,-2)&(-2,-1,-1,0)&(0,-1,-1,2)\\
(0,1,0,0)&(2,-2,-1,1)&(-2,1,-2,0)&(-1,-2,-1,-1)\\
(0,-2,0,2)&(1,-2,-1,-1)&(2,-1,-1,0)&(2,2,0,0)\\
(1,1,2,1)&(1,0,2,-2)&(-1,2,-1,1)&(2,1,1,1)
\end{pmatrix}.
\end{equation}
Call the resulting matrix of linear forms $L_4(x)$. Index the
parameters $a_{ijk}$ by $0\leq i,j\leq3$, $0\leq k\leq3$,
where $a_{ijk}$ multiplies $x_{k+1}$ in entry $(i+1,j+1)$,
and assign column number $c=4(4i+j)+k$.
Order the 35 exponent quadruples $\alpha$ of total degree four in
decreasing lexicographic order. Define the integer Jacobian by
\begin{equation}\label{R12:derivative}
 J_{\alpha,,4(4i+j)+k}
 =[x^\alpha]\;x_{k+1}
 \operatorname{per}\bigl(L_4(x)_{\widehat{i+1},\widehat{j+1}}\bigr).
\end{equation}
Take all 35 rows and the ordered columns
\[
 C_4=(0,1,\ldots,30,32,33,34,35).
\]
The resulting square integer matrix has determinant
\begin{equation}\label{R12:minor4}
 \det J_{C_4}
 =45459683955923995466096193465479594741071872
 \equiv77\pmod{101}.
\end{equation}
Thus \eqref{R12:pencil4}--\eqref{R12:minor4} specify the entire
nonvanishing certificate in the text, independently of a file format.

\paragraph{Explicit rank certificate for $n=5$.}
Use the same coefficient-quadruple convention with the array
\begin{equation}\label{R12:pencil5}
\begin{pmatrix}
(2,2,-1,1)&(2,-2,-2,1)&(-2,-1,1,1)&(1,-1,0,0)&(1,2,0,2)\\
(1,0,0,-1)&(1,2,0,2)&(1,2,-1,2)&(0,2,-1,1)&(-1,2,1,0)\\
(-2,-2,1,0)&(1,-1,-1,2)&(0,1,0,-1)&(2,2,-1,1)&(2,-1,1,-2)\\
(2,0,0,1)&(-1,1,2,2)&(0,0,1,-2)&(1,-1,1,-1)&(-1,0,1,-2)\\
(0,-2,1,-1)&(1,0,-1,1)&(-1,-1,1,-1)&(-1,-1,2,1)&(0,-2,-2,2)
\end{pmatrix}.
\end{equation}
There are 56 degree-five coefficient rows, again in decreasing
lexicographic order, and parameter column $4(5i+j)+k$.
Use the same derivative rule, with $4\times4$ complementary
permanents, and the columns
\[
 C_5=(0,1,\ldots,37,39,40,\ldots,56).
\]
The selected $56\times56$ minor has determinant $30$ modulo 101.

\begin{proof}
Differentiating the permanent with respect to one matrix entry
gives the complementary permanent; differentiating that entry
with respect to $a_{ijk}$ contributes $x_{k+1}$. This proves the
Jacobian rule. Expanding the complementary permanents and
performing exact modular elimination gives the displayed
nonzero residues. For full reproducibility,
\texttt{check\_permanent\_projection.py} records the full pencils,
coefficient matrices, selected columns, and integer minors.
The characteristic-zero determinants are also calculated by
fraction-free Bareiss elimination. As a second check on the
derivative expansion, every one of the 64 quartic columns was
recomputed as the exact finite difference
$\Phi_4(a+e_{ijk})-\Phi_4(a)$; this is exactly the derivative
because the map is affine in each individual parameter.
Twelve quintic columns were checked in the same way.

A nonzero minor modulo 101 is a nonzero integer minor, so the
complex differential is surjective at each displayed point.
Fix the other parameters and retain the selected 35 or 56
parameters. The resulting square polynomial map has invertible
complex differential there. By the inverse function theorem its
image contains a nonempty analytic open subset. A polynomial
vanishing on that image vanishes on an analytic open subset and
is identically zero. This proves Zariski density, hence dominance.

Every permanent pencil is a linear endomorphism substitution
into $\operatorname{per}_n$, and therefore belongs to its
$GL_{n^2}$ orbit closure. This closed variety consequently
contains the full linear space of degree-$n$ forms supported on
the first four variables. A highest-weight polynomial of weight
$\lambda$ with $\ell(\lambda)\leq4$ uses only those coefficient
variables, by the nonnegative-weight support argument of R-07.
If it vanishes on the permanent orbit closure, it vanishes on
that entire coefficient space, so it is the zero polynomial.
Thus the ideal has no highest-weight vector in those weights,
which proves \eqref{R12:perambient}.

Equation \eqref{R12:deficit} now follows by subtracting G-07.
For \eqref{R12:gapone}, the shifted partition is $(4s)$, whose
multiplicity in $\operatorname{Sym}^s(\operatorname{Sym}^4V)$ is
one: its weight space is spanned by $(e_1^4)^s$, and this vector
is highest weight. This also covers $s=0$.
For $n=5$, the four-variable determinant image has dimension at
most $52<56$ by G-06, so its ideal contains a nonzero homogeneous
highest-weight vector. G-05 excludes highest weights of length
at most three. Hence it removes at least one four-row ambient
multiplicity, whereas \eqref{R12:perambient} retains them all.
\end{proof}

These computations are exact finite certificates. The independent
agent recheck uses Ryser's formula at integer evaluation points
and an invertible interpolation matrix, so it does not reuse
the complementary-permanent coefficient-expansion algorithm.
The checks establish the stated degrees four and five; they
do not extrapolate dominance to growing $n$ or supply an
asymptotic permanent-versus-determinant lower bound.

\subsection{R-13: the exact range of quaternary permanent dominance}

\paragraph{Statement R-13 (finite certificates and a dimension obstruction).}
For the polynomial map $\Phi_n$ in R-12,
\begin{equation}\label{R13:range}
 \Phi_n\text{ is dominant over }\mathbb C
 \quad\Longleftrightarrow\quad1\leq n\leq16.
\end{equation}
Consequently, for $2\leq n\leq16$ and every
$\lambda\vdash nd$ with $\ell(\lambda)\leq4$,
\begin{equation}\label{R13:multiplicities}
 m_\lambda(\operatorname{per}_n)=a_\lambda(d[n]).
\end{equation}
For every $4\leq n\leq16$, the unpadded permanent and determinant
of size $n$ have a multiplicity obstruction with exactly four rows
in some coordinate-ring degree. R-12 gives an explicit infinite
family when $n=4$; no degree or partition is asserted here for
$5\leq n\leq16$. This is a certificate-based finite-range result,
with no claim of literature priority or asymptotic separation.

\paragraph{Complete specification of the finite certificates.}
All arithmetic for this paragraph is in $\mathbb F_{101}$, with
integer coefficient representatives in $\{0,\ldots,100\}$.
For each $n=1,\ldots,16$ set
\[
 s_0=20261121+1009n,\qquad
 s_{c+1}=(1664525s_c+1013904223)\bmod2^{32},\qquad
 a_c=\left\lfloor\frac{s_{c+1}}{2^{16}}\right\rfloor\bmod101.
\]
There are $4n^2$ coefficients, with $0\leq c<4n^2$, and the
pencil is
\[
 L_n(x)_{i+1,j+1}=\sum_{k=0}^3a_{4(ni+j)+k}x_{k+1}
 \quad(0\leq i,j<n).
\]
Put $b_n=\binom{n+3}{3}$ and order all integer points
$\alpha\in\mathbb N^4$ with $|\alpha|=n$ in decreasing
lexicographic order. Define the $b_n\times4n^2$ matrix
\begin{equation}\label{R13:evaluationmatrix}
 E_n\bigl(\alpha,4(ni+j)+k\bigr)
 =\alpha_{k+1}\operatorname{per}
   \bigl(L_n(\alpha)_{\widehat{i+1},\widehat{j+1}}\bigr)
 \pmod{101}.
\end{equation}
Choose columns greedily from left to right, retaining a column
exactly when it increases the span of previously retained columns.
Equivalently, scan columns left to right in Gaussian elimination,
choose the first available nonzero pivot row, and clear the rows
below it. This produces $b_n$ columns $C_n$ in every case, and the
square minors, with rows and columns in the stated order, have
the following nonzero residues:
\[
\begin{array}{c|r|r@{\qquad}c|r|r}
n&b_n&\det(E_n)_{C_n}&n&b_n&\det(E_n)_{C_n}\\\hline
1&4&1&9&220&67\\
2&10&11&10&286&67\\
3&20&64&11&364&42\\
4&35&49&12&455&87\\
5&56&85&13&560&25\\
6&84&28&14&680&69\\
7&120&25&15&816&73\\
8&165&59&16&969&8
\end{array}
\]
These uniform range certificates use the displayed generator;
R-12 retains its separate, smaller-entry quartic and quintic pencils.
The directory \texttt{permanent\_range} contains all generated
integer coefficients, evaluation points, chosen columns, and full
matrices, with unambiguous row-major byte encoding and checksums.
Thus the formula above and the code reproduce the certificates
without probabilistic choices or an external data source.

\begin{proof}
For the positive direction, \eqref{R13:evaluationmatrix} is the
derivative of the coefficient map followed by evaluation at the
specified points, by the complementary-permanent derivative
rule in R-12. Its full row rank $b_n$ forces the coefficient-map
differential, whose target dimension is $b_n$, to be surjective.
The nonzero minor modulo 101 is the reduction of a nonzero
integer minor. Hence the characteristic-zero differential has
the same required full rank, and the inverse-function argument
in R-12 proves dominance.

Here are the exact algorithms used to certify every entry and
minor in the table. For a numeric $n\times n$ matrix $M$, let
$F(\varnothing)=1$ and
\[
 F(S)=\sum_{j\in S}M_{|S|,j}F(S\setminus\{j\}).
\]
Induction on $|S|$ shows that $F(S)$ is the permanent formed from
the first $|S|$ rows and columns $S$. Thus $F([n])$ is the
permanent. Reverse differentiation of this recurrence computes
every entry derivative, hence every complementary permanent,
in $O(n2^n)$ field operations per evaluation point. This is
implemented in \texttt{check\_permanent\_range.cpp}.

All $3{,}348{,}864$ entries across the sixteen matrices were then
recomputed from Ryser's different identity
\[
 \operatorname{per}(M)=
 \sum_{S\subseteq[n]}(-1)^{n-|S|}
 \prod_{i=1}^n\left(\sum_{j\in S}M_{ij}\right).
\]
To prove this identity, expand the products as maps
$f:[n]\to[n]$. The coefficient of a fixed map is
$\sum_{S\supseteq\operatorname{im}f}(-1)^{n-|S|}$, which is zero
unless $f$ is onto and is one if it is onto. Such onto maps are
exactly permutations. Differentiating this identity and using
Gray-code row-sum updates gives the independent algorithm in
\texttt{check\_permanent\_range\_ryser.cpp}. Every entry agrees
with the first computation. Finally
\texttt{check\_permanent\_range\_minors.py} recomputes each
selected determinant by a second modular elimination algorithm;
all sixteen nonzero residues agree. These are finite exact
checks, rather than floating-point rank estimates.

For the converse, the torus
\[
 H_n=\{(u_1,\ldots,u_n,v_1,\ldots,v_n)\in(\mathbb C^*)^{2n}:
                  \prod_i u_i\prod_j v_j=1\}
\]
preserves $\Phi_n$ under the entrywise action
$L_{ij}\mapsto u_i v_jL_{ij}$. It has dimension $2n-1$.
On the nonempty open set where every entry $L_{ij}$ is a nonzero
linear form, a stabilizing torus element satisfies
$u_i v_j=1$ for all $i,j$. It is therefore of the form
$u_i=t$, $v_j=t^{-1}$, a one-dimensional stabilizer.
Thus each generic fiber contains an orbit of dimension $2n-2$,
and
\begin{equation}\label{R13:dimension}
 \dim\overline{\operatorname{im}\Phi_n}
       \leq4n^2-2n+2.
\end{equation}
At $n=17$, this is $1124<1140=b_{17}$. For $n\geq18$,
\[
 6\bigl(b_n-(4n^2-2n+2)\bigr)
 =n^3-18n^2+23n-6>0.
\]
Consequently no $n\geq17$ map is dominant. This proves
\eqref{R13:range}.

For $n\geq2$, the ambient variable space has dimension $n^2\geq4$,
so dominance places the entire space of degree-$n$ forms in the
first four variables inside the permanent orbit closure. The
highest-weight support argument of R-12 proves
\eqref{R13:multiplicities}.
For $n\geq4$, G-06 gives a nonzero determinant equation supported
on four variables, and G-05 excludes weights with at most three
rows. Thus some exactly four-row highest weight has determinant
multiplicity below its ambient value. Equation
\eqref{R13:multiplicities} retains that ambient value on the
permanent side when $n\leq16$, giving the asserted obstruction.
\end{proof}

The dimension ceiling concerns full ambient survival in four
variables. It does not determine the comparison between the two
proper four-variable images for $n\geq17$.

\subsection{R-14: a dimension criterion that survives genuine padding}

\paragraph{Statement R-14.}
Let $\mathcal P_{m,4}$ be the Zariski closure of the image of
$\Phi_m$ and put $c_m=\dim\mathcal P_{m,4}$.
For integers $n>m\geq2$, embed the usual padded permanent
$h_{m,n}=z^{n-m}\operatorname{per}_m$ in
$\operatorname{Sym}^n(\mathbb C^{n^2})^*$, with $z$ independent
of the $m^2$ permanent variables. The projection of its orbit
closure to the first four variables is exactly
\begin{equation}\label{R14:paddingvariety}
 \mathcal Z_{m,n}=
 \{\ell^{n-m}q:\ell\in(\mathbb C^4)^*,\ q\in\mathcal P_{m,4}\},
 \qquad \dim\mathcal Z_{m,n}=c_m+3.
\end{equation}
If
\begin{equation}\label{R14:criterion}
 c_m+3>2n^2+2,
\end{equation}
then, for every sufficiently large coordinate-ring degree $d$,
there is a partition $\lambda\vdash nd$ of length exactly four
such that
\[
 m_\lambda(h_{m,n})>m_\lambda(\det_n).
\]
No effective threshold for $d$ is asserted.

For $4\leq m\leq16$, R-13 gives $c_m=\binom{m+3}{3}$, so the
sufficient condition becomes
\begin{equation}\label{R14:explicitrange}
 m<n\leq\left\lfloor\sqrt{\frac12\binom{m+3}{3}}\right\rfloor.
\end{equation}
The nonempty padded ranges are
\[
\begin{array}{c|ccccccccc}
m&8&9&10&11&12&13&14&15&16\\\hline
n&9&10&11&12\text{--}13&13\text{--}15&14\text{--}16&
15\text{--}18&16\text{--}20&17\text{--}22
\end{array}.
\]
In particular, some four-row multiplicity obstruction separates
the unpadded $\operatorname{per}_n$ and $\det_n$ for every
$4\leq n\leq22$; for $17\leq n\leq22$ this follows by
block-diagonal containment from the $m=16$ padded cases.

\begin{proof}
First, $\mathcal P_{m,4}$ is an irreducible affine cone, as the
closure of the image of an irreducible affine space. It contains
an irreducible degree-$m$ polynomial: use the cyclic matrix with
diagonal entries $x_1$, entries $(i,i+1)$ equal to $x_2$ for
$i<m$, entry $(m,1)$ equal to $x_3$, and every other entry zero.
Its only perfect matchings are the diagonal and the full cycle,
so its permanent is $x_1^m+x_2^{m-1}x_3$.
Viewed in $\mathbb C[x_2,x_3][x_1]$, this polynomial is
Eisenstein at the prime $x_3$ and hence irreducible.
The reducible degree-$m$ forms constitute a closed subset:
projectivize each of the finitely many multiplication maps of
degrees $a$ and $m-a$, $1\leq a<m$, and take their closed images.
It follows that a general $q\in\mathcal P_{m,4}$ is irreducible.

Projectivizing the multiplication map
$(\ell,q)\mapsto\ell^{n-m}q$ gives a morphism
\[
 \mathbb P^3\times\mathbb P(\mathcal P_{m,4})
 \longrightarrow\mathbb P(\operatorname{Sym}^n\mathbb C^4)^*.
\]
Its image is closed because its domain is projective; its affine
cone is exactly the set in \eqref{R14:paddingvariety}, including
zero. For a generic irreducible $q$ of degree $m\geq2$, the
linear factor $\ell$ is uniquely determined by the product,
up to scalar, and then so is $q$. Thus the projective map is
generically injective. Its image dimension is
$3+(c_m-1)$ and the affine-cone dimension is $c_m+3$.

Every four-variable restriction of an invertible substitution
in $h_{m,n}$ has the form in \eqref{R14:paddingvariety}, so
closedness gives one inclusion of the projected orbit closure.
Conversely, for each permanent pencil $q=\Phi_m(A)$, any chosen
$\ell$ and its $m^2$ matrix-entry linear forms can be prescribed
simultaneously by an endomorphism substitution, since the
padding variable and the permanent variables are independent.
There is enough room because $m^2+1\leq n^2$.
Such substitutions lie in the orbit closure. Taking limits in
$q$ then shows that the whole closed set $\mathcal Z_{m,n}$ is
contained in the orbit closure, supported on its first four
variables. This proves equality of the projected image.

Let $\mathcal D_{n,4}$ be the determinant four-variable image
closure. By G-06, $\dim\mathcal D_{n,4}\leq2n^2+2$.
The coordinate rings of both cones are standard graded. The
Hilbert function of a nonzero standard graded affine cone of
dimension $c$ agrees for large $d$ with a polynomial of degree
$c-1$ and positive leading coefficient. Thus
\eqref{R14:criterion} implies
\[
 \dim\mathbb C[\mathcal Z_{m,n}]_d
 >\dim\mathbb C[\mathcal D_{n,4}]_d
\]
for every sufficiently large $d$.
Decompose these finite-dimensional $GL_4$ representations into
irreducibles. Their dimensions are the sums of multiplicities
times the positive integers $\dim S_\lambda\mathbb C^4$.
Hence some irreducible multiplicity on the first side is larger.

These $GL_4$ multiplicities agree with the corresponding
$GL_{n^2}$ multiplicities of length at most four in the original
orbit-closure coordinate rings. Indeed, highest-weight vectors
with those weights use only the first-four coefficient variables,
and the intersection of the full ideal with that polynomial
subring is the ideal of its projected image. The additional
raising operators involving later variables vanish on such
vectors. This proves the claimed multiplicity inequality.
G-05 states that the determinant multiplicity is already the
ambient maximum for weights of length at most three, so the
strict inequality must occur at length exactly four.

Substitution of $c_m=\binom{m+3}{3}$ into
\eqref{R14:criterion} gives $2n^2\leq\binom{m+3}{3}$, which is
\eqref{R14:explicitrange} and the displayed finite table.
Finally
\[
 \operatorname{per}\operatorname{diag}(L_m,zI_{n-m})
 =z^{n-m}\operatorname{per}(L_m)
\]
shows that the padded-permanent orbit closure is contained in
the size-$n$ permanent orbit closure. The resulting coordinate
quotient is surjective and therefore cannot increase
multiplicities. The padded gaps at $m=16$ thus imply the stated
unpadded gaps for $17\leq n\leq22$.
\end{proof}

If the expected formula
$c_m=\min\{\binom{m+3}{3},4m^2-2m+2\}$ were proved beyond the
certified range, \eqref{R14:criterion} would turn it into further
padded obstructions up to approximately $n=\sqrt2\,m$.
This expectation is a separate generic-fiber problem; neither
the finite certificates nor the dimension upper bound proves it.
The criterion above is unconditional whenever its stated
dimension lower bound has been established.

\subsection{What the attached GCT IV, VII, and VIII positivity claims supply}

The attached GCT~IV is the later Blasiak--Mulmuley--Sohoni manuscript
corresponding to their 2015 memoir, with filename
\texttt{Apr13KroneckerGCT4.tex}. Its \S12.1 theorem
(\texttt{t nonstandard schur-weyl duality}) proves nonstandard
Schur--Weyl duality for general $\dim V,\dim W$; the later
explicit-basis limitation should not obscure this proved
structural result. Its Theorem~15.21(vii)--(viii)
(\texttt{t main canonical basis}) proves a highest-weight crystal
count for $g(\alpha,\beta,\gamma)$ when
$\ell(\alpha),\ell(\beta)\leq2$, including the symmetric/exterior
split. Theorem~18.2 (\texttt{t consequence for Kronecker}) makes
the count explicit, and Corollary~18.4 gives the generating
function used in R-10. This supplies a positive rule for rectangles
of height two, or width two after simultaneous conjugation.
The construction for the full tensor power remains Conjecture~19.1
(\texttt{cj canonical basis X\^{}r}); \S19.2 describes the failed
flatness of a naive extension beyond the two-row case.

Even in the proved case, counting highest-weight elements does not
mean that every matrix coefficient is nonnegative. GCT~IV,
Example~17.8 (\texttt{ex not positive}), displays the divided-power
coefficient $q^{10}-q^4-q^{-4}+q^{-10}$. This concerns the specified
basis and operator; it is a concrete warning against replacing a
precise canonical-basis positivity statement by unrestricted
coefficient positivity.

GCT~VII (Mulmuley, 2008a, revised extended abstract) proves in
Theorem~1.2 (\texttt{tnonstd}) existence of the nonstandard
quantization, complete reducibility of its polynomial
representations, and its Peter--Weyl statement. The nonstandard
duality and reciprocity needed to recover plethysm are explicitly
Conjectures~1.3--1.4 (\texttt{cduality}, \texttt{creciprocity});
the refined form is Conjecture~6.1 (\texttt{crecirefined}).
The proposed formula
$a_{\lambda,\mu}^{\pi}=\sum_\alpha
m_\lambda^\alpha n_\pi^\alpha$ is conditional on those structures.
The introduction states that correctness of the basis algorithms
and their required properties has not been proved there.

GCT~VIII (Mulmuley, 2008b, revised extended abstract) makes the
additional gates explicit: local crystal existence is Conjecture~2.2
(\texttt{clocalcrystalb}); minimum-degree globalization is
Conjectures~2.5, 3.1--3.2; cellular decomposition is Conjectures~2.6,
3.3 and the specialization-compatible quasi-cell condition is
Conjecture~3.7. Kronecker positivity is Conjectures~2.7 and~3.4,
with specified signs and powers of $q-q^{-1}$ allowed. The general
plethysm setting uses the weaker saturation and nonstandard
positivity Conjectures~2.8--2.9 and~3.5--3.6. Finally, \S6
(\texttt{scomplexity}) also requires efficient combinatorial
labels and crystal operations to obtain a $\#P$ formula; these
complexity properties are themselves proposed rather than proved.

Its discussion of the Riemann hypothesis in the introduction and
\S2.4.2 is a suggested extension of finite-field purity and
perverse-sheaf machinery to the nonflat, noncommutative objects at
issue. It is not a theorem deriving these conjectures from the
classical RH or Dirichlet GRH. Accepting the stipulated analytic
theorems therefore leaves specific algebraic gates: construct the
required realization, prove its cellular and specialization
properties, and control the sizes of its labels and operations.
The supplied papers prove the height-two case and propose this
general route; they supply no unconditional constructive positive
formula for arbitrary rectangles and no theorem controlling the
orbit-closure ranks in R-05--R-06.

\paragraph{Algebraic-geometric dictionary for R-01--R-08.}
R-01 and R-02 are identities of Schur operations and restriction maps
on representation $K_0$, evaluated over $\mathbb C$ for Haar integration.
The tensor homomorphism is the integral group-scheme morphism
$GL_a\times GL_n\to GL_{an}$, and its kernel contains the central
$\mathbb G_m$ described in R-03. R-04 is the cycle-index expression
for symmetric powers in a $\lambda$-ring; the Adams operations preserve
the representation ring, not its cone of actual representations.
For R-05--R-06 the intrinsic algebraic object is the graded orbit-map
image, or the coefficient map in \eqref{R06:coefficients}. Its rank
strata are defined by determinantal (Fitting) ideals. The Gaussian
and Hermitian metrics are analytic devices for the same rank; they
are not asserted to be motives. The Schur idempotent is defined over
$\mathbb Z[1/D!]$, and the integral matrix in \eqref{R06:integer}
avoids treating characteristic-zero semisimplicity as an integral
direct-sum theorem. R-07 is a statement about restrictions of locally
free Schur constructions to the binary-form subscheme and the
scheme-theoretic image of the linear-substitution morphism. Its
formulas can be cleared of their binomial denominators over
$\mathbb Z$, while the basis and orbit-closure claims here are in
characteristic zero. R-08 is additivity in equivariant $K_0$ applied
to a resolution; the numerator remembers the quotient's equations.
Exact SGA/EGA/FGA and cote~133 locators for this dictionary are supplied
in the separate Grothendieck part of the paper.
R-09 identifies a four-factor invariant of
$\operatorname{Sym}(V_1\otimes V_2\otimes V_3\otimes V_4)$ with the
composition of two tensor-product restriction functors; its degree
$nd$ determinant-isotypic coefficient is the fourfold Kronecker
multiplicity. The finite CUE cutoff is the rank restriction on a
Schur functor, not an arithmetic purity hypothesis.

% END SOURCE work/rectangular/findings.tex

% BEGIN SOURCE work/root/findings.tex
\section{Integral and $\lambda$-ring controls on the bridges}\label{sec:root-findings}
The results in this section are deductions checked in this run. They are not claims of historical priority. We distinguish an identity of characters, a decomposition in a Grothendieck group, an actual direct-sum decomposition, and a dimension identity.

\begin{proposition}[G-01: The three meanings of the second operation]\label{G-01}
Let $E$ be a finite $G$-set. Write $b_2(E)$ for the $G$-set of two-element subsets and $s_2(E)$ for the $G$-set of two-element multisets. In the Burnside ring one has
\[
 [s_2(E)]=[b_2(E)]+[E],\qquad [E]^2=[E]+[\operatorname{Inj}(\{1,2\},E)].
\]
The identity $[E]^2=2[b_2(E)]+[E]$ does not hold in general. If $V$ is a representation over a field of characteristic zero, the correct special $\lambda$-ring identities are
\[
 [\operatorname{Sym}^2V]=[V]^2-[\wedge^2V],\quad
 [\operatorname{Sym}^2V]-[\wedge^2V]=\psi^2[V],\quad
 [V]^2=2[\wedge^2V]+\psi^2[V].
\]
Consequently the expression $2\lambda^2(E)+E=E^2$ in cote 133, scan 019, line 53 is valid after taking cardinalities, but is not a universal equivariant or special-$\lambda$-ring identity. The supplied study's assertion that the two equivariant descriptions agree needs this correction.
\end{proposition}
\begin{proof}
Partition a multiset of cardinality two into the diagonal multisets $\{e,e\}$ and the subsets $\{e,e'\}$ with $e\ne e'$. This partition is $G$-equivariant. Similarly partition ordered pairs into diagonal and off-diagonal pairs. For a counterexample take $G=C_2$ and its regular two-element set $E$. Then $E\times E$ is a disjoint union of two free $C_2$-orbits, whereas $b_2(E)$ is one fixed point. Thus the asserted equality would read $2[E]=[E]+2[\mathrm{pt}]$. Applying the fixed-point homomorphism for the whole group gives $0=2$, a contradiction.
For representations, first prove the universal identity on $\mathrm{GL}_r$ and then pull it back along the representation of $G$. Write the character of its defining representation on the diagonal torus as $x_1+\cdots+x_r$. The symmetric and exterior squares have characters $\sum_{i\leq j}x_ix_j$ and $\sum_{i<j}x_ix_j$. Their sum is $(\sum_i x_i)^2$, and their difference is $\sum_i x_i^2=\psi^2[V]$. Character injectivity in characteristic zero proves the formulas. In the same $C_2$ example the permutation representation is $1\oplus\epsilon$, and $\wedge^2V=\epsilon$, whereas linearization of $b_2(E)$ is $1$. This explicitly shows why the subset operation cannot be silently identified with exterior power. Taking dimensions sends $\psi^2[V]$ to $\dim V$ and recovers the numerical identity. None of this determines the ink on a scan which has not been supplied.
\end{proof}

\begin{proposition}[G-02: The terminal exterior-power variable is necessary]\label{G-02}
For $i\geq2$, a universal formula for $\lambda^i(xy)$ in special $\lambda$-rings cannot in general omit $\lambda^i(x)$ or $\lambda^i(y)$. Thus the terminal index in cote 133, scan 006, lines 25--26 must be $i$ if those lines state the universal product formula.
\end{proposition}
\begin{proof}
Set $y=1$, for which $\lambda^1(1)=1$ and $\lambda^j(1)=0$ for $j>1$. A formula using only $x,\lambda^2(x),\ldots,\lambda^{i-1}(x)$ would express $\lambda^i(x)$ in the subring generated by these earlier operations. Evaluate at the defining representation of $\mathrm{GL}_i$ and restrict to its diagonal torus. This would express the elementary symmetric polynomial $e_i$ in $\mathbb Z[e_1,\ldots,e_{i-1}]$, contrary to the algebraic independence of $e_1,\ldots,e_i$ (equivalently the fundamental theorem of symmetric polynomials). Interchange $x$ and $y$ for the other endpoint. The argument establishes the required mathematical reading, not a new visual transcription.
\end{proof}

\begin{proposition}[G-03: Finite $\gamma$-data do not determine multiplicities]\label{G-03}
For every integer $D\geq0$ there are two honest representations $A_D,B_D$ of the split group scheme $\mathbb G_{m,\mathbb Z}$, each of rank $2^D$, which have the same class modulo $F^{D+1}_\gamma$ but different weight multiplicities. Over $\mathbb Q$ their Chern characters agree in degrees $0,\ldots,D$.
\end{proposition}
\begin{proof}
Let $L$ denote the standard character and put $r=D+1$. Define
\[
 A_D=\bigoplus_{\substack{0\leq j\leq r\\j\text{ even}}}L^{\otimes j\,\oplus\binom rj},\qquad
 B_D=\bigoplus_{\substack{0\leq j\leq r\\j\text{ odd}}}L^{\otimes j\,\oplus\binom rj}.
\]
The two binomial parity sums are $2^{r-1}$. In the representation ring $\mathbb Z[L,L^{-1}]$ the difference of the classes is $(1-L)^r$, a product of $r$ rank-zero elements and hence an element of $F^r_\gamma$. The constant weight has multiplicity one in $A_D$ and zero in $B_D$, so the characters differ. With $c=c_1(L)$, the difference of the Chern characters is $(1-e^c)^r$, whose lowest possible degree is $r$. This proves the assertion. For a line bundle on a scheme, applying the associated-bundle map gives the identical statement in its $K_0$. This does not assert that the two resulting bundles are nonisomorphic on every scheme.
\end{proof}

\begin{proposition}[G-04: An integral coefficient model and its exact specialization gate]\label{G-04}
Let $W_{\mathbb Z}$ be a finite free module, let $\rho:\mathrm{GL}_{N,\mathbb Z}\to\mathrm{GL}(W_{\mathbb Z})$ be a homogeneous polynomial representation of positive degree $e$, and let $f\in W_{\mathbb Z}$. Define $X_{\mathbb Z}$ as the affine schematic closure of its orbit. Its coordinate algebra is torsion-free, hence flat over $\mathbb Z$. In degree $d$ its orbit-coefficient map is an explicit finite integer matrix
\[
 A_d:\operatorname{Sym}^d(W_{\mathbb Z}^*)\longrightarrow
 \mathbb Z[g_{11},\ldots,g_{NN}]_{\leq ed},\qquad P\longmapsto P(\rho(g)f).
\]
Its rank over $\mathbb Q$ is the degree-$d$ coordinate-ring dimension. If $\delta_d$ is any nonzero maximal-rank minor of $A_d$, then this dimension and the degree-$d$ kernel commute with base change over $\mathbb Z[1/\delta_d]$. At a prime dividing $\delta_d$, rank must be tested anew; the fiber of the integral closure need not be the schematic orbit closure of the reduced point.
\end{proposition}
\begin{proof}
The coordinate homomorphism of the orbit map has target $\mathbb Z[\mathrm{GL}_N]=\mathbb Z[g_{ij},\det(g)^{-1}]$, a torsion-free $\mathbb Z$-module. Its kernel $I$ defines the schematic closure and its image $\mathbb Z[W]/I$ is torsion-free. A module over the principal ideal domain $\mathbb Z$ is flat if it is torsion-free, proving the first assertion. Since the representation is polynomial, each degree-$d$ homogeneous coordinate polynomial maps to a polynomial of degree at most $ed$ in the entries of $g$. Equality on invertible matrices is equality in the polynomial ring because localization by the nonzero polynomial $\det(g)$ is injective. This proves the finite coefficient-matrix description.
Let $r=\operatorname{rank}_{\mathbb Q}A_d$. All $(r+1)$-minors vanish as integers, while an $r$-minor $\delta_d$ is nonzero. Over $\mathbb Z[1/\delta_d]$, use that invertible block for row and column elimination. The vanishing of larger minors makes the remaining Schur complement zero, so the matrix is equivalent to $\operatorname{diag}(I_r,0)$. Its image and kernel are direct summands there and commute with every base change. Without inversion, an inclusion of lattices need not remain injective modulo a prime: already multiplication by $p$ on $\mathbb Z$ has rank one over $\mathbb Q$ and rank zero modulo $p$. Thus no unsupported all-primes orbit-fiber conclusion follows. Homogeneity sends degree $d$ to degree $ed$, so the kernel is homogeneous and the stated grading is well defined. If $r=0$, take $\delta_d=1$.
\end{proof}

\begin{proposition}[G-05: Ternary saturation and the three-row exclusion]\label{G-05}
Work over $\mathbb C$, and let $n\geq2$, $N=n^2$, and $\mathrm{Det}_n\subseteq\operatorname{Sym}^n(\mathbb C^N)^*$ be the determinant orbit closure. Then $\mathrm{Det}_n$ contains every degree-$n$ form supported on a fixed three-dimensional subspace. Consequently, for every partition $\lambda\vdash nd$ with $\ell(\lambda)\leq3$,
\[
 \operatorname{mult}_{S_\lambda\mathbb C^N}\mathbb C[\mathrm{Det}_n]_d
 =\operatorname{mult}_{S_\lambda\mathbb C^N}\operatorname{Sym}^d(\operatorname{Sym}^n\mathbb C^N).
\]
No multiplicity obstruction with at most three rows can separate any other orbit closure of a degree-$n$ form in the same ambient space, including a padded permanent, from $\mathrm{Det}_n$.
\end{proposition}
\begin{proof}
We use the established linear determinantal representation theorem for smooth plane curves: a smooth plane degree-$n$ curve together with a line bundle $L$ of degree $g-1$ with $H^0(L)=0$ admits an $n\times n$ matrix of linear forms whose determinant is its defining form up to a nonzero scalar (Beauville, 2000, Proposition 3.1(a)). Such an $L$ exists over $\mathbb C$: for $g\geq1$, the effective locus in $\operatorname{Pic}^{g-1}$ is the image of the symmetric power $C^{(g-1)}$, of dimension at most $g-1$, whereas $\operatorname{Pic}^{g-1}$ has dimension $g$; for $g=0$, take the degree-$-1$ line bundle. Scaling one row absorbs the nonzero scalar. Hence every smooth ternary degree-$n$ form is obtained by a linear substitution in $\det_n$. Any endomorphism substitution is a limit of invertible substitutions because $\mathrm{GL}_N$ is Zariski dense in $\operatorname{End}(\mathbb C^N)$; continuity of the polynomial action therefore puts these forms in $\mathrm{Det}_n$. Smooth ternary forms form a dense open set, so its closedness puts every ternary form there.
Let $W=\operatorname{Sym}^d(\operatorname{Sym}^n\mathbb C^N)$, and choose the Borel convention whose highest-weight vector has nonnegative weight $\lambda=(\lambda_1,\lambda_2,\lambda_3,0,\ldots,0)$. In a monomial basis of $W$, zero weight in positions $4,\ldots,N$ forces every coefficient variable to involve only the first three basis vectors. Thus the entire highest-weight space $W_\lambda^{\mathrm{hw}}$ is contained in $\operatorname{Sym}^d(\operatorname{Sym}^n\mathbb C^3)$. A nonzero polynomial in these coefficient variables cannot vanish on all ternary forms. By the first paragraph, the degree-$d$ determinant ideal therefore meets $W_\lambda^{\mathrm{hw}}$ trivially. Exactness and complete reducibility in characteristic zero give the asserted multiplicity equality. Every other closed subvariety has a degree-$d$ coordinate ring which is a quotient of $W$, so its multiplicity cannot exceed this common ambient value. This proves the exclusion. The assertion is about complex characteristic-zero multiplicities, not an unqualified statement in positive characteristic.
\end{proof}

\begin{proposition}[G-06: The four-variable gate is nonvacuous]\label{G-06}
For $n\geq4$, the closure $\mathcal D_{n,4}$ of degree-$n$ quaternary forms which are determinants of $n\times n$ matrices of linear forms has dimension at most $2n^2+2$. Since
\[
 \binom{n+3}{3}-(2n^2+2)=\binom{n-1}{3}>0,
\]
this is a proper subvariety of the space of quaternary forms. The determinant ideal in $\operatorname{Sym}(\operatorname{Sym}^n\mathbb C^{n^2})$ therefore has a nonzero isotypic component with a partition of exactly four rows in some finite degree.
\end{proposition}
\begin{proof}
Consider the polynomial map
\[
 (A_1,A_2,A_3,A_4)\longmapsto\det(x_1A_1+x_2A_2+x_3A_3+x_4A_4)
\]
from an affine space of dimension $4n^2$. The group
$H=\{(P,Q)\in\mathrm{GL}_n^2:\det(P)=\det(Q)\}$ acts by $A_i\mapsto PA_iQ^{-1}$, preserving this map. It has dimension $2n^2-1$. Its generic stabilizer has dimension one. Indeed, for a tuple with $A_1=I$, a diagonal $A_2$ with pairwise distinct eigenvalues, and an $A_3$ with every off-diagonal entry nonzero, a stabilizer has $Q=P$; commuting with $A_2$ makes $P$ diagonal, and commuting with $A_3$ makes it scalar. The scalar pairs $(cI,cI)$ always stabilize every tuple. Upper semicontinuity of stabilizer dimension, or the rank calculation for the commutation equations on this exhibited open set, gives generic dimension one. Thus generic fibers contain orbits of dimension $2n^2-2$, and the dimension-of-fibers theorem yields the stated image bound. The displayed identity follows by expanding the binomial coefficient.
Let $\pi$ set all but the first four variables to zero. For every form in $\mathrm{Det}_n$, its image under $\pi$ lies in $\mathcal D_{n,4}$: approximate the form by invertible substitutions of $\det_n$, apply $\pi$, and use the closedness of $\mathcal D_{n,4}$. Choose a nonzero homogeneous equation $P$ of this proper cone. Then $P\circ\pi$ is a nonzero homogeneous equation of $\mathrm{Det}_n$, using only first-four coefficient variables. The ideal of $\mathcal D_{n,4}$ is $\mathrm{GL}_4$-stable; choose a highest-weight vector in a nonzero irreducible submodule of its degree-$d$ part. Embedded in the full coordinate ring, it is also a $\mathrm{GL}_{n^2}$ highest-weight vector: raising operators involving a variable beyond the fourth annihilate it. Its highest weight has at most four rows, and G-05 rules out at most three. Hence its length is exactly four. Existence of such an equation does not establish that it is nonzero on the padded permanent.
\end{proof}

\begin{proposition}[G-07: An exact four-row determinant-closure formula]\label{G-07}
Set $D=320112$. For every $d\geq0$ and partition $\lambda\vdash4d$ with $\ell(\lambda)\leq4$, padded by zeros to length four, the degree-$d$ determinant-closure multiplicity satisfies
\[
 m_{\mathrm{Det}_4}(\lambda,d)
 =a_\lambda(d[4])-a_{\lambda-(D,D,D,D)}((d-D)[4]).
\]
The second term is defined to be zero if $d<D$ or $\lambda_4<D$. In particular the four-row sector is ambient for $d<D$, and at $d=D$ the multiplicity of $\lambda=(D^4)$ is exactly one less than its ambient plethysm multiplicity.
\end{proposition}
\begin{proof}
Leal, Lozano Huerta, and Vite (2024), Theorem 2 and Table 2, identify the closure of linear determinantal quartic surfaces as an irreducible reduced hypersurface of degree $D$ in $\mathbb P(\operatorname{Sym}^4\mathbb C^4)^*$. Table 2 identifies this component $F_1$ by $a=(5,5,5,5)$, $b=(6,6,6,6)$, so every matrix entry has degree one. Its affine cone has a prime principal ideal $(F)$ in the polynomial ring $A=\operatorname{Sym}(\operatorname{Sym}^4\mathbb C^4)$, with $F$ homogeneous of degree $D$.
The line $\mathbb CF$ is stable under $\mathrm{GL}_4$. Every algebraic one-dimensional character of $\mathrm{GL}_4$ is a power of the determinant. Since scalar $zI$ acts on $A_D$ by $z^{4D}$, this character is $\det^D$. Thus multiplication by $F$ gives the exact sequence of $\mathrm{GL}_4$-modules
\[
 0\longrightarrow \det^D\otimes A_{d-D}\xrightarrow{\ \cdot F\ }A_d
 \longrightarrow(A/(F))_d\longrightarrow0.
\]
For $d<D$ its first term is zero.
Let $\pi$ restrict forms in sixteen variables to their first four variables. The image $\pi(\mathrm{Det}_4)$ is precisely the affine cone just described. One inclusion follows by continuity from the orbit of $\det_4$. For the converse, every four-variable linear determinant is an endomorphism substitution in $\det_4$, hence belongs to its orbit closure; the same is true of every limit of such forms, and such a form is fixed by $\pi$. It follows that the intersection of the full determinant ideal with the first-four coefficient polynomial ring is exactly $(F)$. Highest-weight spaces for partitions of at most four rows use only these coefficient variables, by the weight argument of G-05. Applying complete reducibility to the exact sequence and shifting a highest weight by the determinant character proves the formula. At $d=D$ the source is $\det^D\otimes\mathbb C$, so the deficit is one.
\end{proof}
\paragraph{Source correction for G-07.}
The attached Landsberg source, section \texttt{zhanglee}, theorem \texttt{zlthm}, records $640224$. The later result used here is the degree of the reduced prime divisor. Leal--Lozano Huerta--Vite, section 4.1 immediately before Theorem 2, explicitly gives $P_{20,2}=\tfrac12D_{23,14}$. This is the precise cycle-multiplicity distinction; using the unreduced Noether--Lefschetz coefficient as the polynomial degree doubles the answer. This correction concerns the cited literature, not a purported visual correction of cote 133.

\begin{proposition}[G-08: A uniform low-variable determinantal construction]\label{G-08}
Let $m\geq1$ and $k\geq2$, put $a=\lfloor k/2\rfloor$, $b=\lceil k/2\rceil$, and
\[
 B(m,k)=\binom{m+a}{a}+\binom{m+b}{b}-1.
\]
Over any commutative ring, every degree-$m$ homogeneous polynomial $q$ in $k$ variables admits a matrix of linear forms of size $B(m,k)$ with determinant
\[
 \ell^{B(m,k)-m}q,
\]
where $\ell$ may first be an independent variable and then be specialized to any linear form. The construction is polynomial in the coefficients of $q$ and defined over $\mathbb Z$.
Consequently, over $\mathbb C$, whenever $n\geq B(m,k)$, no multiplicity obstruction of length at most $k$ can separate the orbit closure of an arbitrarily padded degree-$m$ form $\ell^{n-m}f$ from $\mathrm{Det}_n$.
\end{proposition}
\begin{proof}
Split the variables into disjoint lists $x=(x_1,\ldots,x_a)$ and $y=(y_1,\ldots,y_b)$, and write
\[
 q=\sum_{|\alpha|+|\beta|=m}c_{\alpha\beta}x^\alpha y^\beta.
\]
Make a rooted directed tree with one vertex for every $x$-monomial of degree at most $m$. The parent of a nonconstant monomial is obtained by deleting one occurrence of its largest-index variable; orient the edge away from the constant root and label it by that variable. There is exactly one root-to-vertex path, with label product $x^\alpha$. Make the corresponding $y$-tree with every edge reversed, so the unique path from vertex $\beta$ to its constant root has product $y^\beta$. Connect vertex $\alpha$ in the first tree to vertex $\beta$ in the second by an edge of constant label $c_{\alpha\beta}$ exactly when $|\alpha|+|\beta|=m$. Every source-to-sink path uses exactly one connecting edge; its label product is $c_{\alpha\beta}x^\alpha y^\beta$. Therefore its path polynomial is exactly $q$. The graph is acyclic and has
\[
 s=\binom{m+a}{a}+\binom{m+b}{b}=B(m,k)+1
\]
vertices.
Choose a topological ordering with source first and sink last, and let $A$ be its strictly upper triangular weighted adjacency matrix. Then $\det(I-A)=1$ and $(I-A)^{-1}=I+A+\cdots+A^{s-1}$. The $(1,s)$ entry is the source-to-sink path polynomial $q$. The adjugate identity therefore expresses $q$, up to the fixed sign $(-1)^{s+1}$, as the determinant of the $(s-1)\times(s-1)$ minor of $I-A$ obtained by deleting row $s$ and column $1$. Change the sign of one row if needed. Its entries are affine linear in $x,y$.
Replace each constant entry $c$ of this minor by $c\ell$. Its determinant is the degree-$B(m,k)$ homogenization of $q$, hence is $\ell^{B(m,k)-m}q$. This identity is universal over $\mathbb Z[c_{\alpha\beta},x,y,\ell]$. Padding the matrix with $n-B(m,k)$ diagonal entries $\ell$ proves the corresponding size-$n$ statement.
For the multiplicity assertion, restrict any invertible substitution of $\ell^{n-m}f$ to $k$ coordinate variables. The result has the form $L^{n-m}q$, with $q$ a homogeneous degree-$m$ form in these variables. The construction puts this result in $\mathrm{Det}_n$. Continuity gives the same inclusion for the entire padded orbit closure. If $P$ is any highest-weight polynomial of length at most $k$ vanishing on $\mathrm{Det}_n$, it uses only coefficient variables supported on the first $k$ underlying variables (in the sense made precise in G-05), and therefore vanishes on the other orbit closure as well. Thus the highest-weight space of the determinant ideal embeds in that of the padded-form ideal. Subtracting their dimensions from the common ambient highest-weight multiplicity gives $m_{\mathrm{padded}}(\lambda,d)\leq m_{\mathrm{Det}_n}(\lambda,d)$ in every degree.
\end{proof}
\paragraph{Quantified implication.}
For $k=4$, the construction gives $B(m,4)=(m+1)(m+2)-1$. More generally $B(m,k)=O_k(m^{\lceil k/2\rceil})$. Thus for a fixed number of rows this sector is excluded once the padding exceeds the displayed bound. This is an explicit restricted lower-bound obstruction theorem, not a general separation result and not an efficient solution of arbitrary plethysm.

\begin{proposition}[G-09: A sharper low-row exclusion by shared exponent chains]\label{G-09}
Let $m\geq1$ and $k\geq3$. Define the following integer bounds. If $k=2r+1$ is odd, set
\[
 C(m,k)=2\binom{m+r}{r}-1.
\]
If $k=2r+2\geq4$ is even, for $1\leq h\leq m+1$ put $Q=\lfloor m/h\rfloor$, $s_0=\min(h-1,m)$ and
\[
 C_h(m,k)=\sum_{q=0}^{Q}\binom{m-qh+r}{r}+Q(h-1)
          +\sum_{s=0}^{s_0}\binom{m-s+r}{r}-1,
 \qquad C(m,k)=\min_{1\leq h\leq m+1}C_h(m,k).
\]
Over a field, for every nonzero linear form $L$ and homogeneous degree-$m$ form $q$ in $k$ variables, $L^{C(m,k)-m}q$ has a linear determinantal representation of size $C(m,k)$. In coordinates with $L=x_0$ the construction is defined over $\mathbb Z$ and polynomial in the coefficients of $q$. For fixed $k\geq3$,
\[
 C(m,k)=O_k\bigl(m^{(k-1)/2}\bigr).
\]
Consequently, over $\mathbb C$, if $n\geq C(m,k)$, there is no multiplicity obstruction with at most $k$ rows separating any padded degree-$m$ form from $\mathrm{Det}_n$. This strengthens G-08 in the specified range.
\end{proposition}
\paragraph{Prior-result boundary.}
The asymptotic determinantal-representation order used here is established: Boralevi, van Doornmalen, Draisma, Hochstenbach, and Plestenjak (2017), Theorem 1.3, prove order $m^{v/2}$ for a fixed number $v\geq2$ of affine variables, with a matching generic lower bound. Their proof in section 4 uses connected monomial spaces and an even/odd binary-digit split. We do not claim that order as new. The proof below gives the displayed finite integer bound by a different shared-chain construction, records its coefficient-ring scope, and proves the resulting GCT multiplicity exclusion. These are the forms needed here that the cited theorem does not state.
\begin{proof}
Extend $L$ to a basis of linear forms and write it as $x_0$. Dehomogenize $q$ by setting $x_0=1$, obtaining an arbitrary polynomial $p$ of degree at most $m$ in $k-1$ variables. It suffices to construct an affine linear determinant of size $C(m,k)$ equal to $p$. Homogenizing each matrix entry by $x_0$ then gives $x_0^{C(m,k)-m}q$.

For odd $k=2r+1$, split the dehomogenized variables into lists $x$ and $y$ of $r$ variables each. Use the two monomial trees from G-08, each with $\binom{m+r}{r}$ vertices, and connect every pair $(\alpha,\beta)$ with $|\alpha|+|\beta|\leq m$ by its coefficient in $p$. Every source-to-sink path is unique for its two monomial indices, so its path polynomial is $p$. The cofactor construction in G-08 has size one less than the total number of vertices, giving the displayed bound.

For even $k=2r+2$, write the variables as $x_1,\ldots,x_r,y_1,\ldots,y_r,z$ and fix $h$. Decompose the exponent of $z$ uniquely as $e=qh+s$, $0\leq s<h$. On the left, make a directed chain with $Qh$ edges, all labelled $z$, beginning at the source. At its vertex of index $qh$ attach the outward $x$-monomial tree of degrees at most $m-qh$, identifying its constant root with that chain vertex. Every such tree vertex indexed by $(q,\alpha)$ has a unique source path, with product $z^{qh}x^\alpha$. Chain vertices whose indices are not multiples of $h$ are transit vertices and receive no connecting edges. The left graph has
\[
 (Qh+1)+\sum_{q=0}^Q\left[\binom{m-qh+r}{r}-1\right]
 =\sum_{q=0}^Q\binom{m-qh+r}{r}+Q(h-1)
\]
vertices.

On the right, make a chain with vertices $0,\ldots,s_0$, directed toward vertex $0$, the sink, and again label its edges by $z$. At vertex $s$ attach the inward $y$-monomial tree of degrees at most $m-s$, identifying the root with that vertex. Every tree vertex indexed by $(s,\beta)$ has a unique sink path with product $y^\beta z^s$. The right graph has $\sum_{s=0}^{s_0}\binom{m-s+r}{r}$ vertices. Connect left vertex $(q,\alpha)$ to right vertex $(s,\beta)$ with the constant coefficient of $x^\alpha y^\beta z^{qh+s}$ in $p$ when $|\alpha|+|\beta|+qh+s\leq m$.

All connecting edges go from left to right. Thus every source-to-sink path uses exactly one such edge, and the uniqueness of the two tree paths and the division $e=qh+s$ gives exactly the required monomial with its coefficient, once. The path polynomial is $p$. Taking the same cofactor of $I-A$ as in G-08 yields an affine linear determinant of size $C_h(m,k)$. This graph has a path capable of realizing $z^m$ and a connecting edge, so its number of vertices is at least $m+2$; hence $C_h(m,k)\geq m+1$. The odd bound is also greater than $m$. Homogenization therefore uses a nonnegative exponent. Minimizing over the finite choices of $h$ proves the stated representation bound.

For the asymptotic statement in the even case take $h=\lceil\sqrt m\rceil$. The number of binomial summands is at most $m/h+1+h\leq2\sqrt m+2$, each bounded by $\binom{m+r}{r}=O_r(m^r)$, while $Q(h-1)\leq m$. Thus $C(m,k)=O_r(m^{r+1/2})$. The odd case is immediate from its binomial formula. These exponents are $(k-1)/2$ in both cases.

For the GCT consequence restrict a substituted padded form to any fixed $k$ underlying variables. If its restricted padding $L$ is nonzero, apply the determinant construction and pad with $n-C(m,k)$ diagonal copies of $L$. If $L=0$, the form is zero because $n\geq C(m,k)>m$. Thus every restriction lies in $\mathrm{Det}_n$. Apply the highest-weight ideal argument from G-08, which uses exactly this containment, to obtain the multiplicity inequality in every degree.
\end{proof}
\paragraph{Integral scope and growth of the remaining gate.}
With $L=x_0$ the graph and cofactor identity are universal over the coefficient ring. Extending an arbitrary nonzero $L$ to coordinates is a field operation; over a general base one instead works on a chart where a coefficient of $L$ is invertible and chooses the resulting local basis. No global polynomial choice of that basis is asserted. If $n=m^c$ with fixed $c>0$, every fixed $k\geq3$ satisfying $(k-1)/2<c$ is excluded for sufficiently large $m$. This is a quantitative restriction on possible multiplicity obstructions, not a claim that the remaining higher-row obstruction exists.

% END SOURCE work/root/findings.tex

% BEGIN SOURCE work/isolation/findings.tex
% CUE_GCT isolation track. Run CUE-GCT-20261121-I.
% Model: GPT-6 Astra Pro; configured effort is not exposed to this agent.
% Bibliographic keys required: Hoza2026, AGKS2014, ForbesShpilka2013, MulmuleyGCTV, AllenderOgihara1996.

\section{Isolation: the exact transfer, the quotient, and the surviving gate}
\label{sec:isolation-track}

We use a field $F$, $n$ variables, individual degree at most $\delta$, and
an ROABP expression
\[
 f(\mathbf x)=u^{\mathsf T}A_1(x_{\pi(1)})\cdots A_n(x_{\pi(n)})v.
\]
The order $\pi$ may be unknown.  A polynomial-size hitting set here means
one set depending only on $(n,w,\delta,F)$, not on the matrices or their
coefficients, which hits every nonzero member of the class.  Arithmetic
complexity and bit complexity are distinguished throughout.

\paragraph{Background I-B01.}
The stipulated new graph theorem is Hoza (2026), Theorem 1.3, pp.~2--3;
its quantitative version is Theorem 2.4, p.~5.  For every simple layered
digraph of width $W$ and length $N$, some seed in a family of size
$(WN)^{O(1)}$ gives a minimum-weight path unique for every reachable pair.
The weights are nonnegative integers bounded by $B=(WN)^{O(1)}$ and are
computable in logarithmic space.  A path weight sums its internal vertex
weights.  With a binary code $C$ of length $\ell$ and distance $d$, the
formula is
\[
 \rho_x(a,i)=6^{\nu_2(i)}\sum_{j=1}^{\ell}6^{\ell-j}C(a)_j x_j
 \quad(i>0),\qquad \rho_x(a,0)=0,
\tag{I.1}\label{eq:I-hoza}
\]
and the probability of a bad uniform seed is at most
$2^{-d}W^4N^2$.  The graph is fixed before the seed is chosen.
We use this result as supplied; the present work does not claim an
independent verification of its complete theorem.

\paragraph{Background I-B02.}
Agrawal, Gurjar, Korwar, and Saxena (2014), \S3, Definition~5, call
$\omega\in\mathbb Z_{\geq0}^{n}$ basis isolating for
$D=\sum_a v_a\mathbf x^a$ if there is a coefficient basis indexed by $B$
whose monomial weights are distinct, and each $v_a$ with $a\notin B$
belongs to the span of basis coefficients of strictly smaller weight.
Their Lemma~6 preserves every nonzero scalar projection under
$x_i=t^{\omega_i}$; Lemma~8 and Theorem~1 give their quasipolynomial
construction, including unknown variable order.  Those known results
are cited, not claimed as findings here.

\subsection{I-01. A genuine polynomial transfer for vertex-variable path polynomials}
\label{I-01}

\paragraph{Statement.}
Let $G$ range over simple layered digraphs with fixed vertex slots,
width $W$, length $N\ge1$, fixed source $s$ and target $t$, and arbitrary
scalar edge labels $c_e\in F$.  Give every internal vertex slot its own
variable $z_v$, and form
\[
 P_G(\mathbf z)=\sum_{p:s\leadsto t}\left(\prod_{e\in p}c_e\right)
                        \prod_{v\in\operatorname{int}(p)}z_v.
\]
There is a uniform hitting set of cardinality at most
\[
 2^{\ell}\bigl((N-1)B+1\bigr)=(WN)^{O(1)}
\tag{I.2}
\]
over any field containing $(N-1)B+1$ distinct elements.  Its points have
coordinates $z_v=a^{\rho_x(v)}$, where $x$ runs over Hoza's seeds and
$a$ runs over a fixed set of $(N-1)B+1$ field elements.  The same conclusion
holds for independent edge variables after subdividing edges; the new
width is at most $W^2$ and the new length is $2N$.

\paragraph{Proof.}
Delete zero-labelled edges.  In a simple layered digraph two different
$s$--$t$ paths have different sequences of internal vertices and hence
different vertex monomials, except in the length-one case, where there
is at most one edge.  Thus $P_G\ne0$ exactly when some path survives.
Choose a min-isolating seed for the surviving graph using I-B01.  In
$P_G(t^{\rho_x(v)})$, the coefficient of the least path weight is the
product of labels on its unique minimum path, and is nonzero in the
field.  The degree is at most $(N-1)B$.  A nonzero polynomial of that
degree has at most $(N-1)B$ roots.  This proves (I.2), uniformly in the
edge labels.  For edge variables, insert one distinct intermediate
vertex on each edge and give it that edge's variable; give the original
internal vertices auxiliary independent variables as well.  The same
argument provides a hitting set for the expanded vertex-variable
polynomial.  To obtain points only on edge variables, assign to an edge
the product of its intermediate-vertex value and the value of its head
when the head is an internal original vertex.  Each original path then
receives exactly the product of all its expanded internal-vertex values.
This combines exponents additively, leaves a polynomial degree bound,
and proves the edge-variable statement. \hfill$\square$

\paragraph{Scope.}
This is an exact consequence of the new graph theorem.  It does not
identify several vertex or edge variables with the same ROABP variable.
That identification is the substantive next step, treated below.
Over a smaller finite field, the argument gives an extension-field
hitting set after constructing a sufficiently large extension; it is
not a statement that the same number of base-field points exists.

\subsection{I-02. The quotient obstruction is already present in a two-path diamond}
\label{I-02}

\paragraph{Statement.}
Let the monomial map $q:F[z_1,\ldots,z_s]\longrightarrow
F[x_1,\ldots,x_n]$ send $z_j$ to $\mathbf x^{a_j}$, and let $A$ be
the matrix with columns $a_j$.  A diagonal one-parameter action
$z_j\mapsto t^{r_j}z_j$ is induced by a diagonal action
$x_i\mapsto t^{\omega_i}x_i$ precisely when
$r=A^{\mathsf T}\omega$.  If the $a_j$ generate a proper sublattice,
requiring $r$ to annihilate $\ker(A:\mathbb Z^s\to\mathbb Z^n)$
only guarantees a character on the image lattice; an integral
extension to all $\mathbb Z^n$ is an additional lattice condition.
Consequently path isolation does not, in general, descend through $q$.

\paragraph{Proof and counterexample.}
Equivariance on each generator reads
$t^{r_j}\mathbf x^{a_j}=t^{\langle\omega,a_j\rangle}\mathbf x^{a_j}$,
which is equivalent to the displayed matrix equation.  The kernel
description follows from the factorization property of homomorphisms
of free abelian groups, with the stated extension qualification.

For the counterexample take the simple diamond
$s\to u\to t$ and $s\to v\to t$.  Label the two paths so that their
vertex-variable polynomial is $z_u-z_v$.  The map $q(z_u)=q(z_v)=x$
gives the zero polynomial.  Every min-isolating weight on this diamond
has $r_u\ne r_v$, so that its weighted polynomial
$t^{r_u}-t^{r_v}$ is nonzero.  Such a weight does not descend, because
$z_u-z_v\in\ker q$ is not homogeneous for it.  This remains a valid
counterexample in characteristic two, where the two distinct powers
still have nonzero coefficients.  Alternatively give both paths label
$1$ in characteristic zero: the original polynomial is $2x\ne0$,
but every weight induced from $x\mapsto t^{\omega}$ ties its two paths.
Thus compatible minimum-path isolation is not even necessary for
nonzero scalar output, while arbitrary path isolation is not sufficient
to certify nonzero output after the quotient. \hfill$\square$

\paragraph{What is refuted.}
The refuted assertion is the direct identification of a Hoza seed with
an ROABP variable-weight assignment.  Neither polynomial ROABP hitting
sets nor their possible derivation by a different construction is
refuted by this example.

\subsection{I-03. Three exact characterizations of basis isolation}
\label{I-03}

\paragraph{Statement.}
Let $S\subset\mathbb Z_{\geq0}^n$ index distinct monomials and let
$V=\operatorname{span}_F\{v_a:a\in S\}$ have dimension $r$ over an arbitrary field.
For $\omega\in\mathbb Z^n$ set $w_a=\langle\omega,a\rangle$.
For integer weights we use the same order-theoretic definition as
I-B02, allowing negative variable weights; the hitting-set
applications below use nonnegative weights.
The following are equivalent.
\begin{enumerate}
\item $\omega$ is a basis-isolating weight assignment in the sense of
this extended definition.
\item Every nonzero linear functional $\lambda\in V^*$ makes
$\sum_a\lambda(v_a)\mathbf x^a$ have a unique monomial of minimum weight.
\item At each weight level $c$, at most one coefficient $v_a$ with
$w_a=c$ survives modulo
$V_{<c}=\operatorname{span}\{v_b:w_b<c\}$.
\end{enumerate}
For a fixed coefficient basis $B$, write uniquely
$v_a=\sum_{b\in B}\lambda_{ab}v_b$.  It is the basis isolated by
$\omega$ exactly when
\[
 \langle\omega,a-b\rangle>0
 \quad(a\notin B,\ \lambda_{ab}\ne0),
 \qquad
 \langle\omega,b-b'\rangle\ne0\quad(b\ne b'\in B).
\tag{I.3}\label{eq:I-cone}
\]
The first inequalities alone characterize $B$ as the unique
minimum-total-weight coefficient basis.

\paragraph{Proof.}
Assume (1), with isolated basis $B$.  For a nonzero $\lambda$, take
the least-weight basis member $b$ on which $\lambda$ is nonzero.
It is unique.  Every nonbasis coefficient of no greater weight lies
in the span of smaller basis coefficients and is annihilated by
$\lambda$.  This proves (2), including its stronger unique-monomial
assertion.

If (3) fails, choose two nonzero coefficient classes at a level $c$.
There is a functional on $V/V_{<c}$ nonzero on both: for independent
classes assign value $1$ to both and extend a basis; for dependent
nonzero classes use any functional nonzero on one of them.  Pull it
back to $V$.  It vanishes below $c$ and is nonzero on two monomials
at $c$, contradicting (2).  This argument works also over finite
fields.  Thus (2) implies (3).

Under (3), choose the unique surviving monomial at every level where
the coefficient span grows.  These coefficients form a basis, one
per growth level.  All remaining coefficients at a level already
belong to the span of smaller levels.  Iterating through the finitely
many levels proves (1).  Unique coordinates in a fixed basis prove
the equivalence with (I.3).

If $B$ is the unique minimum-total-weight basis and
$\lambda_{ab}\ne0$, then $B-b+a$ is also a basis.  Its total weight
is strictly larger, so $w_a>w_b$.  Conversely, suppose these strict
inequalities hold.  Expand any other basis $C$ in the basis $B$.
Its change-of-basis determinant is nonzero, so at least one
permutation in its determinant expansion is nonzero.  Every common
member of $B\cap C$ is forced to match itself; every member
$a\in C\setminus B$ matches some $b\in B\setminus C$ with
$\lambda_{ab}\ne0$.  Summing the strict inequalities over this
matching gives $\sum_{a\in C}w_a>\sum_{b\in B}w_b$.
Thus $B$ is the unique minimum-total-weight basis. \hfill$\square$

\paragraph{Geometric consequence.}
For a fixed ordering of the weights within $B$, (I.3) is an open
rational polyhedral cone in the cocharacter space $\mathbb R^n$.
The union of these cones is precisely the set of BIWAs.  Equivalently,
the one-parameter subgroup simultaneously gives a monomial initial
form for every nonzero scalar projection of $D$.  Isolating one path
or one preselected scalar output has weaker quantifiers.

\subsection{I-04. Projected matroid vertices and polynomial random weights}
\label{I-04}

\paragraph{Statement.}
Let $\mathcal B$ be the coefficient bases of the data in I-03, and put
\[
 \beta(B)=\sum_{a\in B}a,\qquad
 Q=\operatorname{conv}\{\beta(B):B\in\mathcal B\}\subset\mathbb R^n.
\]
Every vertex of $Q$ is the image of exactly one indexed basis.
Assume $S\subseteq\{0,\ldots,\delta\}^{n}$, $r\ge1$, $\delta\ge1$.
Choose independently
$a_i\in\{1,\ldots,M_1\}$ and $b_i\in\{1,\ldots,M_2\}$.
With $L=n\delta M_2+1$, the weights
\[
 \omega_i=L a_i+b_i
\tag{I.4}
\]
are basis isolating with probability at least
\[
 1-\frac{nr\delta}{M_1}-\frac{\binom r2}{M_2}.
\tag{I.5}
\]
For example $M_1=4nr\delta$ and $M_2=4\max(1,r^2)$ give a positive
constant success probability and polynomially bounded weights.  This
is a randomized existence statement, not an explicit small family.

\paragraph{Proof of the vertex assertion.}
Suppose two distinct bases $B,C$ have the same image $\beta$.
Take $e\in B\setminus C$.  Let $\ell$ be the coordinate functional
of $v_e$ in the basis $B$, and expand $v_e=\sum_{f\in C}c_fv_f$.
Because $1=\ell(v_e)=\sum c_f\ell(v_f)$, some
$f\in C\setminus B$ has both $c_f\ne0$ and $\ell(v_f)\ne0$.
Hence $B-e+f$ and $C-f+e$ are both bases.  Their images are
$\beta+(f-e)$ and $\beta-(f-e)$, which are distinct because the
ground-set exponents are distinct.  Their midpoint is $\beta$,
so $\beta$ is not a vertex.  Every vertex has at least one preimage
from the finite defining set, proving the claim.

\paragraph{Proof of the probability bound.}
First consider any finite set $T\subseteq\{0,\ldots,D\}^{n}$ and
uniform independent weights $a_i\in[M_1]$.  Fix all weights except
$a_i$.  The least weight as a function of a real parameter $t=a_i$
is the lower envelope of at most $D+1$ affine functions, one for
each possible $i$-th coordinate.  Its slopes are among
$0,\ldots,D$, and there are at most $D$ breakpoints at which
different slopes minimize.  If two different points of $T$ minimize,
they differ in some coordinate $i$, and that coordinate weight is
one of these breakpoints.  Conditioning on the other coordinates
and applying a union bound gives probability at most $nD/M_1$.
Apply this with $T=\{\beta(B)\}$ and $D=r\delta$.
With probability at least $1-nr\delta/M_1$, there is a unique
minimizing point of $Q$.  It is a vertex and, by the preceding
argument, has a unique basis preimage $B$.

Conditional on the first weights and this $B$, two different
monomials in $B$ have equal $b$-weight with probability at most
$1/M_2$: condition on every $b_i$ but one whose exponent difference
is nonzero.  There is at most one possible value of that remaining
integer weight.  A union bound over pairs costs
$\binom r2/M_2$.  By I-03, the unique first-weight minimum basis
satisfies $a\cdot m>a\cdot b$ whenever the $b$-coordinate of
$v_m$ is nonzero and $m\notin B$.  Each difference is an integer
at least one.  The perturbation changes it by at most
$n\delta M_2$, so multiplication by $L=n\delta M_2+1$ preserves
strict positivity.  It also preserves inequalities between basis
members with unequal first weights, and separates any remaining
ties by the second weights.  Equation (I.3) proves (I.5).
All randomness and weights in this proof are integers; no assumption
on the characteristic of the coefficient field enters the argument.
\hfill$\square$

\paragraph{The remaining derandomization problem.}
The desired generator must hit these cones for coefficient matroids
arising from products of univariate matrices, with a family and
maximum weights both polynomial in $(n,w,\delta)$.  Hoza's theorem
hits path-polytope cones in independent vertex coordinates.  I-02
shows why those cocharacter spaces cannot be silently identified.
I-04 narrows the required object to normal cones of a specified
projected representable-matroid polytope, with the finite
within-basis tie hyperplanes removed.

\subsection{I-05. Positive Gram determinants certify bases, not compressed rank}
\label{I-05}

\paragraph{Statement.}
For complex coefficient vectors, choose coordinates $V=\mathbb C^r$
and form
\[
 G_D(\mathbf z)=\det\left(\sum_{a\in S}v_av_a^*\mathbf z^a\right).
\]
Then
\[
 G_D(\mathbf z)=\sum_{B\in\mathcal B}
                  |\det C_B|^2\mathbf z^{\beta(B)},
\qquad \operatorname{Newt}(G_D)=Q,
\tag{I.6}
\]
where $C$ has columns $v_a$.  At a vertex of $Q$, the coefficient
is the squared volume of its unique basis from I-04.  Nevertheless,
positivity of $G_D(t^\omega)$ does not imply that the coefficient
span of $D(t^\omega)$ retains its dimension.

\paragraph{Proof.}
Apply the Cauchy--Binet identity to
$C\operatorname{diag}(\mathbf z^a)C^*$.
Every nonzero summand is strictly positive in its coefficient, so
there are no cancellations between bases with the same exponent
sum; their convex hull is exactly the Newton polytope.  I-04
identifies its vertex coefficients.  For the negative assertion,
take $D(x,y)=e_1x+e_2y$ and $\omega=(1,1)$.  Its Gram determinant
is $xy$ and specializes to $t^2$, whereas
$D(t,t)=t(e_1+e_2)$ has coefficient span of dimension one instead
of two.  The unique coefficient basis has two equal member weights;
this is exactly the extra condition in I-03. \hfill$\square$

\paragraph{Field and geometric qualification.}
The conjugate transpose and positivity in (I.6) are statements on
a complex vector space with a Hermitian structure.  They are not
positivity statements in an arbitrary characteristic.  The
determinantal/Cauchy--Binet identity itself has an integral
two-matrix form with independent dual variables.  Specializing the
second matrix to a conjugate transpose is a real, not a complex
algebraic, operation.  This qualification is essential in translating
the formula into a scheme over $\mathbb Z$.

\subsection{I-06. Exact black-box implementation of nonlinear exponent weights can be exponential}
\label{I-06}

\paragraph{Statement.}
Let $K$ be any field, let $t$ be transcendental over it, and define
a linear functional on polynomials of individual degree at most $2d$
by
\[
 \mathcal L(\mathbf x^\alpha)=\prod_{i=1}^n h_i(\alpha_i),
 \qquad h_i(e)\in K(t).
\]
Put $H_i=(h_i(a+b))_{0\le a,b\le d}$.
If one exact universal implementation is
\[
 \mathcal L(f)=\sum_{j=1}^M c_j f(\xi_j),
\tag{I.7}
\]
where the points and coefficients are fixed independently of $f$
over any extension of $K(t)$, then
\[
 M\ \ge\ \prod_{i=1}^n\operatorname{rank}(H_i).
\tag{I.8}
\]
In particular $\mathcal L(\mathbf x^\alpha)=
t^{\sum_i\alpha_i^2}$ requires $M\ge(d+1)^n$.
For degree two, the nonlinear exponent filters supplied by every
Hoza seed that separates three exponent symbols also have the
lower bound $M\ge2^n$.

\paragraph{Proof.}
Index a moment matrix by
$\alpha,\beta\in\{0,\ldots,d\}^n$ and set
$H_{\alpha\beta}=\mathcal L(\mathbf x^{\alpha+\beta})$.
It is the tensor product $\bigotimes_iH_i$, with rank the product
of the ranks.  Each evaluation in (I.7) contributes the rank-one
matrix $c_j(\xi_j^\alpha)_\alpha(\xi_j^\beta)_\beta^{\mathsf T}$,
so its rank is at most $M$.  Extension of scalars does not change
the rank, proving (I.8).

For $h_i(e)=t^{e^2}$, factor the $(a,b)$ entry as
$t^{a^2}t^{b^2}(t^{2a})^b$.  The remaining matrix is Vandermonde
on the distinct elements $1,t^2,\ldots,t^{2d}$ and has full rank,
in every characteristic.

For the Hoza assertion take three distinct binary codewords
$C(0),C(1),C(2)$ and put
$h_x(e)=\sum_j6^{\ell-j}C(e)_j x_j$ as integers.
If $h_x(0)+h_x(2)=2h_x(1)$, then
\[
 \sum_j6^{\ell-j}x_j\bigl(C(0)_j+C(2)_j-2C(1)_j\bigr)=0.
\]
The digits lie between $-2$ and $2$.  The first nonzero digit
dominates the absolute sum of all later digits, since
$2\sum_{k=0}^{q-1}6^k<6^q$.  Thus every digit vanishes, which for
binary digits says the three selected code bits agree in every
selected coordinate.  Hence all three values $h_x(e)$ coincide.
Conversely constant values give that equality.  Any seed isolating
the three-arm graph is nonconstant on its three internal vertices,
so the equality fails.  At layer $i$, the scalar multiplier
$6^{\nu_2(i)}$ in (I.1) preserves this failure.  Consequently the
$2\times2$ Hankel determinant of the exponent filter
\[
 e\longmapsto t^{6^{\nu_2(i)}h_x(e)}
\]
is the difference of two distinct monomials, and has nonzero
determinant over $K(t)$.  Equation (I.8) now gives $2^n$.
\hfill$\square$

\paragraph{Scope of the lower bound.}
This rules out the proposed exact linear interpolation implementation
of independently reweighted monomial letters by a polynomial number
of evaluations of the original black box.  It is not a lower bound
on general ROABP PIT, adaptive reconstruction, or approximate
integration.  Even requiring (I.7) only for width-one ROABPs implies
the same identities on every monomial and hence the same lower
bound by linearity; width-one PIT itself remains elementary.

\subsection{I-07. Haar averaging repairs invariance by discarding state isolation}
\label{I-07}

\paragraph{Statement.}
For a fixed complex $m\times m$ matrix $D$, the function
$A\mapsto\operatorname{tr}(DA)$ is invariant under simultaneous
conjugation if and only if $D$ is scalar.  Moreover
\[
 \int_{U(m)}UDU^*\,dU=\frac{\operatorname{tr}D}{m}I.
\tag{I.9}
\]
Thus independently averaging fixed diagonal state-weight insertions
in a trace word replaces each insertion by a scalar and recovers
a scalar multiple of the original trace word.  If the same unitary
is used for two insertions, the exact formula for $m\ge2$ is
\[
 \int\operatorname{tr}(A UDU^* B UEU^*)\,dU
 =\alpha\operatorname{tr}(AB)+\beta\operatorname{tr}A\operatorname{tr}B,
\tag{I.10}
\]
where
\[
 \alpha=\frac{m\operatorname{tr}D\operatorname{tr}E-
                    \operatorname{tr}(DE)}{m(m^2-1)},\qquad
 \beta=\frac{m\operatorname{tr}(DE)-
                    \operatorname{tr}D\operatorname{tr}E}{m(m^2-1)}.
\]
Even positive semidefinite $D,E$ need not make these coefficients
nonnegative: for orthogonal rank-one projections,
$\beta=-1/[m(m^2-1)]$.

\paragraph{Proof.}
Conjugation invariance for every $A$ gives
$\operatorname{tr}((g^{-1}Dg-D)A)=0$ for all $A$.  The trace pairing
is nondegenerate, so $D$ commutes with every invertible matrix and
is scalar.  The reverse direction is immediate.  Haar invariance
makes the left side of (I.9) commute with every unitary; the
commutant is scalar, and its trace determines the scalar.  Repeated
integration gives the independent-insertion claim.

For completeness, in (I.10) define
$T=\int(UDU^*)\otimes(UEU^*)\,dU$.  The commutant of
$U\otimes U$ is spanned by the identity $I$ and tensor flip $S$.
Here is an elementary verification: commutation with diagonal unitaries
forces matrix entries to preserve the unordered pair of tensor indices;
permutation unitaries make the two-dimensional off-diagonal blocks
$\left(\begin{smallmatrix}b&c\\c&b\end{smallmatrix}\right)$ and
the diagonal blocks a common scalar $a$.  A rotation mixing two basis
vectors forces $a=b+c$.  The resulting operator is $bI+cS$.
Thus $T=\alpha I+\beta S$.  Its trace and its trace after multiplication
by $S$ are respectively
\[
 m^2\alpha+m\beta=\operatorname{tr}D\operatorname{tr}E,
 \qquad m\alpha+m^2\beta=\operatorname{tr}(DE).
\]
Solving gives the stated values.  Finally
$\operatorname{tr}(A X B Y)=
\operatorname{tr}((A\otimes B)(X\otimes Y)S)$;
contracting $T$ gives (I.10).  For orthogonal rank-one projections,
$\operatorname{tr}D=\operatorname{tr}E=1$ and
$\operatorname{tr}(DE)=0$, proving the sign assertion.
\hfill$\square$

\paragraph{Meaning for the proposed transfer.}
Hoza's internal-state weights enter matrix products as fixed
diagonal insertions.  Their lack of conjugation invariance cannot
be repaired by independent Haar averaging while preserving the
state information.  Correlated averaging changes the trace word
into a combination of different trace invariants and can reintroduce
negative coefficients.  Neither averaging is the missing
basis-isolation construction.

\subsection{I-08. The exact matrix-invariant consequence, conditional on the remaining gate}
\label{I-08}

\paragraph{Statement.}
Suppose one has an explicit hitting set $\mathcal H\subset F^{m^2}$
of polynomial size, constructible in polynomial arithmetic time,
for ROABPs of depth $m^2$, width at most $2m^2$, and individual
degree less than $r$.  Let $F$ have characteristic zero and put
\[
 M(x)=\sum_{j=0}^{r-1}M_jx^j,\qquad
 T_{\ell,\alpha}(\mathbf M)=
 \operatorname{tr}\bigl(M(\alpha_1)\cdots M(\alpha_\ell)\bigr),
 \quad1\le\ell\le m^2,quad\alpha\in\mathcal H.
\tag{I.11}
\]
These at most $m^2|\mathcal H|$ invariants are an explicit
polynomial-size separating family for simultaneous conjugation.
Over algebraically closed $F$, the invariant ring is finite over
the subring they generate.  This is precisely the specialization
of Forbes and Shpilka (2013), Construction~3.7, Theorem~3.8 and
Corollary~3.10, to a polynomial hitting set.  It supplies the
matrix-invariant instance of explicit Noether normalization in
the sense of that paper.  It does not by itself extract an
algebraically independent system of parameters of cardinality
equal to the quotient dimension.

\paragraph{Derivation and exact source dependence.}
The coefficient of $x_1^{a_1}\cdots x_\ell^{a_\ell}$ in
$\operatorname{tr}(M(x_1)\cdots M(x_\ell))$ is the trace word
$\operatorname{tr}(M_{a_1}\cdots M_{a_\ell})$.  The latter words,
for $\ell\le m^2$, generate the characteristic-zero invariant ring
as recorded in Forbes--Shpilka Theorem~1.3 and Corollary~3.6.
For two tuples $A,B$, the difference of their trace-generating
polynomials is an ROABP of width at most $2m^2$, by their Lemma~3.2.
The hitting-set assumption therefore detects exactly whether
some trace generator separates $A$ and $B$.  Their Theorem~3.9
(the stated specialization of Derksen--Kemper Theorem~2.3.12)
then gives integrality, and finite generation upgrades this to
a finite morphism.  The only new hypothesis in this specialization
is the polynomial hitting-set gate; it has not been proved here.
\hfill$\square$

\paragraph{Consequences that do not follow.}
This does not prove general arithmetic PIT, arbitrary explicit
Noether normalization, permanent circuit lower bounds, or
$P\ne NP$.  Forbes--Shpilka already prove deterministic NC
orbit-closure intersection in Theorem~1.15 using white-box PIT.
Their discussion after Theorems~1.13--1.15 explicitly separates
the matrix-invariant case from the stronger general explicit-variety
problem.  In the supplied revised GCT~V source, the general-PIT
equivalence is Theorem~5.14, and the matrix-specific ROABP replacement
is \S7.5, Lemma~7.9.  The earlier GCT~V theorem numbering cited by
Forbes--Shpilka must not be transferred to the revised source.  No new polynomial hitting set
may be inserted into (I.11) until its own gate is closed.

\subsection{I-09. An exact positive transfer-matrix certificate remains white box}
\label{I-09}

\paragraph{Statement.}
For a complex ROABP in known matrix-product form
\[
 f(\mathbf x)=u^{\mathsf T}\prod_{i=1}^n
                  \left(\sum_{a=0}^{\delta}A_{i,a}x_i^a\right)v
             =\sum_{\alpha}c_{\alpha}\mathbf x^{\alpha},
\]
its coefficient squared-norm generating polynomial is
\[
 \sum_{\alpha}|c_\alpha|^2\mathbf z^\alpha
 = (u\otimes\bar u)^{\mathsf T}
     \prod_{i=1}^n\left(\sum_{a=0}^{\delta}
                    A_{i,a}\otimes\overline{A_{i,a}}\,z_i^a\right)
                     (v\otimes\bar v).
\tag{I.12}
\]
In particular $f=0$ if and only if the right side at
$\mathbf z=\mathbf1$ is zero.  Its matrices have width $w^2$.
The equality also gives the usual torus identity
\[
 \int_{[0,2\pi]^n}|f(e^{i\theta_1},\ldots,e^{i\theta_n})|^2
                \prod_i\frac{d\theta_i}{2\pi}
       =\sum_\alpha|c_\alpha|^2.
\tag{I.13}
\]

\paragraph{Proof.}
The coefficient $c_\alpha$ is
$u^{\mathsf T}A_{1,\alpha_1}\cdots A_{n,\alpha_n}v$.
Multiplicativity of tensor products identifies its squared
absolute value with the product of the corresponding tensor
matrices and boundary vectors.  Summing independently over each
$\alpha_i$ gives (I.12).  The sum is a finite sum of nonnegative
real numbers, and it is zero exactly when every coefficient is
zero.  In (I.13), expansion of the square and integration of
$e^{i\langle\alpha-\beta,\theta\rangle}$ kills exactly the terms
$\alpha\ne\beta$. \hfill$\square$

\paragraph{Why this does not close the black-box gate.}
The transfer matrices use the individual coefficient matrices
$A_{i,a}$, which a black box does not supply.  The torus integral
is another exact expression, but the product quadrature that is
exact for all such degree bounds uses $(\delta+1)^n$ nodes.
No lower bound on all specialized ROABP quadratures is asserted.
I-06 proves a lower bound for the different, explicitly specified
nonlinear exponent filter.  Positivity in (I.12) closes the
white-box zero-testing question only; it does not supply the
missing uniform polynomial set of original evaluations.

\subsection{I-10. Schemes, cocharacters, flags, and flat limits}
\label{I-10}

\paragraph{Statement and dictionary.}
Let $R$ be a domain carrying the relevant coefficients, and let
$T=\mathbb G_{m,R}^{n}$.  The integer weight vector $\omega$
is the cocharacter
$\lambda_\omega:\mathbb G_{m,R}\to T$,
$t\mapsto(t^{\omega_1},\ldots,t^{\omega_n})$.
For nonzero $f=\sum_a c_a\mathbf x^a$ and
$\mu=\min_{c_a\ne0}\langle\omega,a\rangle$, define
\[
 \widetilde f(t,\mathbf x)=
 \sum_a c_a t^{\langle\omega,a\rangle-\mu}\mathbf x^a.
\tag{I.14}
\]
This is a polynomial over $R[t]$, its special fiber is the
weight-initial form, and its generic fiber is the original
polynomial under a torus change of coordinates.  A unique
minimizing exponent gives a monomial special fiber.  Over a field,
the hypersurface family defined by (I.14) is flat over the
$t$-line.  For a homogeneous ideal, the corresponding weight
homogenization is the usual flat Rees/initial-ideal degeneration;
this is stronger than merely selecting a least monomial of one
polynomial and must not be conflated with it.

For a vector polynomial, I-03 identifies a complete coefficient
flag whose successive quotients have one supported weight
monomial.  With an explicit coefficient matrix, nonvanishing
maximal minors define a matroid stratum of the appropriate
Grassmannian.  The exponents of its Pl\"ucker coordinates under
$\lambda_\omega$ are the sums $\beta(B)$.  Its torus orbit has
weight polytope $Q$; I-04 says a vertex with nonzero Pl\"ucker
coordinate corresponds to a unique indexed coefficient basis.
The BIWA cones refine its normal fan by the individual
within-basis weight comparisons in (I.3).

\paragraph{Proof of the flatness assertion and the dictionary.}
Formula (I.14) has no negative power of $t$ and is not divisible
by $t$.  Over a field $F$, for every irreducible polynomial
$p(t)$ the image of $\widetilde f$ modulo $p(t)$ is nonzero:
for $p=t$ use its nonzero initial form; otherwise $t$ is invertible
in $F[t]/(p)$ and its monomial coefficients remain nonzero.
Thus no nonzero $p(t)$ divides $\widetilde f$ in the UFD
$F[t,\mathbf x]$.  If $p(t)g\in(\widetilde f)$, primality of
the irreducible factors of $\widetilde f$ implies
$g\in(\widetilde f)$.  Hence the quotient by
$(\widetilde f)$ is torsion-free over the PID $F[t]$, and
therefore flat.  The generic-fiber statement follows by
inverting $t$ in (I.14).  Exterior products of coefficient
columns are precisely Pl\"ucker coordinates, and their weights
add, giving $\beta(B)$.  Equation (I.3) then identifies the
normal-cone refinement.  This proves each part of the stated
dictionary directly. \hfill$\square$

\paragraph{Integral qualification.}
The torus and monomial maps are defined over $\mathbb Z$.
A specified coefficient matroid and its rank flag persist on an
open localization where the finitely many designated nonzero
minors remain nonzero; they need not persist in every residue
characteristic.  Over a general base $R$, flatness of a
hypersurface quotient also needs the relevant coefficient and
fiber hypotheses.  The field proof above is not a claim of
automatic flatness over arbitrary $R$.  The Grothendieck locator
table in the main paper should be used for the exact SGA/FGA/EGA
placement of these notions, rather than inventing a theorem
number for these elementary derived statements.

\subsection{I-11. Exterior and subspace-state repairs do not yet produce a graph generator}
\label{I-11}

\paragraph{Checked finite assertions.}
For coefficient ambient dimension $k$, the standard coordinate
space $\bigwedge^rF^k$ has dimension $\binom kr$.  Over
$\mathbb F_q$, there are at least $q^{r(k-r)}$ distinct
$r$-dimensional subspaces of $F_q^k$: the graphs of the
$q^{r(k-r)}$ linear maps $F_q^r\to F_q^{k-r}$ are all distinct.
Thus a literal graph with one state for each such subspace is
not polynomial in $k$ when $r$ grows with $k$.

Exterior coordinates also retain cancellation.  For two equal
vectors $v_1=v_2=e_1+e_2$, the determinant is
$1\cdot1-1\cdot1=0$, although its expansion has two nonzero
permutation terms.  Making one term uniquely minimum by
path-dependent weights changes the determinant rather than
certifying its nonvanishing.  Tensoring $r$ independently chosen
word trajectories gives a standard realization with $w^{O(r)}$
states; for fixed $r$ this removes that particular size cost,
but the cancellation and weight-descent problems remain.

\paragraph{Proof and scope.}
The exterior dimension follows from the basis wedges indexed
by $r$-element subsets of a fixed $k$-element basis.  The graph
of a linear map is an $r$-dimensional subspace, and projection
onto the first summand recovers the map, proving distinctness.
The determinant example is immediate.  The tensor product
construction records one of $w$ state positions for each of
the $r$ trajectories, so has $w^r$ positions before additional
boundary coordinates; determinant signs are imposed at the
boundary and do not disappear.  These are size and correctness
facts about the specified constructions, not lower bounds on
all possible reductions.  A scalar output can itself have
coefficient-space dimension one, so constant coefficient rank
alone cannot make the general ROABP hitting-set task disappear.
\hfill$\square$

\subsection{I-12. Resolving coefficient support is an arithmetic counting gate}
\label{I-12}

\paragraph{Background I-B05.}
Allender and Ogihara (1996), \S2, p.~4, define the accepting-minus-rejecting
gap model; \S5, Definition~2, p.~11, defines $C_{=}L$ by the vanishing
of a $\operatorname{GapL}$ function.  Their Theorem~1, p.~2, identifies
$\operatorname{GapL}$ with logarithmic-space reductions to integer
determinant, and Theorem~14, p.~12, gives the corresponding singularity
completeness result.  The following ROABP formulation is proved directly.

\paragraph{Statement.}
Restrict to explicitly given layered ROABPs over $\mathbb Z$
(equivalently over $\mathbb Q$) with coefficients in
$\{0,1,-1\}$ and polynomially bounded width and length.  The input
lists the coefficient entries and the queried exponent vector;
univariate edge labels may be given densely or as canonical sparse
lists with binary exponents.  Deciding
whether a specified monomial coefficient is nonzero is complete,
under deterministic logarithmic-space reductions, for the class
\[
 \operatorname{coC}_{=}L
   =\{\{x:g(x)\ne0\}:g\in\operatorname{GapL}\}.
\]
Here $\operatorname{GapL}$ is defined by the difference of the
numbers of accepting and rejecting paths of polynomial-time
nondeterministic logarithmic-space machines.  Hardness already
holds for ROABPs whose polynomial is one monomial times a scalar.
Therefore a logarithmic-space uniform polynomial-size reduction
which represents exactly nonzero coefficient support by ordinary
directed reachability would imply
$\operatorname{coC}_{=}L\subseteq NL$.

\paragraph{Proof.}
For membership, guess a state path in the layered ROABP.  At layer
$i$ choose the term of the prescribed exponent from its univariate
edge label.  A zero term contributes nothing; otherwise track the
parity of the negative signs.  The difference between positive and
negative paths is exactly the specified monomial coefficient.
State indices, a layer counter, and a sign bit require logarithmic
space.  Neutral or invalid branches can be arranged to contribute
one accepting and one rejecting path, hence zero gap.

For hardness, start with a polynomial-time nondeterministic
logarithmic-space machine computing a gap $g(x)$.  Its time-unrolled
configuration graph has polynomial width and length and is
constructible from $x$ in logarithmic space.  Pad halted computations
by deterministic transitions, retaining their accepting or rejecting
tag, so every path has the same length.  Mark the final edge
of accepting paths by $1$ and of rejecting paths by $-1$;
all preceding edges have label $1$.  Separate parallel transitions
by an intermediate layer if necessary.  Multiply every edge label
in layer $i$ by a fresh variable $x_i$.  Every source-to-target path
then has the same variable monomial $\prod_i x_i$, with its
accepting or rejecting sign, so the output is
$g(x)\prod_i x_i$.  This is a known-order multilinear ROABP with
coefficients $0,1,-1$.  Its unique possible monomial coefficient
is nonzero exactly when $g(x)\ne0$.  Finally, composition with the
postulated graph reduction puts every such language in $NL$.
\hfill$\square$

\paragraph{Relation to the other gates.}
This is a precise complexity formulation of the operation which
would have to cancel equal-monomial paths before graph isolation.
It is not a separation of complexity classes.  Nor does it rule
out a different polynomial-time algebraic construction: the
explicit ROABP supplies its matrices, and white-box polynomial
identity testing is already possible.  In scheme language the
coefficient is a regular function of the edge coefficients;
testing its vanishing uses the kernel of a coefficient map, not
merely whether its parameterizing path set is empty.  The diamond
in I-02 is the smallest instance of this distinction.

\subsection{I-13. Exact universal energy cubature needs the full coefficient dimension}
\label{I-13}

\paragraph{Statement.}
Let $n\ge1$, $\delta\ge0$, and $D=(\delta+1)^n$.  Fix points
$\xi_1,\ldots,\xi_Q\in\mathbb C^n$ and a complex $Q\times Q$ matrix $C$,
independent of the polynomial.  Suppose that every width-one product
\[
 f(\mathbf x)=\prod_{i=1}^n\left(\sum_{a=0}^{\delta}z_{i,a}x_i^a\right),
 \qquad z_{i,a}\in\mathbb C,
\]
satisfies the exact identity
\[
 \sum_{\alpha\in\{0,\ldots,\delta\}^n}|[\mathbf x^\alpha]f|^2
  =\sum_{j,k=1}^Q C_{jk}\,\overline{f(\xi_j)}f(\xi_k).
\tag{I.13}\label{eq:I-cubature}
\]
Then $Q\ge D$.  In particular, signed cubature
$\sum_j c_j|f(\xi_j)|^2$ requires at least $D$ nodes.  More generally,
a formula with $M$ fixed paired terms
$\sum_{j=1}^M c_j\overline{f(\xi_j)}f(\eta_j)$ requires $M\ge D$.
The bound is attained for the diagonal formula by the Cartesian product
of the $(\delta+1)$st roots of unity in each coordinate, all with weight
$1/D$.

\paragraph{Proof.}
Index coordinates by $\alpha\in\{0,\ldots,\delta\}^n$ and put
$a_\alpha=\prod_i z_{i,\alpha_i}$.  Let the evaluation matrix be
$E_{j\alpha}=\xi_j^\alpha$.  The asserted identity is
\[
 a^*(I_D-E^*CE)a=0
\quad\text{for every decomposable coefficient vector }
a=\bigotimes_i(z_{i,0},\ldots,z_{i,\delta}).
\]
Expanding gives a polynomial identity in the variables
$z_{i,a}$ and their conjugates.  Its monomials
\[
 \prod_i\overline{z_{i,\alpha_i}}z_{i,\beta_i},
 \qquad (\alpha,\beta)\in\{0,\ldots,\delta\}^{2n},
\]
are linearly independent as functions of the complex variables.
Indeed, writing each $z=u+iv$ identifies the polynomial ring in
formal $z,\bar z$ with the polynomial ring in $u,v$ by an invertible
complex linear change of variables; a polynomial vanishing on every
real tuple $(u,v)$ has every coefficient zero.  Distinct pairs
$(\alpha,\beta)$ give distinct displayed monomials.  Consequently
$I_D=E^*CE$, so
$D=\operatorname{rank}(I_D)\le\operatorname{rank}(E)\le Q$.
For paired terms the same coefficient comparison represents $I_D$
as a sum of $M$ rank-at-most-one matrices, proving $M\ge D$.

Finally, averaging $\xi^\alpha\overline{\xi^\beta}$ over the stated
root-of-unity grid gives $1$ when $\alpha=\beta$ and $0$ otherwise:
in each coordinate the difference lies between $-\delta$ and
$\delta$, whose only multiple of $\delta+1$ is zero.  Expanding
$|f|^2$ therefore proves the matching formula, for every polynomial
of the stated individual degree, including the width-one class.
\hfill$\square$

\paragraph{Scope and consequence for I-09.}
Torus orthogonality identifies the left side of (I.13) with
$\int_{(S^1)^n}|f(z)|^2\,dz$.  Thus the positive white-box energy
computation in I-09 cannot be replaced by a polynomial number of
fixed black-box samples through any exact quadratic formula of the
displayed kind, even for width one and fixed $\delta\ge1$ as $n$ grows.
This is a lower bound for universal exact cubature, not for hitting
sets: for width one the substitution $x_i=t$ gives the nonzero
univariate product $\prod_i f_i(t)$ of degree at most $n\delta$.
Evaluation at $n\delta+1$ distinct scalars already gives a polynomial
hitting set for that class.  The rank argument makes no claim about
adaptive procedures or other nonlinear uses of evaluations.

\subsection{I-14. Joint compression is a Grassmannian transversality condition}
\label{I-14}

\paragraph{Statement.}
Write $D=\sum_{a\in S}v_a\mathbf x^a$, choose coordinates on
$V=\operatorname{span}\{v_a\}$, and let $C=(v_a)_{a\in S}$ be the
resulting full-row-rank $r\times s$ matrix, where $s=|S|$.
For a finite family $\mathcal W$ of integer weights, form a matrix
$P_{\mathcal W}$ with rows indexed by $S$ and columns indexed by
the pairs $(\omega,c)$ that occur as a monomial weight.  Substitutions
are interpreted as Laurent polynomials if necessary.  Set
\[
 (P_{\mathcal W})_{a,(\omega,c)}
   =\begin{cases}1&\langle\omega,a\rangle=c,\\0&\text{otherwise.}\end{cases}
\]
The following conditions are equivalent:
\begin{enumerate}
\item the coefficients of all the substitutions
$D(t^{\omega_1},\ldots,t^{\omega_n})$, taken jointly, span $V$;
\item $\operatorname{rank}(CP_{\mathcal W})=r$;
\item $\operatorname{im}(C^{\mathsf T})\cap
\ker(P_{\mathcal W}^{\mathsf T})=\{0\}$;
\item for every nonzero $\lambda\in V^*$, at least one substituted
scalar polynomial $\lambda(D(t^\omega))$ is nonzero.
\end{enumerate}
Equivalently, for some $r$-element set $T$ of columns of $P_{\mathcal W}$,
\[
 \sum_{\substack{B\subseteq S\\|B|=r}}
   \det(C_B)\det((P_{\mathcal W})_{B,T})\ne0.
\tag{I.14}\label{eq:I-joint-Plucker}
\]
Thus the failure locus on $\operatorname{Gr}(r,F^s)$ is the
intersection with the linear equations obtained by setting every
expression in (I.14) to zero.  These depend on Pl\"ucker coordinates,
including their values, rather than only on the coefficient matroid.

\paragraph{Proof.}
The $(\omega,c)$ column of $CP_{\mathcal W}$ is exactly
$\sum_{\langle\omega,a\rangle=c}v_a$, the coefficient at that
weight after substitution.  This proves (1)$\Leftrightarrow$(2).
Since $C^{\mathsf T}:F^r\to F^s$ is injective, the transpose of
$CP_{\mathcal W}$ is injective exactly when the intersection in
(3) is zero.  A proper subspace of $V$ has a nonzero annihilating
linear functional, proving (1)$\Leftrightarrow$(4).  Finally,
Cauchy--Binet applied to each $r$-column submatrix of
$CP_{\mathcal W}$ gives (I.14).  These are linear forms in the
Pl\"ucker coordinates $\det(C_B)$, proving the geometric description.

Joint preservation can hold even though no individual member is
basis isolating.  Take
\[
 D=e_1(x_1-x_2)+e_2(x_2-x_3),\qquad
 \mathcal W=\{(0,0,1),(1,0,0)\}.
\]
The two substituted vector polynomials are respectively
$(1-t)e_2$ and $(t-1)e_1$.  Each has coefficient rank one and so
fails I-03, but their coefficients jointly span the original
two-dimensional space.  To see why the matroid alone is insufficient,
over a field of characteristic different from two compare coefficient
rows $(1,1)$ and $(1,-1)$ on the support $\{x_1,x_2\}$ under weight
$(1,1)$.  They have the same rank-one coefficient matroid, while
their compressed rows are respectively $2$ and $0$.
\hfill$\square$

\paragraph{A sufficient black-box gate.}
If $S\subseteq\{0,\ldots,\delta\}^n$, all weights are nonnegative
and at most $B$, and joint preservation holds, evaluating the
$|\mathcal W|$ curves at $n\delta B+1$ distinct field elements
hits every nonzero scalar projection of $D$.  Indeed, condition
(4) supplies a nonzero univariate polynomial of degree at most
$n\delta B$.  This exact sufficient condition is weaker than
requiring one BIWA.  It does not remove cancellation from the
linear Pl\"ucker expressions in (I.14).

\subsection{I-15. Bounded affine decodings of the seed masks have width-two obstructions}
\label{I-15}

\paragraph{Statement.}
Suppose a proposed family of variable weights has the form
\[
 \omega(s)=b+Ts,\qquad
 b\in\mathbb Z^n,\quad T\in\mathbb Z^{n\times k},
\]
where every entry of $b$ and $T$ has absolute value at most $M$.
The seed set can be any subset of $\mathbb Z^k$.  If
\[
 2^n>(2nM+1)^{k+1},
\tag{I.15}\label{eq:I-affine-mask-count}
\]
there are distinct $\alpha,\beta\in\{0,1\}^n$ such that the
nonzero multilinear binomial $\mathbf x^\alpha-\mathbf x^\beta$
vanishes identically on every curve $\mathbf x=t^{\omega(s)}$.
This binomial has a width-two ROABP in every variable order.
In particular the whole family fails even the joint condition I-14.

\paragraph{Proof.}
Map the $2^n$ Boolean exponent vectors $\alpha$ to
\[
 \left(b^{\mathsf T}\alpha,T^{\mathsf T}\alpha\right)\in
 \{-nM,\ldots,nM\}^{k+1}.
\]
By (I.15), two distinct vectors have the same image.  Then
$\langle\omega(s),\alpha-\beta\rangle=0$ for every seed $s$,
so the two substituted powers of $t$ coincide.  This is a Laurent
identity if negative weights are allowed, and a polynomial identity
when the weights are nonnegative.  Distinct monomials have nonzero
coefficients in their difference over any field, including
characteristic two.  The diagonal matrix product with factors
$\operatorname{diag}(x_i^{\alpha_i},x_i^{\beta_i})$ and boundaries
$(1,-1)$ and $(1,1)^{\mathsf T}$ computes the binomial.  Reordering
the commuting diagonal factors gives every variable order.
\hfill$\square$

\paragraph{A necessary dimension for any bounded weight family.}
More generally, let $\mathcal W\subseteq\mathbb Z^n$ have all
coordinates of absolute value at most $B\ge1$, and let
$k=\dim_{\mathbb Q}\operatorname{span}_{\mathbb Q}\mathcal W$.
If $\mathcal W$ hits every multilinear binomial by monomial curves,
then necessarily
\[
 k\ge\frac{n}{\log_2(2nB+1)}.
\]
Indeed, select $k$ members of $\mathcal W$ forming a rational basis
and apply the same pigeonhole argument to their integer dot products
with Boolean exponent vectors.  Equality on this basis implies
equality on every member of $\mathcal W$, independently of how
the seed was decoded.  Thus polynomially bounded output weights
computed as polynomials of any fixed degree $h$ in
$\ell=O(\log n)$ Boolean seed bits also fail for sufficiently large
$n$: their span has dimension at most
$\sum_{j=0}^h\binom{\ell}{j}=O((\log n)^h)$, below the necessary
bound.  The dimension estimate follows by grouping the coefficient
vectors of the multilinear seed monomials; on Boolean seeds powers
of a bit reduce to that bit.

\paragraph{Application to the actual mask formula.}
Every weight in Hoza's formula (I.1) is an integer linear form in
the common $\ell$ seed bits.  Hence assigning those forms to
variables, or taking any fixed linear averaging or decoding with
polynomially bounded integer coefficients, gives $\omega(s)=b+Ts$
with $k=\ell$ and polynomially bounded entries.  A rational average
with a common polynomially bounded denominator has the same property
after clearing that denominator.  Since $\ell=O(\log(WN))$, if
$W,N$ and the resulting entry bound are polynomial in the number
$n$ of variables, condition (I.15) holds for all sufficiently large
$n$.  This rules out this specified seed-to-variable decoding,
even if all its members are used jointly.

For the most direct alphabet decoding the counterexample is explicit
without counting.  Let
$H_s(a)=\sum_j6^{\ell-j}C(a)_js_j$.  Matching letters $0,1$ by
an affine exponent map at layer $i$ gives
\[
 \omega_i(s)=6^{\nu_2(i)}\bigl(H_s(1)-H_s(0)\bigr).
\]
Layers $1$ and $3$ receive the same weight, so $x_1-x_3$ vanishes
for every seed.  This uses the actual layer dependence in (I.1)
and holds for every choice of its code.  There is no contradiction
with minimum-path isolation: the graph weights themselves need not
descend to these variable weights.  The argument does not rule out
nonlinear decoders whose outputs span sufficiently many independent
directions, or unbounded integer weights.

\subsection{I-16. The finite-field row-state lift and its uniform evaluation limit}
\label{I-16}

\paragraph{Statement.}
Over $F=\mathbb F_q$, an explicitly given width-$w$ ROABP
\[
 f=u^{\mathsf T}\prod_{i=1}^n
       \left(\sum_{a=0}^{\delta}A_{i,a}x_i^a\right)v
\]
has an exact deterministic coefficient-support automaton with
at most $q^w$ states per layer.  A state is a row vector
$r\in F^{1\times w}$; its letter-$a$ transition at layer $i$ is
$r\mapsto rA_{i,a}$.  The initial state is $u^{\mathsf T}$, and
a final state is accepting exactly when $rv\ne0$.
Nevertheless, if $0\le\delta<q$, every base-field hitting set
$H\subseteq F^n$ for all width-one individual-degree-$\delta$
polynomials satisfies
\[
 |H|\ \ge\ \left(\frac{q}{q-\delta}\right)^n.
\tag{I.16}\label{eq:I-small-field-lower}
\]
For $\delta=q-1$ all $q^n$ points are necessary and sufficient.
Thus polynomial state-space size alone cannot produce a polynomial
uniform base-field hitting set, even when $w=1$ and $q$ is fixed.

\paragraph{Proof.}
On the word $(a_1,\ldots,a_n)$, the automaton's final scalar is
$u^{\mathsf T}A_{1,a_1}\cdots A_{n,a_n}v$, precisely the monomial
coefficient.  Matrix multiplication has already performed all path
cancellation before the state is tested.  To obtain a simple layered
graph, subdivide each transition with an intermediate vertex labelled
by its current row and its letter; append a common sink after all
accepting states.  The width is at most $(\delta+1)q^w$ and the
length is $2n+1$.  It is constructible from the coefficient matrices
in time polynomial in this displayed size and the field-operation
cost.  Hoza applies to this coefficient-dependent graph, with the
corresponding polynomial-in-$q^w$ seed family; its path weights
still depend on the row states.

For the lower bound, choose independently a $\delta$-element subset
$R_i\subseteq F$ for each coordinate and consider
\[
 f_{R_1,\ldots,R_n}(\mathbf x)
     =\prod_{i=1}^n\prod_{a\in R_i}(x_i-a).
\]
These are $\binom{q}{\delta}^{n}$ distinct nonzero width-one polynomials.
A fixed point $\xi\in F^n$ hits exactly
$\binom{q-1}{\delta}^n$ of them, since each $R_i$ must avoid
$\xi_i$.  Covering all these polynomials therefore requires
\[
 |H|\binom{q-1}{\delta}^n\ge\binom{q}{\delta}^{n},
\]
which is (I.16).  The full grid $F^n$ hits every polynomial of
individual degree less than $q$: induction on the number of variables,
using the univariate root bound, shows that a polynomial vanishing
on that grid must be zero.  This also gives the matching upper bound
when $\delta=q-1$.  For $\delta\ge q$, even that grid cannot be a
hitting set for all formally nonzero polynomials, because
$x_1^q-x_1$ vanishes on it. \hfill$\square$

\paragraph{Precise scope.}
The lower bound concerns evaluations in the coefficient field.
It does not apply to extension-field evaluations, which can behave
very differently for a fixed base-field coefficient class.
The row-state lift is useful as an exact finite arithmetic simulation
and avoids the erroneous unweighted support graph of I-02.
It does not construct the graph from an evaluation oracle or make
its state-dependent path weights into variable cocharacters.

\subsection{I-17. A concrete uniform hitting set for invertible factors}
\label{I-17}

\paragraph{Background and statement.}
Invertible-factor ROABP hitting sets are already a result of
Agrawal et al. (2014), Theorem~3, pp.~5 and~19, and their
Lemma~27, p.~22, gives matrix coefficient concentration from
invertible constant terms.  The following is a self-contained
row-space specialization with constant boundary vectors; no
historical-priority claim is made.

Let $n,w\ge1$ and $\delta\ge1$.  Suppose
\[
 f=u^{\mathsf T}A_1(x_{\pi(1)})\cdots A_n(x_{\pi(n)})v
\]
has square $w\times w$ factors of individual degree at most
$\delta$, each with nonzero determinant polynomial.
Choose sets $S,B\subseteq F$ with
$|S|=nw\delta+1$, $|B|=\delta+1$, and $0\in B$.
Then the following single set, independent of the factors and order,
hits every nonzero such $f$:
\[
 \mathcal H=\{a\mathbf1+b: a\in S,\ b\in B^n,
                         |\operatorname{supp}(b)|\le w-1\}.
\tag{I.17}\label{eq:I-invertible-hit}
\]
Its cardinality is at most
\[
 (nw\delta+1)\sum_{j=0}^{\min(w-1,n)}\binom nj\delta^j.
\]
It is polynomial in $n,\delta$ for every fixed $w$, works in arbitrary
characteristic with the stated field-size condition, and requires
no knowledge of $\pi$.

\paragraph{Proof.}
The nonzero polynomial $\prod_i\det A_i(t)$ has degree at most
$nw\delta$, so some $a\in S$ makes every $A_i(a)$ invertible.
Reindex the variables in their true order for this proof only.
Set $P_0=I$, $P_i=A_1(a)\cdots A_i(a)$, and
\[
 M_i(y)=P_{i-1}A_i(a+y)P_i^{-1}.
\]
Then $M_i(0)=I$ and telescoping gives
\[
 f(a\mathbf1+\mathbf y)
    =u^{\mathsf T}M_1(y_1)\cdots M_n(y_n)P_nv.
\]
For nonzero $f$ we have $u\ne0$.  Let $R_i\subseteq F^{1\times w}$
be the span of the coefficient rows of
$u^{\mathsf T}M_1(y_1)\cdots M_i(y_i)$, and put $d_i=\dim R_i$.
Because the next constant term is $I$, these spaces are nested and
$d_0=1$, $d_i\le w$.

We claim that $R_i$ has a basis consisting of actual coefficient
rows whose monomials have support at most $d_i-1$.  This is true for
$i=0$.  Inductively retain a basis for $R_{i-1}$.  If the dimension
does not grow, that basis suffices.  Otherwise extend it using rows
$rM_{i,a}$, where $r$ is a retained coefficient row and $a>0$.
These rows span all new coefficients, and each is itself a coefficient
row with support at most $d_{i-1}$.  Since $d_i>d_{i-1}$, every
retained or added basis row has support at most $d_i-1$.
This proves the claim.

At $i=n$, at least one coefficient row has nonzero product with
$P_nv$.  Therefore some basis coefficient of support at most $w-1$
also has nonzero product.  Restricting the shifted scalar polynomial
to those variables, with the others zero, gives a nonzero polynomial
of individual degree at most $\delta$.  The grid $B$ on those
variables hits it by repeated univariate root counting.  This point
lies in (I.17).  The set (I.17) is invariant under coordinate
permutations, so its construction never needed the true order.
Counting supports proves its size bound. \hfill$\square$

The support bound is sharp for the normalized row-space argument:
for $w\ge2$ and $n=w-1$, factors $I+x_iE_{i,i+1}$ and boundaries $e_1,e_w$
compute $x_1\cdots x_{w-1}$.  For $w=1$, the theorem is the familiar
diagonal univariate test.  Generically singular factors require a
different argument, as the following explicit family shows.

\subsection{I-18. Regularization must preserve a unit Pl\"ucker coordinate}
\label{I-18}

\paragraph{Statement and counterexample.}
Let $m\ge2$, $n=2m$, and
\[
 f(\mathbf x)=\prod_{j=1}^m(x_{2j}-x_{2j-1}).
\]
This is a nonzero multilinear width-two ROABP in the displayed
variable order.  It vanishes at every point $a\mathbf1+b$ with
$|\operatorname{supp}(b)|<m$, for every scalar $a$.
Consequently the sampling restriction in I-17 cannot extend to
all width-two ROABPs by regularizing their matrices and specializing
the regularization parameter to zero.

More precisely, let $R=F[[\varepsilon]]$, $K=F((\varepsilon))$,
$C(\varepsilon)\in R^{r\times s}$ have
$\operatorname{rank}C(0)=r$, and let $P\in R^{s\times h}$ be any
integral compression.  Then
\[
 \operatorname{rank}_K(CP)=r
 \quad\Longleftrightarrow\quad
 \text{some $r$-minor of $CP$ has finite valuation},
\]
whereas
\[
 \operatorname{rank}_F(C(0)P(0))=r
 \quad\Longleftrightarrow\quad
 \text{some $r$-minor of $CP$ has valuation zero}.
\tag{I.18}\label{eq:I-unit-Plucker}
\]
The latter is equivalent to $\operatorname{im}_R(CP)=R^r$.
Generic span equality alone is insufficient.

\paragraph{Proof.}
Use the two singular matrix factors
\[
 L(x)=\begin{pmatrix}-x&1\\0&0\end{pmatrix},\qquad
 R(y)=\begin{pmatrix}1&0\\y&0\end{pmatrix}.
\]
Their product has $(1,1)$ entry $y-x$ and all other entries zero.
Taking the ordered product of these pairs with boundaries $e_1,e_1$
computes $f$.  At $a\mathbf1+b$ each pair difference is
$b_{2j}-b_{2j-1}$.  If fewer than $m$ coordinates of $b$ are
nonzero, one of the $m$ disjoint pairs has both entries zero,
so the product vanishes.

Replace every matrix factor by that factor plus $\varepsilon I$.
All regularized factors are generically invertible over $K$,
because their determinant is monic in $\varepsilon$.
Call the resulting scalar polynomial $f_\varepsilon$.
At the zero point a regularized pair is
\[
 (L(0)+\varepsilon I)(R(0)+\varepsilon I)
  =\begin{pmatrix}\varepsilon(1+\varepsilon)&\varepsilon\\
                   0&\varepsilon^2\end{pmatrix},
\]
so $f_\varepsilon(0)=\varepsilon^m(1+\varepsilon)^m\ne0$ in $K$.
Yet every common-shift sample of perturbation support less than $m$
has zero specialization at $\varepsilon=0$.

For the matrix assertion, a minor is nonzero over $K$ exactly when
it has finite valuation, and its reduction is nonzero over $F$
exactly when its valuation is zero.  This proves the two rank
criteria.  A unit $r$-minor gives $r$ columns invertible over $R$,
so their span is $R^r$.  Conversely, if the columns span $R^r$,
their reductions span $F^r$, which supplies a unit minor.
\hfill$\square$

\paragraph{The resulting special-fiber gate.}
In the scalar example the module generated by all coefficients of
$f_\varepsilon$ is $R$, since some coefficient of $f_0=f$ is nonzero.
For any finite collection of the displayed fixed $F$-valued samples
that contains the zero point, their values generate a nonzero ideal $N\subseteq
\varepsilon R$.  Hence $K N=K$ and its saturation
$(KN)\cap R$ is $R$, although all original sampled values are zero.
Saturating by dividing by a power of $\varepsilon$ thus does not
supply a nonzero evaluation of the original black box.
For vector coefficients the correct next target is a compression
with a unit Pl\"ucker coordinate, as in (I.18), or an equivalent
proof of coefficient-span preservation on the special fiber.
This refines the generic Grassmannian condition in I-14; it is not
established merely by flatness of a surrounding Rees family.

\subsection{Open gate I-G01 and the order of further work}
\label{I-G01}

The surviving main gate is an explicit polynomial family of
integer cocharacters of polynomial magnitude which, for every
width-$w$ matrix product with disjoint univariate factors,
enters at least one cone (I.3), or satisfies the exact weaker joint
coefficient-span condition I-14 for the family.  If this is
proved, I-B02 gives polynomial hitting sets, and I-08 gives the
precisely delimited matrix-invariant conclusion.  No finding
above claims that this gate is closed.

The constructions already attempted and resolved as stated are:
direct state weights (I-02); monomial-alphabet lifting followed
by exact evaluation interpolation (I-06); positive Gram basis
enumeration (I-05); independent and correlated Haar repair
(I-07); exterior/permutation expansion and literal subspace
configuration graphs (I-11); and positive coefficient-norm
transfer matrices (I-09), exact universal energy cubature (I-13),
and the arithmetic support test (I-12).  The further repairs now
resolved are joint span preservation (I-14), bounded affine mask
decoding (I-15), exact finite-field row states (I-16), the restricted
invertible-factor sampler (I-17), and epsilon regularization with
its necessary special-fiber condition (I-18).  The positive result remaining from
direct graph isolation is I-01.  A next substantive construction
would have to preserve the monomial quotient and coefficient
dependences while compressing the matroid-cone search, rather
than merely isolate another term in a cancellation-prone
expansion.

% END SOURCE work/isolation/findings.tex

% BEGIN SOURCE work/purity/findings.tex
% Purity/zero-status contribution, run CUE_GCT-20261121-purity.
% New proved deductions, not claims of priority over the literature.
\section{Purity, evaluation maps, and the precise derandomization gate}

\subsection{P-01: a quantitative Gram criterion}
\paragraph{Statement P-01.}
Let $\Omega$ be a finite probability space, let $\phi_1,\ldots,\phi_s:\Omega\to\mathbb C$, and suppose
\[
 \mathbb E|\phi_j|^2\ge a>0,\qquad
 |\mathbb E\phi_i\overline{\phi_j}|\le b\quad(i\ne j).
\]
Then every $c\in\mathbb C^s$ satisfies
\[
 \mathbb E\left|\sum_jc_j\phi_j\right|^2
 \ge \{a-(s-1)b\}\sum_j|c_j|^2.
\]
Consequently $a>(s-1)b$ implies that the coefficient-to-evaluation map $\mathbb C^s\to\mathbb C^\Omega$ is injective and that every nonzero combination is nonzero at a sampled point.
\paragraph{Proof.}
Expand the square. The diagonal sum is at least $a\sum_j|c_j|^2$, and the absolute value of the off-diagonal sum is at most $b\sum_{i\ne j}|c_i||c_j|$. Cauchy--Schwarz gives $(\sum_j|c_j|)^2\le s\sum_j|c_j|^2$, so the off-diagonal sum is at most $b(s-1)\sum_j|c_j|^2$. This proves the inequality. If a nonzero combination vanished at every sampled point, the left side would vanish while the right side would be positive, a contradiction. $\square$

For multilinear $f=\sum_{S\in\mathcal S}c_Sx_S$ on the sign cube and an $\varepsilon$-biased distribution (meaning $|\mathbb E x_T|\le\varepsilon$ for nonempty $T$), take $\phi_S=x_S$. Then $a=1$, $b=\varepsilon$, because $x_Sx_T=x_{S\triangle T}$. Thus $\varepsilon(|\mathcal S|-1)<1$ is a sufficient hitting criterion. The sharper coefficient-dependent criterion is
\[
 \mathbb E|f|^2\ge (1+\varepsilon)\|c\|_2^2-\varepsilon\|c\|_1^2>0.
\]
This criterion depends on coefficient energy or sparsity; ROABP width alone is not such a bound.

\subsection{P-02: an explicit purity-based hitting set in the sparse regime}
\paragraph{Statement P-02.}
For positive integers $n,s$, choose an odd prime $p>16n^2s^2$. Let $\chi$ be the quadratic character of $\mathbb F_p$, and put $b_i(t)=\chi(t+i)$ when $t+i\ne0$, and $b_i(t)=1$ otherwise, for $1\le i\le n$. The multiset
\[
 H_{n,s}=\{(b_1(t),\ldots,b_n(t)):t\in\mathbb F_p\}\subset\{\pm1\}^n
\]
is a hitting set for every nonzero $s$-sparse multilinear polynomial over any field of characteristic zero. Taking the first such prime gives cardinality $O(n^2s^2)$ and an elementary deterministic construction in time polynomial in $n,s$.
\paragraph{Proof.}
The quadratic-character construction and its bias bound are AGHP, \S4, Theorem~1, Construction~2, and Proposition~2 (P-B01). For completeness, the short form required here is as follows. For nonempty $T\subseteq[n]$, the squarefree polynomial $\prod_{i\in T}(t+i)$ has $|T|$ distinct roots since $p>n$. Weil's character-sum bound gives a sum of absolute value at most $(|T|-1)\sqrt p$. Replacing each zero character value by $1$ changes at most $|T|$ terms by at most $1$. Therefore
\[
 \left|\frac1p\sum_t\prod_{i\in T}b_i(t)\right|
 \le \varepsilon:=\frac{n-1}{\sqrt p}+\frac np.
\]
Our choice gives $\varepsilon(s-1)<1$, so P-01 proves the assertion over $\mathbb C$. For any set of at most $s$ monomials the evaluation matrix has integer entries; injectivity over $\mathbb C$ implies an integer maximal minor is nonzero, hence it remains nonzero over every characteristic-zero field. Bertrand's theorem supplies a prime less than twice the chosen threshold (enlarging an endpoint by one if necessary). Trial division and Euler's criterion already yield polynomial time in the numerical parameters $n,s$; no GRH is needed. $\square$

The word ``sparse'' is essential. This applies to sparse ROABPs, but an $n$-variate width-one ROABP can have $2^n$ nonzero coefficients.

\subsection{P-03: two small-width counterexamples and the attempted repair}
\paragraph{Statement P-03(a).}
Over a field of characteristic different from $2$, a subset $H\subseteq\{\pm1\}^n$ hits all multilinear width-one ROABPs if and only if $H=\{\pm1\}^n$.
\paragraph{Proof.}
If $a\notin H$, the nonzero polynomial
\[
 f_a(x)=\prod_{i=1}^n(1+a_ix_i)
\]
is a width-one ROABP. For $x\in\{\pm1\}^n$ it is nonzero exactly when $x=a$, where its value is $2^n$. Thus $H$ misses it. Conversely the sign-cube evaluation matrix for multilinear monomials is the tensor product of $\begin{pmatrix}1&1\\1&-1\end{pmatrix}$ and is invertible when $2$ is invertible. Hence the full cube hits all nonzero multilinear polynomials. $\square$

\paragraph{Statement P-03(b).}
Allowing a common affine transformation $x_i=u+vb_i(t)$ of every sign point does not repair the construction for general ROABPs. The polynomial
\[
 V(x_1,x_2,x_3)=(x_2-x_1)(x_3-x_1)(x_3-x_2)
\]
is a nonzero width-three ROABP of individual degree two and vanishes at all such points, even when $u,v$ range over the entire coefficient field.
\paragraph{Proof.}
Among three coordinates taking values in $\{u-v,u+v\}$, two coincide. Thus one factor vanishes. Nonzeroness follows from its displayed factorization in a polynomial ring over a field. The explicit ROABP is
\[
 [1\quad x_1\quad x_1^2]
 \begin{pmatrix}
 x_2&x_2^2&0\\-1&0&x_2^2\\0&-1&-x_2
 \end{pmatrix}
 \begin{bmatrix}x_3^2\\-x_3\\1\end{bmatrix}=V.
\]
Expansion verifies this identity in $\mathbb Z[x_1,x_2,x_3]$, so it holds in every characteristic. $\square$

These are refutations of particular proposed transfers, not refutations of polynomial-size ROABP hitting sets over a larger field or of weight generators with richer output alphabets.

\subsection{P-04: curve sampling after a surviving restriction}
\paragraph{Statement P-04.}
Let $C/\mathbb F_q$ be a smooth projective geometrically connected curve of genus $g$, let $r_i\in\mathbb F_q(C)$, and choose effective divisors $D_i$ containing their polar divisors. For multidegrees $d_i$, put $D=\sum_i d_iD_i$ and let $B$ be the union of supports of the $D_i$. If
\[
 q+1-2g\sqrt q-\#B(\mathbb F_q)>\deg D,
\]
then evaluations at $C(\mathbb F_q)\setminus B$ hit every polynomial $f$ of individual degrees at most $d_i$ whose pullback $f(r_1,\ldots,r_n)$ is nonzero in $\mathbb F_q(C)$.
\paragraph{Proof.}
Every monomial in the pullback has poles bounded by $D$, and the same is true of their sum. If the pullback is a nonzero constant it never vanishes. Otherwise its principal divisor has degree zero, so the degree of its zero divisor equals that of its pole divisor and is at most $\deg D$. In particular it vanishes at at most $\deg D$ rational points outside $B$. Weil's bound (P-B02) supplies at least $q+1-2g\sqrt q-\#B(\mathbb F_q)$ available rational points, which is larger. $\square$

This separates the two requirements: purity supplies enough rational points; it does not prove that the restriction to $C$ is nonzero. An explicit algorithm also needs an explicit model of $C$ and an algorithm for enumerating its rational points with the stated complexity. For $C=\mathbb P^1$ this enumeration is immediate, and the assertion reduces to univariate root counting.

\subsection{P-05: a low-width annihilator for a union of curves}
\paragraph{Statement P-05.}
Let $C_j/k$ ($1\le j\le M$) be smooth projective geometrically connected curves and let $r_{1j},r_{2j}\in k(C_j)$ have poles bounded by a common effective divisor $D_j$ of degree $\delta_j\ge1$. There is a nonzero polynomial $P\in k[X,Y]$, of degree at most $2\sum_j\delta_j$ in each variable, with
\[
 P(r_{1j},r_{2j})=0\quad\text{for every }j.
\]
Moreover $P$ has a two-variable ROABP of width at most $2\sum_j\delta_j+1$. Thus a union of such curve images does not hit that entire width/degree class.
\paragraph{Proof.}
For each $j$, put $b_j=2\delta_j$. The $(b_j+1)^2$ monomials $r_{1j}^a r_{2j}^b$ ($0\le a,b\le b_j$) lie in $L(2b_jD_j)$. The elementary divisor-space bound $\dim_k L(E)\le\deg E+1$ for effective $E$ gives
\[
 \dim L(2b_jD_j)\le2b_j\delta_j+1=4\delta_j^2+1<(2\delta_j+1)^2.
\]
For clarity, this divisor-space bound follows after scalar extension to an algebraic closure: imposing vanishing at any $\deg E+1$ distinct points away from $E$ is injective, since a nonzero section cannot have more zeros than poles. Base extension preserves dimension. Hence there is a nontrivial $k$-linear relation, defining a nonzero formal polynomial $P_j(X,Y)$ of bidegree at most $(b_j,b_j)$ that vanishes on $C_j$. Set $P=\prod_jP_j$. The polynomial ring is a domain, so $P\ne0$; its individual degrees are at most $B_0:=\sum_jb_j$. Finally write $P(X,Y)=\sum_{a=0}^{B_0}X^a Q_a(Y)$ and represent it by the row $[1,X,\ldots,X^{B_0}]$ times the column $(Q_0(Y),\ldots,Q_{B_0}(Y))^{\mathsf T}$. This is an ROABP of width $B_0+1$. $\square$

The quantitative conclusion is a lower bound on a specific curve-based construction. It does not exclude a family whose total pole degree is much larger than the target ROABP width or degree.

\subsection{P-06: what general Frobenius orthogonality proves}
\paragraph{Statement P-06.}
Let $U/\mathbb F_q$ be a smooth geometrically connected affine curve. Let $\mathcal F_1,\ldots,\mathcal F_s$ be geometrically irreducible lisse $\overline{\mathbb Q}_\ell$-sheaves, pure of weight zero for a fixed complex embedding, pairwise geometrically nonisomorphic. Set $t_i(x)=\operatorname{tr}(\operatorname{Frob}_x\mid\mathcal F_i)$, and
\[
 B=\max_{i,j}\dim H_c^1(U_{\overline{\mathbb F}_q},\mathcal F_i\otimes\mathcal F_j^\vee).
\]
If $q>s^2B^2$, the functions $t_1,\ldots,t_s$ on $U(\mathbb F_q)$ are linearly independent. More precisely, with $Q=\#U(\mathbb F_q)$,
\[
 \frac1Q\sum_x\left|\sum_i c_it_i(x)\right|^2
 \ge\frac{q-sB\sqrt q}{Q}\sum_i|c_i|^2.
\]
\paragraph{Proof.}
Purity of weight zero identifies the complex conjugate of a Frobenius trace with the trace of the dual representation. For $\mathcal G_{ij}=\mathcal F_i\otimes\mathcal F_j^\vee$, the trace formula gives
\[
 \sum_{x\in U(\mathbb F_q)}t_i(x)\overline{t_j(x)}
 =\sum_{a=0}^2(-1)^a\operatorname{tr}(\operatorname{Frob}_q\mid H_c^a(U_{\bar{\mathbb F}_q},\mathcal G_{ij})).
\]
Here $H_c^0=0$. By duality and geometric irreducibility, $H_c^2$ is zero for $i\ne j$ and is $\overline{\mathbb Q}_\ell(-1)$ for $i=j$: in the latter case its invariant scalar endomorphisms are fixed by arithmetic Frobenius. Thus the $H_c^2$ contribution is $q\delta_{ij}$. Deligne's weight theorem bounds the absolute values of the $H_c^1$ eigenvalues by $\sqrt q$, so the error in every Gram entry is at most $B\sqrt q$. If $E$ is the error matrix, $|c^*Ec|\le B\sqrt q(\sum|c_i|)^2\le sB\sqrt q\sum|c_i|^2$. Dividing by $Q$ gives the result; the displayed strict inequality also implies $Q>0$. $\square$

For an explicit conductor budget, let $\overline U$ have genus $g$ and let $b=\#(\overline U\setminus U)(\overline{\mathbb F}_q)$ count geometric boundary points. Write $r_i=\operatorname{rank}\mathcal F_i$ and $\sigma_{ij}=\sum_{x\in\overline U\setminus U}\operatorname{Swan}_x(\mathcal F_i\otimes\mathcal F_j^\vee)$. Then the Grothendieck--Ogg--Shafarevich formula gives the exact expression
\[
 \dim H_c^1(U_{\bar{\mathbb F}_q},\mathcal G_{ij})
 =\delta_{ij}+r_ir_j(2g-2+b)+\sigma_{ij}.
\]
Indeed $\chi_c(U,\mathcal G_{ij})=r_ir_j(2-2g-b)-\sigma_{ij}$, while the proof above gives $h_c^0=0$ and $h_c^2=\delta_{ij}$. This proves the displayed expression without discarding the conductor term. The Euler--Poincar\'e source is SGA~5, Expos\'e~X, Theorem~7.1 and Corollary~7.12, formula~(7.13), with the $\ell$-adic form also written in Katz--Sarnak, proof of Lemma~10.1.12, p.~296 (P-B14). In particular, if all tensor sheaves are tame and $r_i\le R$, then $B\le1+R^2\max(0,2g-2+b)$ is a sufficient bound. A claimed polynomial-size trace sampler must bound these geometric quantities polynomially as well as $s$.

The trace formula and weights are cited as P-B03; the passage from cohomological estimates to the explicit Gram inequality is the deduction here. The condition uses the dimensions of the relevant cohomology groups. Saying only ``equidistribution'' without controlling these quantities does not supply the claimed complexity. Tensor products used to model many monomials can increase both their number $s$ and $B$.

\subsection{Geometric dictionary for P-01--P-06}
For a monomial set $\mathcal S$, the coefficient space is the affine scheme $\operatorname{Spec}\mathbb Z[c_S:S\in\mathcal S]$. Evaluation on a finite set is a linear morphism of free modules. Injectivity over characteristic zero in P-01/P-02 is witnessed by a nonzero integral maximal minor, so it holds after base change to $\mathbb Z[1/\Delta]$ for some nonzero minor $\Delta$, but need not hold in every positive characteristic. P-03(a) is the idempotent/Lagrange basis of the finite etale sign-cube algebra over $\mathbb Z[1/2]$. P-03(b) is the alternating exterior product of three sections of the rank-three bundle generated by $1,X,X^2$; its vanishing is rank at most two.

In P-04 the map $k[x_1,\ldots,x_n]\to k(C)$ is pullback along a rational morphism to affine space, and the degree budget is a divisor controlling the corresponding section of $\mathcal O_C(D)$. P-05 concerns the nontrivial kernel of that pullback, detected by dimension of global sections. In P-06 the functions are images under the trace-of-Frobenius map on the Grothendieck group of constructible $\ell$-adic complexes. Tensor product becomes pointwise multiplication, duality becomes complex conjugation for weight zero, and compactly supported pushforward to $\operatorname{Spec}\mathbb F_q$ gives the trace sum. Purity bounds the nonconstant part of this pushforward. These assertions do not posit an unconstructed motive for an arbitrary ROABP. Exact trace/weight locators are SGA $4\frac12$, Rapport sur la formule des traces, Theorem~3.2; Deligne, Weil II, Theorem~3.3.1 and \S3.3.10; and Katz--Sarnak, \S\S9.0.13--9.0.16. The general equidistribution statement is Katz--Sarnak, Theorem~9.2.6 (compare Weil II, Theorem~3.5.3). Further SGA/EGA locators are supplied in the global dictionary.

\subsection{P-07: the next construction---bounded-support isolating weights}
\paragraph{Statement P-07.}
Let $S\subseteq\{0,\ldots,d\}^n$ have at most $s$ elements. Fix a prime
\[
 p>\max\{2,d,n\tbinom{s}{2}\}.
\]
For each $t\in\mathbb F_p$, let
\[
 w_i(t)=1+\operatorname{rep}_p(t^i)\in\{1,\ldots,p\},\qquad 1\le i\le n.
\]
Some member of this explicit $p$-element family assigns pairwise different integer weights $w(t)\cdot\alpha$ to the elements of $S$. In particular it isolates their minimum. Over any field containing $ndp+1$ distinct scalars, the substitutions $x_i=z^{w_i(t)}$ followed by evaluation at these scalars form a hitting set of size at most $p(ndp+1)$ for all $s$-sparse polynomials of individual degree at most $d$.
\paragraph{Proof.}
For distinct $\alpha,\beta\in S$, equality of their integer weights implies
\[
 0=\sum_i(\alpha_i-\beta_i)(1+t^i)\quad\text{in }\mathbb F_p.
\]
This polynomial in $t$ is nonzero: at least one nonconstant coefficient $\alpha_i-\beta_i$ is nonzero modulo $p$, because its absolute value is at most $d<p$. Its degree is at most $n$, so it has at most $n$ roots. The union of all pair-collision sets has size at most $n\binom{s}{2}<p$. Choosing $t$ outside that union proves the weight assertion. Under that substitution, every nonzero coefficient remains on a distinct power of $z$, so the resulting univariate polynomial is nonzero, of degree at most $ndp$. It cannot vanish at all $ndp+1$ distinct scalar values. The union of all evaluation points therefore hits every polynomial in the stated class. $\square$

This gives a polynomial-size family for bounded sparsity and allows arbitrary coefficients and cancellation; it does not require purity, because elementary root counting is already stronger for the polynomial collision predicates involved. Its geometric form is a family of one-parameter subgroups of $\mathbb G_m^n$, and the successful member lies in the interior of a normal cone exposing a unique vertex of the support's Newton polytope. For general width-$w$ ROABPs the support size can be exponential even at $w=1$, so substituting $s=\operatorname{poly}(w,n,d)$ here is unjustified. The remaining transfer must exploit coefficient-space structure, as in basis isolation, rather than enumerate support pairs.

\subsection{P-08: a precise conditional arithmetic encoding reduction}
\paragraph{Statement P-08 (conditional reduction).}
Let $\mathcal C_N$ be a polynomial class. Suppose the following data can be constructed in time polynomial in $N$: a common finite Galois extension $E_N/\mathbb Q$, its nonzero discriminant, and a map $p\mapsto h_N(p)$ from unramified rational primes to evaluation points. Suppose $[E_N:\mathbb Q]$ and $\log|\Delta_{E_N}|$ are polynomially bounded, that $h_N(p)$ is computable in polynomial time in $N+\log p$, and that for every nonzero $f\in\mathcal C_N$ some conjugacy class $C_f\subseteq\operatorname{Gal}(E_N/\mathbb Q)$ has
\[
 \operatorname{Frob}_p\in C_f\quad\Longrightarrow\quad f(h_N(p))\ne0
\]
for every unramified $p$. Under GRH for $\zeta_{E_N}$ these data yield a polynomial-size, polynomial-time hitting set for $\mathcal C_N$.
\paragraph{Proof.}
Lagarias--Odlyzko, Corollary~1.2 (P-B04), gives an unramified prime in every conjugacy class with
\[
 p\le c\,\log^2|\Delta_{E_N}|\,\log^4\log|\Delta_{E_N}|
\]
for an effective absolute constant, with the trivial extension treated separately. Increasing the bound to cover the finitely many small discriminants gives a polynomial bound $B(N)$. Enumerate all primes up to this bound by trial division, discard primes dividing the supplied discriminant, and evaluate $h_N$ on the rest. The set has polynomial size and is polynomial-time constructible; it contains a point from $C_f$ for every nonzero $f$, proving the claim. $\square$

If $E_N/\mathbb Q$ is abelian, factorization into Dirichlet $L$-functions makes this GRH assumption a consequence of the stipulated Alp\"oge theorem. For arbitrary $E_N$, it follows conditionally from Masullo's asserted Hecke GRH, by using the trivial Hecke character of $E_N$. The extension and evaluation map must be common to the polynomial class: an extension depending on the hidden polynomial would not construct a black-box hitting set. No construction satisfying these hypotheses for general ROABPs is supplied by the stated RH/GRH theorems or by the finite-field purity argument. The exact remaining gate is the explicit common encoding with polynomial discriminant, together with the Frobenius-to-nonvanishing implication. Deterministic quadratic nonresidues and modular square roots are the consequences assigned to Prompt~28 and are not rederived here.

\subsection{P-09: approximate positive cubature already costs the full tensor dimension}
\paragraph{Statement P-09 (adapted trace argument).}
Let $n\ge1$, $\delta\ge0$, $D=(\delta+1)^n$, and let $\xi_1,\ldots,\xi_M\in\mathbb C^n$ have positive real weights $w_j$. Suppose that constants $0<a\le b$ satisfy
\[
 a\|\operatorname{coeff} f\|_2^2
 \le \sum_{j=1}^M w_j|f(\xi_j)|^2
 \le b\|\operatorname{coeff} f\|_2^2
 \tag{P09.1}
\]
for every width-one product $f(x)=\prod_{i=1}^n p_i(x_i)$ with $\deg p_i\le\delta$. Then
\[
 M\ge \frac{a}{b}D.
 \tag{P09.2}
\]
In particular a relative error $0\le\varepsilon<1$ forces
$M\ge D(1-\varepsilon)/(1+\varepsilon)$. Nonnegative weights are allowed after zero weights are removed.
\paragraph{Proof.}
Put $v(z)=(1,z,\ldots,z^\delta)$ and
$e_j=\bigotimes_{i=1}^n v(\xi_{j,i})\in\mathbb C^D$.
Coefficient vectors are paired bilinearly with $e_j$ for evaluation.
Since $e_j\ne0$ and is decomposable, the vector
$c^{(j)}=\overline{e_j}/\|e_j\|_2$ is the coefficient vector of a width-one product, of coefficient norm one. Its value at $\xi_j$ is $\|e_j\|_2$. All summands in (P09.1) are nonnegative, so the upper bound applied to this product gives
\[
 w_j\|e_j\|_2^2\le b.
\]
Apply the lower bound to each of the $D$ monomials $x^\alpha$, $\alpha\in\{0,\ldots,\delta\}^n$, each itself a normalized width-one product, and sum. This gives
\[
 aD\le\sum_\alpha\sum_j w_j|\xi_j^\alpha|^2
 =\sum_j w_j\|e_j\|_2^2\le Mb,
\]
which proves the result. Equivalently, averaging over independently chosen Haar unit vectors in the $n$ coefficient factors gives covariance $I_D/D$ and the same trace bound. The finite basis sum avoids any integration assumption. $\square$

The proof also applies to arbitrary decomposable measurement vectors, not only polynomial evaluations. It is related to the rank-one measurement obstruction for ordinary $\ell_2$ restricted isometry: Cai and Zhang (2015, Lemma~2.1, p.~109; P-B15) prove an ambient-dimension obstruction for their Gaussian matrix model and use a different $\ell_1$ measurement condition instead. The weighted tensor statement above is proved here directly; no priority claim is made. Any proposed purity sampler that uniformly approximates coefficient energy for all bounded-width ROABPs includes width one and therefore meets this barrier. The barrier concerns the quantitative two-sided energy requirement; it does not exclude hitting sets.

\subsection{P-10: positivity alone, tensor-train kernels, and exact width-one sampling}
\paragraph{Statement P-10.}
Let $n,w\ge1$, $\delta\ge0$, and let $\mathcal T_w\subseteq\mathbb C^{(\delta+1)^n}$ be the cone of coefficient tensors whose consecutive prefix/suffix flattenings all have rank at most $w$, for a fixed variable order. Let $A_Hc=(f_c(\xi_j))_{j=1}^M$ and let all sampling weights be positive. The following are equivalent:
\[
 H\text{ hits every nonzero tensor in }\mathcal T_w;
 \qquad \ker A_H\cap\mathcal T_w=\{0\};
 \qquad \inf_{c\in\mathcal T_w,\,\|c\|_2=1}
       \sum_jw_j|f_c(\xi_j)|^2>0.
 \tag{P10.1}
\]
The cone is exactly the coefficient class of width-at-most-$w$ ROABPs in that order. If $a$ and $b$ are the minimum and maximum sampled energies on its unit sphere, every such hitting set satisfies
\[
 \frac ba\ge\frac{(\delta+1)^n}{M}.
 \tag{P10.2}
\]
If $q_w=\dim\mathcal T_w$, every complex hitting set has at least $q_w$ points, and a nonempty Zariski-open set of $q_w$-tuples of evaluation points consists of hitting sets. Moreover
\[
 q_w\le n(\delta+1)w^2.
 \tag{P10.3}
\]
These last assertions establish existence, with integer choices possible, but do not provide an efficient deterministic construction or a coordinate-height bound.

For width one there is an exact explicit answer over every field containing at least $n\delta+1$ elements: $n\delta+1$ distinct diagonal points $(t,\ldots,t)$ hit all nonzero products, and no set with at most $n\delta$ arbitrary evaluation points does so.
\paragraph{Proof.}
Flattening rank bounds are the vanishing of homogeneous minors, so $\mathcal T_w$ is closed. Rank factorization gives the standard tensor-train description. One way to see the converse explicitly is to let $U_i$ be the column space of the $i$th flattening. Then
$U_i\subseteq U_{i-1}\otimes\mathbb C^{\delta+1}$: contractions in the remaining factors produce the columns, and the preceding left support contains each such contraction. Choosing bases in these spaces writes their inclusions as cores of sizes at most $w\times(\delta+1)\times w$, with one-dimensional boundary cores. Their contraction is the original coefficient tensor. Conversely every core contraction factors each prefix flattening through its bond space, giving the rank bound. This is the exact rank-factorization result of Oseledets (2011, Theorem~2.1, pp.~2298--2299; P-B17), restated for polynomial coefficients. Zero-padding the cores gives a surjective polynomial parameterization with at most $n(\delta+1)w^2$ parameters, proving (P10.3).

The first two conditions in (P10.1) are identical by definition. With positive weights, a sampled energy vanishes exactly when all evaluations vanish. The intersection of the closed cone with the unit sphere is compact. Thus absence of a zero of the continuous energy function there is equivalent to a strictly positive minimum. Its maximum is finite as well. Since $\mathcal T_w$ contains all product tensors, P-09 gives (P10.2).

For the generic assertion let $X=\mathbb P(\mathcal T_w)$, of dimension $q_w-1$. Imposing $M<q_w$ homogeneous linear equations on its affine cone leaves dimension at least $q_w-M\ge1$, by the principal ideal dimension bound; over $\mathbb C$ the common zero set therefore contains a nonzero tensor. This proves the lower bound. Evaluation hyperplanes have no common base point on $X$, because a nonzero polynomial is nonzero at some point of $\mathbb C^n$. At each step choose an evaluation hyperplane containing no irreducible component of the remaining projective intersection. Such a choice exists: for each of finitely many components, choose one nonzero coefficient point on it; its polynomial excludes a proper hypersurface of possible evaluation points. Intersecting with the chosen hyperplane lowers every positive component's dimension by one and removes all zero-dimensional components. After $q_w$ steps the intersection is empty. Finally the incidence subset
\[
 \{([c],\xi_1,\ldots,\xi_{q_w}):
      [c]\in X,\ f_c(\xi_j)=0\text{ for all }j\}
 \subseteq X\times(\mathbb A^n)^{q_w}
\]
is closed and its projection is closed because $X$ is projective. The choices just made show that this projected bad set is proper. Its complement is the asserted nonempty open set. Integer tuples are Zariski dense in affine complex space, so some integer choice lies in that open set. The argument does not bound how far one must search.

For the explicit width-one upper bound, the restriction
$f(t,\ldots,t)=\prod_i p_i(t)$ is nonzero and has degree at most $n\delta$, hence cannot vanish at all $n\delta+1$ distinct scalars. For the lower bound, partition any collection of $M\le n\delta$ sample indices into $n$ sets $J_i$, each of size at most $\delta$, and set
\[
 f(x)=\prod_{i=1}^n\prod_{j\in J_i}(x_i-\xi_{j,i}).
\]
This is a nonzero product of monic univariate factors, has individual degree at most $\delta$, and vanishes at every listed point. Empty factors equal one. $\square$

The width-one construction and dimension threshold agree with Parthasarathy's maximal completely entangled subspaces (2004, Propositions~1.3--1.4 and Theorem~1.5; P-B16). The kernel of the diagonal evaluation map is such a subspace of dimension $D-n\delta-1$. In this terminology ``completely entangled'' means precisely that the kernel contains no nonzero decomposable tensor; it is not a physical assumption on the polynomials. For larger tensor-train ranks, (P10.1) is the precise weaker gate. Proving positivity directly could avoid the exponential constant-distortion requirement, but generic existence and compactness do not construct that gate efficiently.

\subsection{P-11: an explicit number-field point is already an unconditional hitting-set encoding}
\paragraph{Statement P-11.}
Let $\mathcal C\subseteq\mathbb Q[x_1,\ldots,x_n]$ have individual degree at most $\delta$. Suppose a common algebraic number $\theta$, polynomials $A_i(T)\in\mathbb Z[T]$ of degree at most $a$, and an integer $b\ne0$ are supplied such that
\[
 f\ne0,\ f\in\mathcal C\quad\Longrightarrow\quad
 f\bigl(A_1(\theta)/b,\ldots,A_n(\theta)/b\bigr)\ne0.
 \tag{P11.1}
\]
Put $d=n\delta$. The completely explicit integer set
\[
 H=\{(uA_1(t),\ldots,uA_n(t)):
        0\le u\le d,\ 0\le t\le da\}
 \tag{P11.2}
\]
hits $\mathcal C$ and has size at most $(d+1)(da+1)$. If $b=1$, the single grid $t=0,\ldots,da$ with $u=1$ suffices. These conclusions need no GRH, no discriminant bound, no coefficient-height bound on $f$, and no algorithm for arithmetic in $\mathbb Q(\theta)$.

Conversely suppose an explicitly given integer hitting set $H_1,\ldots,H_M$, $M\ge1$, hits a class $\mathcal C\subseteq\mathbb Q[x]$ of total degree at most $d$. Polynomial-time interpolation and a prime $\ell>d(M-1)+1$ construct data as in (P11.1) with $\theta$ a primitive $\ell$th root of unity. The degree, coordinate degrees, coefficient bit lengths, and logarithmic discriminant are polynomial in $d,M$ and the bit lengths of the given points, choosing $\ell=O(d(M-1)+2)$. Thus the two explicit encodings are polynomially interconvertible over $\mathbb Q$.
\paragraph{Proof.}
For nonzero $f\in\mathcal C$, define
\[
 G_f(U,T)=f(UA_1(T),\ldots,UA_n(T)).
\]
It is nonzero, since its evaluation at $(U,T)=(1/b,\theta)$ is nonzero by hypothesis. Its separate degrees are at most $d$ in $U$ and $da$ in $T$. If it vanished on the entire rectangular grid in (P11.2), then, for each fixed $t$, univariate interpolation at the $d+1$ values of $u$ would force all its $U$-coefficients to vanish at $t$. Interpolation at the $da+1$ values of $t$ would then make every coefficient polynomial zero, a contradiction. When $b=1$, $G_f(1,T)$ is itself nonzero at $\theta$ and has degree at most $da$, proving the smaller bound. If $|\operatorname{coeff} A_i|\le2^K$, the coordinates in (P11.2) have bit length at most
\[
 K+\left\lceil\log_2(d+1)\right\rceil
  +\left\lceil\log_2(a+1)\right\rceil
  +a\left\lceil\log_2\max(1,da)\right\rceil+2,
\]
so polynomial input sizes give polynomial bit complexity as well as polynomial cardinality.

For the reverse direction interpolate rational polynomials $a_i(T)$ of degree at most $M-1$ with $a_i(j)=H_{j+1,i}$ for $0\le j<M$. A common denominator $b=(M-1)!$ clears their coefficients, because each Lagrange denominator is $(-1)^{M-1-j}j!(M-1-j)!$, which divides $b$. Hence $A_i=ba_i$ are integral and have polynomial coefficient bit lengths; explicitly, if the coordinates of $H$ have absolute values at most $2^K$, the Lagrange product expansion gives coefficient bit length $O(K+M\log(M+1))$. For nonzero $f\in\mathcal C$, the polynomial $f(a_1(T),\ldots,a_n(T))$ is nonzero because it is nonzero at at least one interpolation node. Its degree is at most $d(M-1)$. Choose a prime $\ell>d(M-1)+1$, and let $\theta=\zeta_\ell$. Its minimal polynomial over $\mathbb Q$ is $\Phi_\ell(T)=1+T+\cdots+T^{\ell-1}$: irreducibility follows from Eisenstein applied after $T\mapsto T+1$. Its degree exceeds $d(M-1)$, so it cannot divide the nonzero rational polynomial above. This proves (P11.1). Bertrand's bound and trial division choose $\ell$ of the stated size. The polynomial discriminant has absolute value $\ell^{\ell-2}$: differentiate $(T^\ell-1)/(T-1)$ at its roots and use $\prod_{a=1}^{\ell-1}(1-\zeta_\ell^a)=\Phi_\ell(1)=\ell$ in the product of the derivatives. The field discriminant divides this order discriminant. On every conjugate of $\theta$, $|a_i(\theta)|$ is bounded by the sum of the absolute values of its coefficients. Thus all claimed representation and height bounds are polynomial. $\square$

In particular, if a primitive-element presentation of a degree-$e$ number field is given with coordinate polynomials of degree below $e$, the remaining arithmetic gate is exactly the nonvanishing condition (P11.1), and P-11 uses $a\le e-1$. The argument does not assert that primitive-element coordinates of small degree and polynomial bit size can be obtained from an arbitrary implicit field description. It also does not extend the reverse direction to unrestricted complex coefficients: minimal-polynomial divisibility is a statement over $\mathbb Q$.

There are concrete restricted encodings. For width-one products over $\mathbb Q$, take $\theta=\zeta_\ell$ with $\ell-1>\delta$ and every coordinate equal to $\theta$; no nonzero factor of degree at most $\delta$ vanishes there. This fails already for the width-two polynomial $x_1-x_2$ when $n\ge2$. P-07 provides an explicit integer hitting set of polynomial size and polynomial coordinate bit length for bounded sparsity, so the reverse construction gives a polynomial-degree cyclotomic encoding for that larger restricted class as well. These constructions do not give a new route to general ROABP PIT: they show exactly why supplying a common effective nonvanishing encoding has already supplied essentially the missing hitting set. The strictly prime-dependent encoding in P-08 need not have the explicit polynomial coordinate representation assumed here and is not covered by this simplification.

% END SOURCE work/purity/findings.tex

% BEGIN SOURCE work/purity/zero_status.tex
\section{What the stipulated RH and GRH inputs actually discharge}

\paragraph{Input convention.}
Write A for the supplied Alp\"oge theorem: RH for $\zeta$ and GRH for all Dirichlet $L$-functions. The relevant exported declarations are \texttt{riemannHypothesis}, \texttt{LFunction\_eq\_zero\_re}, and \texttt{LFunction\_eq\_zero\_re\_of\_isPrimitive} in \texttt{ALPOGE\_Main.lean}; their predicates are in \texttt{ALPOGE\_Statements.lean}. We use A as stipulated and do not claim to have rebuilt its imported archive. Write M for Masullo's claimed GRH for Hecke $L$-functions over every number field, cited conditionally as the prompt's \S1.2. The supplied index does not contain a numbered \S1.2 statement or the full proof, so no theorem number is invented for M. Theorem~13.3.1.1 (PDF p.~522 in the index) is a trace-formula claim; Conjecture~14.1.2.1 (PDF p.~529) explicitly concerns automorphic $\mathrm{GL}_n$ GRH and simplicity. Neither is silently substituted for M or for a proof of general automorphic GRH.

\subsection{P-Z01: Montgomery's supported pair correlation}
For an even Schwartz function $f$ with $\operatorname{supp}\widehat f\Subset(-1,1)$, set $W(u)=4/(4+u^2)$ and $N_0(T)=T\log T/(2\pi)$. Montgomery's theorem and Fourier identity (1973, \S1, pp.~181--182, equations~(1)--(4); P-B05), combined with A, give
\[
 \lim_{T\to\infty}\frac1{N_0(T)}
 \sum_{0<\gamma,\gamma'\le T}
 f\left(\frac{\log T}{2\pi}(\gamma-\gamma')\right)W(\gamma-\gamma')
 = f(0)+\int_{\mathbb R}f(u)\left[1-\left(\frac{\sin\pi u}{\pi u}\right)^2\right]du.
\]
The sum is over ordered zero indices with multiplicity. Subtracting the identical-index diagonal gives the sine-kernel two-point density. To check the normalization directly, the right side equals
$\widehat f(0)+\int_{-1}^1|u|\widehat f(u)\,du$ under the support restriction, exactly the contribution of Montgomery's delta term at frequency zero and its $|u|$ term. This is now unconditional under the input convention. The unrestricted-support assertion, the complete nearest-neighbour spacing law, and all-height GUE statistics require additional estimates. The missing analytic input is control of the off-diagonal prime correlations occurring beyond the supported range; critical-line location does not provide that estimate in the cited theorem.

\subsection{P-Z02: the Hughes--Rudnick hypothesis correction}
Put $x_{\chi,j}=\gamma_{\chi,j}\log q/(2\pi)$ and $W_f(\chi)=\sum_jf(x_{\chi,j})$. Hughes--Rudnick (2003), Introduction, expressly permits complex $\gamma_{\chi,j}$ in $\rho=1/2+i\gamma_{\chi,j}$ (P-B06). Its analytic mean, variance, and restricted moments are therefore already unconditional before A. For admissible $f$ (Definition~2.1), averaged over the $q-2$ nonprincipal characters modulo a prime $q$, Theorems~3.1, 4.1, and 5.1 respectively give
\[
 \mathbb E W_f\longrightarrow\int f,\qquad
 \operatorname{Var}(W_f)\longrightarrow\int_{-1}^1|u|\widehat f(u)^2\,du,
\]
with supports $[-2,2]$ and $[-1,1]$, and Gaussian centred moments through order $m$ when $\operatorname{supp}\widehat f\subseteq[-\alpha,\alpha]$ and $m\alpha<2$. A makes these ordinates real. It also discharges the explicit GRH in Theorem~8.1, Corollary~8.2 ($\liminf_q\min_{\chi\ne\chi_0}x_{\chi,1}\le1/4$), and Theorem~8.3 (its positive-proportion bound for $\beta\ge0.633$). Here $x_{\chi,1}$ is the first nonnegative scaled ordinate. A finite collection of Gaussian moments is not a Gaussian limit theorem at microscopic scale: the same paper calls the phenomenon mock-Gaussian, and the full random-matrix moments are not all Gaussian.

\subsection{Dependency table}
\begingroup\small
\begin{longtable}{p{0.25\textwidth}p{0.33\textwidth}p{0.34\textwidth}}
\hline
Statement and locator & Exact RH/GRH input or other gate & Status under A and M\\\hline
Montgomery (1973), \S1 theorem and (3)--(4), supported pair correlation & RH for $\zeta$; unchanged $(-1,1)$ Fourier support & Discharged by A (P-Z01).\\
Rudnick--Sarnak (1996), Theorem~1.2, $m=1$ & RH for the fixed primitive Dirichlet $L$-function (or $\zeta$); support $\sum|\xi_j|<2$ & Discharged by A. Theorem~1.1 already proves a smoothed analytic version with complex ordinates.\\
Hughes--Rudnick (2003), Theorems~3.1,4.1,5.1 & No GRH for the analytic statistic; A makes all ordinates real & Already unconditional formulas; real-zero interpretation now supplied by A (P-Z02).\\
Hughes--Rudnick (2003), \S8 & Dirichlet GRH & Discharged by A; original support and constants retained.\\
Gao--Zhao (2020), Theorems~1.1 and~1.3 & GRH for the quadratic/quartic Hecke families over $\mathbb Q(i)$ and auxiliary Hecke characters in Lemma~2.8 & Conditional on M; symplectic density for quadratic family with support $(-2,2)$, unitary for quartic with $(-20/19,20/19)$. A alone does not cover these general Hecke characters.\\
Iwaniec--Luo--Sarnak (2000), Theorems~1.1--1.3 & Standing RH for target cuspidal $\mathrm{GL}_2$ functions, plus Dirichlet RH for extended support; see p.~56, and pp.~88--89 for auxiliary symmetric-square estimates & Beyond A and M for the general family: missing RH for $L(s,f)$. Extra symmetric-square/Rankin--Selberg hypotheses used in individual estimates can be avoided by the averaged replacements explicitly noted in \S4.\\
Iwaniec--Luo--Sarnak (2000), Theorems~1.4--1.5 & RH for target $L(s,\operatorname{sym}^2f)$ ($\mathrm{GL}_3$); auxiliary hypotheses as specified in \S4 & Beyond A and M; missing target symmetric-square RH.\\
Iwaniec--Luo--Sarnak (2000), Theorem~1.9 & Weighted, parity-unsplit analytic one-level statistic, support $(-2,2)$ & Already unconditional; target zero ordinates may be complex.\\
Iwaniec--Luo--Sarnak (2000), Theorem~1.10 & Hypothesis S, a prime exponential-sum estimate & Beyond A and M; the precise missing input is Hypothesis S, not merely another GRH label.\\
Katz--Sarnak, Theorem~9.2.6; Weil II, Theorem~3.5.3 & Purity, the actual geometric/arithmetic monodromy hypotheses, and the stated field-extension limit & Already unconditional in that scope. Growing rank/genus/conductor requires separate uniformity.\\
Sawin (2020), Hypothesis~1.3, Theorem~1.4 (label \texttt{main}), Remark~1.9(1),(5) & $H(n,r_+,r_-,n-1)$ follows from duality and gives fixed-parameter $O(q^{-1/2})$; stronger degree-aspect conclusions require vanishing and stronger uniform Betti bounds & Basic $q$-aspect comparison already unconditional. General stronger cohomology and uniformity gates remain; the unconditional unbalanced case is recorded below.\\
Keating--Snaith / CFKRS general number-field moment predictions & Arithmetic moment asymptotics and off-diagonal terms; CFKRS labels its zeta assertion a conjecture & A and M do not discharge the moment recipe. The finite-$N$ CUE integrals are separate exact theorems.\\
Berry--Keating $H=xp$ and BGS universality heuristics & A self-adjoint spectral realization with the required boundary/domain conditions and appropriate chaotic universality theorem & Beyond A and M; RH alone is not that construction.\\\hline
\end{longtable}
\endgroup

Sawin also proves an unconditional unbalanced case: Corollary~1.8 (repeated at source label \texttt{combined-result}) treats $r_-=1$ and $r_+=r$, with $n>2r+1$, the same stated characteristic exclusions, and error
\[
 O\!\left(\sqrt q\,\bigl(q^{-(p-2r)/(2pr)}(r+2)^{r+1}\bigr)^n n^{r+1}\right),
 \qquad p=\operatorname{char}\mathbb F_q.
\]
For fixed $r,p,q$ satisfying $q^{(p-2r)/(2pr)}>(r+2)^{r+1}$, this error decays exponentially in $n$ up to the displayed polynomial factor. Thus some degree-aspect cases are already proved; the unresolved general moment claim is broader than this restricted theorem.

For Rudnick--Sarnak in degree $m>1$, Theorem~1.2 also needs the target RH and, for $m>3$ in the stated source, their coefficient Hypothesis~(1.7). M covers a Hecke case only after all the relevant identification and other hypotheses have been verified; it is not a blanket automorphic GRH theorem. For ILS, ``Hecke cusp form'' refers to a $\mathrm{GL}_2$ automorphic object; it must not be confused with a one-dimensional Hecke character over a number field.

\paragraph{Quantum-chaos source discipline.}
The connection recalled from memory is the symmetry-class heuristic: chaotic quantum systems without an antiunitary time-reversal symmetry suggest unitary statistics, while spinless time-reversal invariance with $T^2=+1$ suggests orthogonal statistics (the $T^2=-1$ symmetry class is symplectic). The original BGS paper (Bohigas, Giannoni, \& Schmit, 1984, pp.~1--4) concerns chaotic quantum spectra; it does not construct the zeta-zero operator. The bibliographic locator and Berry--Keating's explicit $H_{\rm cl}=XP$ speculation were checked against the primary publication records. Berry and Keating (1999), \emph{The Riemann zeros and eigenvalue asymptotics}, SIAM Review 41, 236--266, and their \emph{$H=xp$ and the Riemann zeros} contribution in \emph{Supersymmetry and Trace Formulae: Chaos and Disorder}, pp.~355--367, are the relevant sources. Detailed internal section/equation locators for the latter were not inspected here, so they are not fabricated.

There is a distinction from bare operator existence. Once real zero ordinates $\gamma_j$ are given, the diagonal operator $He_j=\gamma_je_j$ on $\ell^2$, with domain $\{(a_j):\sum_j\gamma_j^2|a_j|^2<\infty\}$, is self-adjoint. Indeed testing the adjoint against the coordinate vectors forces the same diagonal action and domain. This merely packages the supplied spectrum. It does not supply a specified differential or pseudodifferential Hamiltonian, arithmetic trace formula, boundary conditions, or a proof of chaotic local statistics. Those are the substantive missing constructions in the Berry--Keating/Hilbert--P\'olya direction.

% END SOURCE work/purity/zero_status.tex

\part{The Grothendieck restatement}
\section{A separate Grothendieck dictionary: objects, operations, and proof transport}
\label{sec:grothendieck-dictionary}

This part restates the checked deductions in the requested geometric language. A locator identifies the framework in which a proof lives; it is not a claim that Grothendieck stated a modern complexity result. The original manuscript images are not in the attachment. Cote and scan references below therefore refer to the supplied transcriptions, and the mathematical corrections do not certify an unseen handwritten reading.

\subsection{Exact framework locators and scope}

\paragraph{D-S3: group schemes and weights.}
For an affine group scheme, its representations are comodules over its Hopf algebra. The relevant reference is SGA~3, Expos\'e~I, Proposition~4.7.2. For a diagonalizable group $D(M)$, such comodules are $M$-graded modules, by Proposition~4.7.3; the following remark gives the corresponding graded-algebra description of affine schemes with action. Proposition~4.7.4 gives equivariant splittings when the underlying sequence splits. These statements, inspected in the original Expos\'e~I PDF, give the integral meaning of every torus weight and cocharacter below. For stabilizers, SGA~3, Expos\'e~VI$_B$, Proposition~4.1 gives upper semicontinuity of group-fiber dimension and explicitly cites EGA~IV$_3$, Theorem~13.1.3. The latter locator is corroborated through that primary SGA text; its entire EGA proof has not been reread here.

\paragraph{D-S6: $K_0$, exterior powers, Adams operations, and filtration.}
SGA~6, Expos\'e~V, Definition~2.1 and Example~2.2(b), construct the pre-$\lambda$ structure of locally projective modules over a ringed topos; Lemma~2.2.1 uses a filtration on exterior powers of an exact sequence. Definition~2.4 and equation~(2.4.2) specify the universal product and composition identities. Expos\'e~VI, Corollary~3.2 is the special-$\lambda$-ring assertion for the geometric $K'$ in its stated setting. Expos\'e~V, Proposition~3.2 defines $\gamma$ from $\lambda$; section~3.10, equation~(3.10.1), defines the augmentation filtration. Section~7.1, equation~(7.1.1), and Proposition~7.5 give Adams operations and their ring identities. The electronic transcriptions of the cited Expos\'e~V passages were inspected; the exact VI~3.2 locator was located in the text's index and corroborating primary references. We do not describe an integral character identity as an integral semisimple decomposition.

For split reductive group schemes a useful complementary reference is Serre (1968), Theorem~5, section~3.7, and the explicit $\mathrm{GL}_n$ character ring in section~3.8. The base-change statements in that theorem concern the representation Grothendieck group, not preservation of a characteristic-zero basis of simple objects in every characteristic. Nor do they automatically answer the separate finite-symmetric-group lattice question recorded in the supplied study of cote~133.

\paragraph{D-FGA: images, quotients, and projective parameter spaces.}
For a homogeneous orbit map over $\mathbb Z$, take its schematic image by the kernel of the coordinate homomorphism, as in G-04. Degreewise kernels and images are given by finite matrices and their minors. EGA~II, sections~2.4--2.5, in particular Proposition~2.5.1, gives the projective scheme and associated sheaf of a graded module. The electronic EGA translation was inspected for this locator. FGA, Bourbaki Expos\'e~221, Theorem~3.1 gives the Quot parameter space, Theorem~3.2 its fixed-Hilbert-polynomial projective form, and Proposition~3.8 its Grassmannian realization in the required degree range. The original FGA scan was inspected. These results place the families in parameter spaces; they supply no polynomial bound for finding the relevant minors or their nonvanishing loci.

\paragraph{D-ET: traces, weights, and their arithmetic size.}
SGA~4$\frac12$, \emph{Rapport sur la formule des traces}, Theorems~3.1--3.2 and Proposition~3.3, gives the determinant expression for the $L$-function, the trace formula, and the logarithmic-derivative identity. SGA~5, Expos\'e~X, Theorem~7.1 and Corollary~7.12, formula~(7.13), gives the Euler--Poincar\'e conductor accounting. Deligne's Weil~II, Theorem~3.3.1 and section~3.3.10, supplies the compact-support weight estimate used in P-06; Theorem~3.5.3 gives the equidistribution setting. The Weil~II locators were checked against Katz and Sarnak (1999), sections~9.0.13--9.0.16 and Theorem~9.2.6. The original large Weil~II PDF was not successfully opened here. This is a statement-and-dependency check, not an independent verification of the complete proof of Weil~II.

\subsection{Rectangular characters and orbit-closure multiplicities}

\paragraph{R-01. Tensor pullback of a Schur operation.}
For locally free $E,F$ of ranks $a,n$, the tensor morphism $\mathrm{GL}(E)\times\mathrm{GL}(F)\to\mathrm{GL}(E\otimes F)$ is defined over $\mathbb Z$. The class $s_{(d^n)}[E\otimes F]$ restricts, after passage to characteristic zero, to the sum of $[S_\mu E\otimes S_\nu F]$ with coefficients $g((d^n),\mu,\nu)$. The Haar projector is the complex realization of extraction of that character component. Its finite CUE dimension is a rank of a Schur construction; its normalized limiting constant need not be an integer rank. This is precisely the proof of R-01 transported through D-S3 and D-S6. Cote~133, scan~028, lines~6--17, provides the tensor-power and symmetric-operation language, while scan~019 distinguishes rank from the full character class.

\paragraph{R-02. Two different meanings of positivity.}
The tableau expansion is an integral character identity in the universal special $\lambda$-ring, with the Weyl alternating factor extracting a highest weight. Over $\mathbb C$, the invariant summand of $S_\lambda(E\otimes F)\otimes\det(E)^{-d}\otimes\det(F)^{-d}$ is an actual vector space. A Hermitian metric identifies its projector trace with a sum of squared absolute values. The former is an integral identity; the latter uses complex conjugation and compact integration. D-S6 and cote~133, scan~028, lines~19--24 (the direct-summand question), are the exact conceptual distinction. No cancellation-free combinatorial basis is inferred from the analytic squares.

\paragraph{R-03. Central quotient, disconnected stabilizer, and rank.}
The tensor morphism has the central subgroup $t\mapsto(tI,t^{-1}I)$ in its kernel, a statement of group schemes over $\mathbb Z$. A compact block-Levi branching rule is therefore a different restriction functor. The transpose component of the determinant stabilizer is additional group-scheme data beyond the connected tensor subgroup. The $S_4$ counterexamples compare distinct invariant functors and survive as finite character computations. D-S3 places these groups and actions; cote~133, scan~019, lines~45--56, is the corresponding warning that augmentation loses equivariant information. The negative Weingarten coefficient is a complex integration fact, not a negative rank of a locally free sheaf.

\paragraph{R-04. Plethysm as an operation on representation $K_0$.}
The cycle-index formula expresses $\sigma^d(\sigma^n[E])$ using the Adams operations $\psi^j$. D-S6, V~7.1 and V~7.5, makes their algebraic status precise. They are ring operations; $g\mapsto g^j$ need not be a group homomorphism, and $\psi^j$ need not preserve the cone of honest representations. The $U(2)$ mixed-moment counterexample therefore translates to the statement that separate character distributions do not determine the joint operation. The wreath-subgroup invariants give the characteristic-zero tensor-functor realization. Cote~133, scan~028, lines~12--17 and~26--35, places tensor powers and induction products, with each $\rho_i$ in its own homogeneous degree.

\paragraph{R-05. Image modules, not merely stabilizer invariants.}
The intrinsic object is the image of the degree-$d$ orbit-coefficient map. Its dual is the span of the orbit of $f^{\otimes d}$. Over $\mathbb C$, complete reducibility identifies the multiplicity in each Schur summand with the rank of its multiplicity matrix. The Haar covariance and its reduced Gram matrices have exactly those images; their principal exterior coefficients are squared Pl\"ucker norms. This transports the entire rank proof to determinantal rank strata and Grassmannians in D-FGA. The Hermitian form is an auxiliary real structure and is not asserted to define a motive. Cote~133, scan~028's question about a direct summand is answered here only after passing to characteristic zero.

\paragraph{R-06. Integral matrices and special fibers.}
The polynomial orbit map is defined over $\mathbb Z$ after the specified symmetrizer scalings. Its finite coefficient matrix is an integral map of free modules. The Gaussian covariance is another complex Gram form on its image, while the integral character-sum matrix computes the required Schur-isotypic rank after scalar extension. The Schur projectors require the stated factorial denominators. G-04 specifies the open localization on which the ranks and kernels specialize correctly. D-FGA records the associated rank and Quot strata; D-S6 records the tensor and Schur operations. This translation preserves the exponential matrix sizes, so changing language does not remove the computational gate.

\paragraph{R-07. Binary-form restriction and highest-weight sections.}
The determinant-pencil substitution is a morphism of coefficient schemes, defined after clearing the displayed binomial denominators. The sections $Q_j$ are highest-weight polynomial functions, and independence is proved by their restrictions to the chosen diagonal-pencil chart. Over $\mathbb C$ every binary form is a product of linear forms, so the binary coefficient subscheme lies in the determinant orbit closure. Highest-weight sections of length at most two therefore cannot belong to its ideal. This is the same proof expressed as ideal restriction in D-FGA and weights in D-S3. Cote~133, scan~028 supplies the polynomial-functor viewpoint; it does not supply a new multiplicity formula by rank alone.

\paragraph{R-08. The quotient class in equivariant $K_0$.}
For the polynomial ring $S=\operatorname{Sym}(W)$ and orbit quotient $A=S/I$, an equivariant free resolution yields $[A]=\sum_i(-1)^i[S\otimes T_i]$ in the appropriate graded Grothendieck group. Symmetric powers produce the Molien denominator and the alternating resolution produces the numerator. Exactness, not pointwise sign, proves the formula. D-S6, V~2.2(b), and D-FGA place the additivity and coherent quotient. Cote~133, scan~019, lines~9--20, contains the exterior/symmetric generating-series relation. Ignoring the numerator forgets the defining ideal and is not an operation allowed by $K_0$ additivity.

\paragraph{R-09. A four-factor invariant and a finite-rank cutoff.}
The squared rectangular multiplicity sum is the determinant-isotypic component of $\operatorname{Sym}(E_1\otimes E_2\otimes E_3\otimes E_4)$ in degree $nd$, obtained by composing two tensor restriction functors. The $U(N)$ cutoff means that Schur functors with more than $N$ rows vanish on a rank-$N$ bundle. The stable ratios denominator is therefore an identity of symmetric operations on four tensor factors. D-S3 and D-S6 carry the proof without an arithmetic purity hypothesis. Cote~133, scan~028, is the corresponding tensor-power locus. Nonnegative character norms are exact, but no effective algorithm for the remaining coefficient extraction has been added by the translation.

\paragraph{R-10. The quadratic case and the component group.}
Over $\mathbb C$, the $2\times2$ determinant is a nondegenerate quadratic form in four variables. Its orbit closure is the whole quadratic-form coefficient space. The connected group is the tensor image of $\mathrm{SL}_2\times\mathrm{SL}_2$; the full quadratic stabilizer includes the other orthogonal component. Accordingly the all-odd four-row part present in the ordinary rectangular Kronecker coefficient is removed from the symmetric coefficient and the coordinate ring. D-S3 records the group and component distinction, while D-S6 identifies the even-part Schur decomposition. Cote~133, scan~019's augmentation does not retain this distinction.

\paragraph{R-11. Low-degree surfaces and restricted ideals.}
The linear and quadratic determinant constructions are coefficient-scheme morphisms. The cubic statement uses a dense open set of smooth cubic surfaces and Beauville's determinantal representation theorem; closedness extends the containment. Block padding is a morphism over the coefficient ring, and the highest-weight ideal inclusion is the same D-S3/D-FGA argument as G-08. It excludes four-row multiplicity obstructions for padded forms of degree at most three without asserting a global orbit containment for arbitrary many-variable cubics. Cote~133, scan~028 provides the Schur-functor interpretation of the resulting coordinate modules.

\paragraph{R-12. A Jacobian certificate and its characteristic-zero consequence.}
The permanent-pencil coefficient map is a morphism of affine schemes over $\mathbb Z$. A maximal Jacobian minor nonzero modulo the stated prime is a nonzero integer and therefore nonzero over $\mathbb Q$. In characteristic zero the differential-rank criterion gives dominance, so the four-variable projection of the permanent orbit closure fills the coefficient space in the certified sizes. D-FGA then translates the hypersurface quotient sequence of G-07 into the exact four-row multiplicity difference. The finite-field calculation is a certificate for an integer minor; it is not an appeal to equidistribution or an assertion of semisimplicity modulo that prime. The tensor-power language is cote~133, scan~028; the rank/character distinction remains scan~019.

\subsection{Isolation as a torus action and a problem of coefficient modules}

\paragraph{I-01. Independent variable coordinates.}
The vertex-variable path polynomial is a section of the structure sheaf of an affine toric coordinate space. Its path monomials are distinct characters of its split torus. A Hoza weight is a cocharacter, and uniqueness of the minimum produces a one-dimensional initial character with a nonzero coefficient. D-S3 identifies the action and grading; I-10 supplies its flat hypersurface family over a field. Cote~133, scan~019, is only a finite-set analogy here, not a claim that path existence already computes a representation multiplicity.

\paragraph{I-02. Descent through a monomial quotient.}
The monomial map is a map of affine toric schemes whose character-lattice map is the integer matrix $A$. A cocharacter is induced by an integral weight on the full variable torus exactly when its degree vector lies in the image of $A^{\mathsf T}$. Homogeneity of the binomial kernel gives the grading on the image algebra; extending it to the full variable torus additionally requires the integral image-lattice condition in I-02. For example $z\mapsto x^2$ has zero kernel, but weight one on $z$ is not induced by an integer weight on $x$. The diamond is an explicit failure even of kernel homogeneity. This is D-S3, I~4.7.3, applied to the quotient algebra. Cote~133, scan~019 supplies no rule permitting independent path gradings to be identified with the quotient grading.

\paragraph{I-03. A complete coefficient flag.}
Filter the coefficient module by the span of monomials below each cocharacter weight. Basis isolation says that at each growth degree precisely one indexed coefficient survives, and that these surviving degrees are distinct. The dual formulation quantifies over every nonzero quotient functional. This is a statement about filtered locally free modules over a field and their graded pieces, not just the presence of one initial path. D-S3 gives the associated grading, and D-FGA places the open minor conditions on the relevant flag and Grassmann schemes. Cote~133, scan~028's direct-summand question motivates the distinction between a filtration and a chosen splitting.

\paragraph{I-04. The Pl\"ucker weight polytope.}
Take the coefficient row plane in a Grassmannian and its nonzero Pl\"ucker coordinates. Their torus weights are the sums of the corresponding monomial exponents. The vertex-fiber assertion is the symmetric-exchange proof stated in I-04; the probabilistic estimate counts walls crossed in one coordinate. The resulting open cones are normal cones refined by basis-weight inequalities. D-S3 identifies their cocharacters and D-FGA, Expos\'e~221, Proposition~3.8 places the Grassmannian. No finite-field purity theorem chooses an integer point in all these cones automatically.

\paragraph{I-05. Real Gram positivity and algebraic minors.}
With two independent dual coefficient matrices, Cauchy--Binet is an integral identity of sections of a determinant line. On the real locus where one is the conjugate transpose of the other, it becomes a sum of nonnegative squared Pl\"ucker coordinates. The counterexample shows that this determinant line does not retain the whole coefficient-image flag after colliding weights. D-FGA supplies the rank-stratum language, and D-S6 supplies the determinant as top exterior power. Cote~133, scans~019 and~028, respectively mark the distinction between numerical positivity and actual tensor-module information.

\paragraph{I-06. Secant rank of an evaluation functional.}
An evaluation at a point is a rank-one decomposable functional on the bounded-degree polynomial space. The moment matrix of the proposed nonlinear exponent functional has tensor-product rank. Expressing it as a fixed sum of evaluations would place it in a secant locus of bounded rank; the computed determinant excludes that locus below the stated exponential size. This is an algebraic rank statement over $K(t)$, hence belongs to the determinantal geometry in D-FGA. D-S3 also explains why a nonlinear exponent function is not a character-lattice homomorphism. No motive is needed for this counterexample.

\paragraph{I-07. Taking invariants changes the inserted data.}
The invariant functor on complex representations is realized by the Reynolds projector. An independent twirl sends each state insertion to its invariant scalar component. The two-insertion calculation resolves the commutant into the identity and transposition operations and exhibits a negative coefficient. Thus invariants are extracted in the representation category of D-S3/D-S6, with the full tensor operation retained. Cote~133, scan~028, describes why the symmetric-group action on tensor powers matters; scan~019 warns against replacing that action by dimension.

\paragraph{I-08. A finite quotient morphism, conditional on a hitting set.}
The trace-generating functions are regular invariant functions on the affine scheme of matrix tuples. A hitting set separating their coefficient polynomials gives a separating family in the quotient coordinate algebra; the cited invariant-theory theorem makes that algebra finite over the generated subalgebra. This is a finite morphism of affine schemes, not automatically a choice of algebraically independent parameters of minimal cardinality. D-FGA gives the general morphism and parameter-space setting. The unresolved explicitness condition is exactly the one retained in I-08 and GCT~V; the dictionary does not discharge it.

\paragraph{I-09. A tensor contraction computing a Hermitian norm.}
The matrix product computes a regular coefficient map. Tensoring with its conjugate and contracting equal monomial indices gives its squared norm; the torus integral extracts the zero character in that tensor product. D-S3 explains the character extraction. Over an integral algebraic base one must use independent dual variables, not treat complex conjugation as a regular morphism. The metric realization lives over $\mathbb R$ or $\mathbb C$. It is the same coefficient module as before, and access to its presenting matrices remains a white-box assumption.

\paragraph{I-10. Rees modules and specialization.}
The cocharacter gives an increasing filtration on coefficient spans. Its Rees module has the associated graded as its special fiber and the original module on the punctured line. The explicit hypersurface homogenization in I-10 is flat over a field by the torsion-free PID argument in its proof; general ideal degenerations require the entire weight-homogenized ideal. D-S3 identifies the grading and D-FGA identifies the flat family and its Hilbert/Quot points. The associated graded preserves total filtered rank, whereas adding colliding coefficients in a substituted polynomial can reduce their span, as I-05 proves. These are different maps.

\paragraph{I-11. Exterior powers do not remove equations.}
The exterior dimension and graph-of-linear-map count describe the sizes of coordinate and Grassmannian models. The two-term zero determinant is already an equation of the rank-deficient locus. Assigning different path weights to its cancelling terms changes that equation. D-S6 gives exterior powers and D-FGA gives the Grassmannian, but neither turns a signed determinant expansion into unsigned reachability. Cote~133, scan~028's direct-summand margin remains a relevant distinction between a functorial module and its termwise expansion.

\paragraph{I-12. Support is the nonvanishing locus of a regular function.}
A specified coefficient of an integer ROABP is a regular polynomial in its edge parameters. Its nonvanishing open subset is determined by a signed path sum, whereas ordinary reachability detects whether a parameterizing path set is empty. The GapL reduction is a computation of that regular function; the conditional complexity-class inclusion follows only under the specified logarithmic-space reduction. D-FGA locates its zero scheme. Cote~133, scan~019 supplies the nearby conceptual distinction between the cardinality of a finite set and a signed class; it states no complexity theorem.

\paragraph{I-13. Evaluation quadrature and the Segre variety.}
Coefficient tensors of width-one products form the affine cone over a Segre variety. Equality of the quadratic forms on all these tensors, with their complex conjugates, forces the moment-matrix identity proved in I-13. Its rank then bounds the number of evaluation vectors. This restates the entire proof as an assertion about a Segre cone, its dual evaluation vectors, and determinantal rank loci in D-FGA. The Hermitian positivity is not an integral $\lambda$-ring assertion. The much smaller hitting set in the same finding concerns only kernel avoidance, not equality of quadratic forms.

\subsection{Purity, arithmetic sampling, and the limits of numerical shadows}

\paragraph{P-01. A Gram form on a finite space of sections.}
The coefficient support spans a vector space of functions. Their pairwise correlations are entries of a Gram matrix; its spectral lower bound is the elementary estimate in P-01. In geometric applications these functions may be sections or Frobenius traces, but no sheaf is manufactured by naming the Gram matrix. D-ET applies only after an actual trace-function realization exists. The distinction from rank augmentation in cote~133, scan~019, is that all cross terms have been retained and bounded.

\paragraph{P-02. A Kummer-character construction.}
The signs arise from the rank-one quadratic Kummer sheaf pulled back by $t\mapsto t+i$ on the line with its zeros removed. Products over distinct indices remain governed by the explicitly squarefree polynomial, so the curve character estimate controls the off-diagonal trace sums. The altered values at the excluded points contribute the stated elementary error. D-ET, together with the exact Alon--Goldreich--H\aa stad--Peralta locator, carries the proof. This yields the stated sparse multilinear hitting set; it does not create a polynomial bound on arbitrary coefficient supports.

\paragraph{P-03. The kernel of restriction.}
The sign-cube counterexample is a nonzero section whose restriction vanishes on the selected finite subscheme; the common-affine repair is annihilated by the explicit Vandermonde section. These are regular-function identities and survive independently of improved point-counting estimates. They belong to the affine-scheme language of D-FGA. Cote~133, scan~019 is an analogy for information lost by passing to a smaller numerical image, not a source of these new counterexamples.

\paragraph{P-04. Pullback to a curve and its zero divisor.}
For the curve morphism specified by the rational functions $r_i$, the section $f(r_1,\ldots,r_n)$ has polar divisor bounded by the displayed $D$. If it is nonzero, its zeros have total degree at most $\deg D$. The trace formula and curve RH then provide enough rational points. D-ET controls the count; the pullback's nonvanishing is a separate scheme-theoretic condition. This proof uses curves and sheaves actually specified in the statement, so a motivic interpretation is available for their cohomology, not for an unspecified ROABP.

\paragraph{P-05. A curve family can lie in a small polynomial kernel.}
Dimension comparison in the spaces $H^0(C,\mathcal O_C(aD))$ produces an algebraic relation between the two chosen functions. Multiplying the relations for the finite curve family gives its annihilating section and the explicit bivariate ROABP. D-FGA places the image and its ideal; D-ET cannot remove a function from that ideal by increasing the field size. The relation is proved directly with its degree and width bounds and is not a lower bound against every geometric sampler.

\paragraph{P-06. Purity with a complete conductor budget.}
The relevant class is $[\mathcal F_i\otimes\mathcal F_j^\vee]$ in the Grothendieck group of constructible Weil sheaves. The compactly supported trace gives its Gram entry. Geometric irreducibility and inequivalence determine $H_c^2$, affine vanishing removes $H_c^0$, and Weil~II bounds $H_c^1$. SGA~5 supplies the displayed Swan-conductor expression for its dimension. This is exactly D-ET, with every rank and conductor retained. A claim of polynomial explicitness additionally requires that these data and the enumeration of the points be polynomially bounded. Cote~133, scan~028 supplies a tensor-category analogy, but not the missing sheaf realization.

\paragraph{P-07. Sparse cocharacters over the integer lattice.}
The weight family is an explicit finite set in $\operatorname{Hom}(\mathbb G_m,\mathbb G_m^n)=\mathbb Z^n$. The collision polynomial over $\mathbb F_p$ is nonzero by its degree and coefficient bounds, so a seed outside the union of its root sets makes the support characters distinct. D-S3 translates the construction exactly; only the elementary finite-field root bound is used. Thus the result is independent of an additional purity or RH hypothesis.

\paragraph{P-08. Arithmetic covers with a specified nonvanishing encoding.}
The common number field determines a finite \emph{\'etale} cover of $\operatorname{Spec}\mathbb Z$ after inverting its discriminant. Frobenius conjugacy classes are the arithmetic data on that cover. Effective Chebotarev produces a small unramified prime only after the stated common map from primes to evaluations and its nonvanishing property are supplied. D-ET provides the parallel finite-field trace perspective; it is not a proof of that arithmetic encoding. The abelian case uses the stipulated Dirichlet GRH, while the general stated Hecke extension remains conditional on Masullo. No claim about arbitrary automorphic $\mathrm{GL}_n$ GRH is made.

\paragraph{P-Z01 and P-Z02. Realizations of zero statistics.}
These are analytic consequences of the precisely stated RH predicates and supported explicit-formula theorems. Their expression by unitary integrals does not make the zeta zeros the eigenvalues of a constructed geometric Frobenius. D-ET explains the function-field comparison only in the families with an established sheaf and monodromy group. There is no cote~133 scan containing a classical-RH-to-GCT implication; scans~019 and~028 supply representation operations, not that missing implication.

\subsection{The integral controls G-01--G-09}

\paragraph{G-01. Burnside classes, representation classes, and rank.}
The Burnside-ring counterexample is an equivariant finite-set computation. Linearization sends a $G$-set to a permutation module, but the two-subset operation does not become exterior power: the $C_2$ sign in the proof witnesses this failure. In D-S6 the correct relation is $x^2=2\lambda^2(x)+\psi^2(x)$, and only augmentation replaces $\psi^2(x)$ by the rank. This is the precise mathematical correction to cote~133, scan~019, line~53 and the supplied study's equivariant reading.

\paragraph{G-02. The universal terminal variable.}
Specializing one factor to the tensor unit in the universal product identity forces the top elementary symmetric function to remain present. SGA~6, V~2.4, equation~(2.4.2), explicitly includes that terminal variable. The proof in G-02 is its necessity in the universal character ring, so the correction to scan~006, lines~25--26, is mathematical and independent of paleography.

\paragraph{G-03. Equal finite filtration data, unequal representations.}
The two sums of line characters are honest locally free representations of $\mathbb G_{m,\mathbb Z}$, whose difference is the $(D+1)$st power of an augmentation-zero element. D-S6, V~3.10, places the difference in the required filtration step. The associated-bundle construction transports this identity to any scheme carrying the line bundle; it does not force the bundles to remain distinct on every base. Thus finite Chern or $\gamma$ data cannot recover all the character coefficients needed in scans~019 and~028.

\paragraph{G-04. The integral image and its good rank locus.}
The schematic orbit closure is defined over $\mathbb Z$ by a kernel, and its image algebra is torsion-free. In a fixed degree a nonzero maximal minor gives an explicit principal open on which the coefficient matrix is equivalent to a split rank matrix. This is a direct proof of base change for the displayed module, with D-FGA providing the image, Grassmannian and Quot framework. At the excluded primes the flat integral model's fiber may differ from the orbit closure of the reduced point. This is why integral $K_0$ identities from scan~028 do not imply unchanged characteristic-zero multiplicities in all fibers.

\paragraph{G-05. Plane curves and the three-row sector.}
The nonempty complement of the theta divisor in $\operatorname{Pic}^{g-1}(C)$ supplies Beauville's line bundle, hence a linear determinantal equation. Density of smooth curves gives the scheme-theoretic containment needed for every ternary coefficient section. Highest-weight sections of length at most three belong to that coefficient subalgebra, so its determinant ideal has none. D-FGA records the ideal restriction and parameter spaces; D-S3 records the weights. The numerical positive count is not extracted from Chern classes of the curve alone.

\paragraph{G-06. Stabilizer dimensions and four-row equations.}
The polynomial matrix-pencil map is invariant under the explicitly stated subgroup scheme $H$ of $\mathrm{GL}_n^2$. Its generic stabilizer is the scalar subgroup by the concrete matrix calculation. D-S3, VI$_B$~4.1 and its cited EGA theorem justify the dimension behavior, and the resulting proper image has a nonzero ideal. D-FGA and highest-weight restriction give an equation of exactly four rows, using G-05 to exclude shorter weights. This supplies no nonvanishing on a particular padded form until that extra statement is proved.

\paragraph{G-07. A Cartier divisor and a determinant twist.}
The reduced determinantal-quartic locus is a prime hypersurface. Its equation spans the character $\det^D$, so its ideal is the line-character twist of the ambient graded algebra shifted by $D$. The exact sequence in G-07 is the equivariant $K_0$ statement of a principal Cartier divisor. D-S6 gives additivity, D-FGA gives the projective hypersurface and quotient, and D-S3 identifies the character. It produces the shifted plethysm subtraction exactly; the corrected reduced degree is a literature issue, not a new transcription of cote~133.

\paragraph{G-08. A universal morphism of coefficient schemes.}
The two monomial trees and their cofactor determinant give a matrix over the universal integer coefficient ring. Homogenization and block padding therefore commute with any base change. Over $\mathbb C$ the containment of all low-variable restrictions yields an inclusion of highest-weight ideal spaces, hence the multiplicity inequality. This is D-FGA applied to an explicit scheme morphism, with the weights read through D-S3 and the symmetric powers through D-S6. Its complexity bound is part of the construction, not a consequence of representability alone.

\paragraph{G-09. A local padding coordinate and smaller universal matrices.}
Once the nonzero padding form is one coordinate, dehomogenization reduces the number of affine variables by one. The shared exponent chains are again defined over the universal coefficient ring; their exact vertex count gives the stated finite size. Changing an arbitrary padding form to that coordinate is local on its nonzero coefficient charts. D-FGA carries these local morphisms and their closure statement, and the D-S3 highest-weight argument yields the improved row exclusion. The known asymptotic uniform-representation theorem is cited separately; neither the chart construction nor its GCT consequence requires RH.

\subsection{What the manuscript record does and does not establish}
The supplied cote~133 study also reports decisions about the name on scan~008, a characteristic-two pre-$\lambda$ issue on scan~006, and a crossed-out Stiefel--Whitney assertion on scan~009. Those reports are retained as reports of earlier readers, with their own scope. No new visual decision about the name or the deleted text is made in this run. The five attached LaTeX transcriptions contain no original raster pages. The separate transcription table records the forced algebraic readings and the distinction between a textual ambiguity, a mathematical correction, and a claim whose ink cannot be checked here.

\part{Unchecked targets and unresolved gates}
\section{Unchecked targets and the current smaller gates}

This part contains no asserted proof of the unresolved targets. A conditional reduction in the checked part means that the implication was proved; it does not mean its antecedent was established. The live question register and dated ledger give the current states, including additional deductions made after this text was first assembled.

\paragraph{U-01: general coefficient extraction and determinant-boundary ranks.}
The exact rectangular Haar transforms, ratios kernel, positive finite-group projectors and Gram-rank identities are established in R-01--R-09. The remaining general gate is a uniformly efficient positive/combinatorial realization of the relevant multiplicity spaces, together with effective control of the orbit-map ranks or quotient numerator. Low-row restrictions have been removed from possible obstruction searches in the proved ranges G-05 and G-08--G-09, while G-07 and R-12 solve a specific four-row quotient calculation exactly. The general nonstandard quantum-group conjectures and their algorithmic requirements are listed with source locators in the checked dependency audit.

\paragraph{U-02: a polynomial ROABP hitting set.}
The strongest current gate is explicit joint coefficient-span preservation in I-14, or the sufficient stronger cone-hitting condition in I-03--I-04. Direct path weights fail to respect the quotient; exact nonlinear exponent simulation and exact/constant-distortion energy cubature have the construction-specific lower bounds recorded below. The affine seed-decoder and singular-factor examples test further repairs. A candidate that regularizes the matrices must preserve a unit Pl\"ucker coordinate on the original special fiber, as I-18 makes precise. Neither Hoza's graph theorem nor the new analytic inputs currently supplies that candidate.

\paragraph{U-03: matrix invariants and GCT~V.}
Once U-02 is solved with the stated field and construction bounds, I-08 gives the polynomial-size separating family and finite matrix-quotient morphism. The implication is complete; the polynomial hitting-set antecedent is not. General explicit Noether normalization and general PIT have stronger hypotheses than this matrix-specific consequence. No circuit lower bound or $P$ versus $NP$ separation is asserted here.

\paragraph{U-04: a geometric or arithmetic encoding with polynomial budgets.}
The sparse and actual-sheaf constructions are proved. For arbitrary ROABPs the remaining task is to construct a sampler whose evaluation kernel avoids the bounded-width tensor variety, without requiring near-isometry of the entire coefficient norm. P-10 makes that kernel criterion exact. P-11 removes a detour: a sufficiently explicit common algebraic point is itself a hitting-set encoding, independently of GRH. A Frobenius construction would have to supply new nonvanishing structure; applying Chebotarev after assuming it does not solve the original problem.

\paragraph{U-05: unrestricted number-field random-matrix laws.}
The dependency table gives the supported statements now discharged under the stipulated inputs. Full pair correlation, spacing laws, the general arithmetic moment recipe, arbitrary automorphic target RH, and a zeta-zero quantum Hamiltonian retain their explicitly named analytic or spectral prerequisites. No literature theorem is strengthened beyond its stated support or uniformity range by changing an RH label.

\paragraph{Next work and stopping discipline.}
The attempts files record the constructions actually tried and their failure locations. The next substantive work is on the remaining joint-compression and special-fiber gates and on the restricted projected-orbit ideals. Packaging and repeated exposition are not counted as movement. Checkpoint and review times are actual elapsed times recorded in the ledger; no scheduled review is claimed before it has occurred.

\appendix
\part{References, exact locators, and audit scope}
\section{Rectangular bridge: sources and exact
locators}\label{rectangular-bridge-sources-and-exact-locators}

All new R statements are deductions made and checked in this run; no
claim of priority over the literature is made. The older numerical
formulas are explicitly credited. Attached LaTeX was read directly, not
inferred from the prompt summary.

\subsection{R-B01 --- Bump--Gamburd rectangular product
identity}\label{r-b01--bumpgamburd-rectangular-product-identity}

Bump, D., \& Gamburd, A. (2006). On the averages of characteristic
polynomials from classical groups. \emph{Communications in Mathematical
Physics, 265}, 227--274. \url{https://doi.org/10.1007/s00220-006-1503-1}

Exact locator: §3.1, Proposition 4, equation source label
\texttt{e:produ}; theorem label \texttt{cpaverage}. Attached source
\texttt{LIT\_Bump\_Gamburd\_averages\_of\_characteristic\_polynomials/acp20.tex},
lines 994--1039. General parameters are \texttt{N,L,K}; the polynomial
is \texttt{s\_(N\^{}L)} in \texttt{K+L} variables. This supports R-01.

\subsection{R-B02 --- Bump--Gamburd swaps and
ratios}\label{r-b02--bumpgamburd-swaps-and-ratios}

Same work. Exact locators: §2.3, Lemma 1 (\texttt{lem:min1}), §3.1,
Proposition 4 (\texttt{cpaverage}), and §4 (\texttt{comf}) for
rectangular Weyl coset/Laplace sums. §3.3, Theorem 3 (\texttt{t:1}),
equation \texttt{unitaryratio} for inverse characteristic-polynomial
ratios, with \texttt{N≥Q,R} and the denominator parameters of modulus
less than one. The specialization \texttt{K=L=0,Q=R=n²} used in R-09 is
within those hypotheses when \texttt{N≥n²}. Attached source lines
1539--1563 and 1650--1735.

\subsection{R-B03 --- Finite CUE moments, hook content, Schur
orthogonality}\label{r-b03--finite-cue-moments-hook-content-schur-orthogonality}

Same work. §3.1, Corollary 1 (\texttt{ratioskeatingsnaith}), equation
\texttt{ksformula}; §2.2, equations \texttt{ratiosschurorthog} and
\texttt{e\_hhl}. Finite moment equals the rectangular representation
dimension; the leading asymptotic coefficient need not be an integer.

\subsection{R-B04 --- Arithmetic CFKRS status and swap
formula}\label{r-b04--arithmetic-cfkrs-status-and-swap-formula}

Conrey, J. B., Farmer, D. W., Keating, J. P., Rubinstein, M. O., \&
Snaith, N. C. (2005). Integral moments of L-functions. \emph{Proceedings
of the London Mathematical Society, 91}(1), 33--104.
\url{https://doi.org/10.1112/S0024611504015175}

Version matching attachment: arXiv:math/0206018v3 (2005),
\url{https://arxiv.org/pdf/math/0206018} . Exact locators: Conjecture
1.5.1 (\texttt{thm:zeta2kconjecture}), Theorem 1.5.2
(\texttt{thm:UN2k}), §2.1 recipe, §2.2 equation \texttt{conj:M}, Lemma
2.5.1 (\texttt{thm:concisesumunitary}). The author\textquotesingle s
older PDF at Waterloo is v1 and uses different section numbers; do not
cite its §3.2 as if it matched the attached §2.5. The statements are
checked as conjecture versus theorem, not assumed proved from RH.

\subsection{R-B05 --- Schur--Weyl and exact
projectors}\label{r-b05--schurweyl-and-exact-projectors}

Collins, B., \& Śniady, P. (2006). Integration with respect to the Haar
measure on unitary, orthogonal and symplectic group.
\emph{Communications in Mathematical Physics, 264}(3), 773--795.
\url{https://doi.org/10.1007/s00220-006-1554-3}

Primary source: \url{https://arxiv.org/pdf/math-ph/0402073} . Exact
locators: Theorem 2.1 (Schur--Weyl duality), equation (3) and the
central idempotents preceding equation (4), Proposition 2.2 (Haar
conditional expectation as orthogonal projection).

\subsection{R-B06 --- Weingarten calculus with small-dimension
correction}\label{r-b06--weingarten-calculus-with-small-dimension-correction}

Same work: Proposition 2.3, equation (9), includes the necessary
\texttt{ell(lambda)≤dimension} cutoff; Corollary 2.4, equation (11), is
the entry integration formula. The example after Proposition 2.5 gives
the negative transposition value \texttt{−1/{[}N(N²−1){]}}. No
unrestricted stable Weingarten inverse is substituted at singular small
dimensions.

\subsection{R-B07 --- Rains theorem}\label{r-b07--rains-theorem}

Rains, E. M. (2003). Images of eigenvalue distributions under power
maps. \emph{Probability Theory and Related Fields, 125}(4), 522--538.
\url{https://doi.org/10.1007/s00440-002-0250-2}

Primary source: \url{https://arxiv.org/pdf/math/0008079} . Exact
locators: Theorem 1.3, equation (8), and Corollary 1.4. They describe
one power\textquotesingle s eigenvalue law. They do not assert
independence of different powers of the same matrix.

\subsection{R-B08 --- GCT multiplicity criterion and rectangular
obstruction
limit}\label{r-b08--gct-multiplicity-criterion-and-rectangular-obstruction-limit}

Ikenmeyer, C., \& Panova, G. (2017). Rectangular Kronecker coefficients
and plethysms in geometric complexity theory. \emph{Advances in
Mathematics, 319}, 40--66.
\url{https://doi.org/10.1016/j.aim.2017.08.024}

Primary source: \url{https://arxiv.org/pdf/1512.03798} (v2); exact
locators: Theorem 1.4 (\texttt{thm:strategy}), equations (1.3) and
(1.5), Theorem 1.6 (\texttt{thm:together}), and Theorem 2.1
(\texttt{thm:manivel\_rect}). Theorem 1.6 requires
\texttt{n\textgreater{}3m\^{}4}. Source hard-coded cross-references
occasionally refer to earlier numbering, so numerical locators were
confirmed against the rendered primary PDF.

\subsection{R-B09 --- Orbit versus closure and no occurrence
obstructions}\label{r-b09--orbit-versus-closure-and-no-occurrence-obstructions}

Bürgisser, P., Ikenmeyer, C., \& Panova, G. (2019). No occurrence
obstructions in geometric complexity theory. \emph{Journal of the
American Mathematical Society, 32}(1), 163--193.
\url{https://doi.org/10.1090/jams/908}

Primary source matching attachment:
\url{https://arxiv.org/pdf/1604.06431v3} . Exact locators: Theorem 1.4
(\texttt{th:main}), \texttt{n≥m\^{}25}; §1(b), paragraphs defining the
GCT coefficients and bounding them by the symmetric rectangular
Kronecker coefficient of the orbit, then by the ordinary rectangular
coefficient. The old v1 theorem number 1.5 must not be silently used for
the supplied v3 source.

\subsection{R-B10 --- Known stable numerical
positivity}\label{r-b10--known-stable-numerical-positivity}

Manivel, L. (2011). On rectangular Kronecker coefficients. \emph{Journal
of Algebraic Combinatorics, 33}, 153--162.
\url{https://arxiv.org/abs/0907.3351}

Primary PDF: \url{https://arxiv.org/pdf/0907.3351} . Exact locators:
Theorem 1(2); §3.3, Corollary 2(1), which explicitly states that for
\texttt{d≥q}, \texttt{g((nd−q,q),(d\^{}n),(d\^{}n))} is the number of
partitions of \texttt{q} into parts between 2 and n. §3.4 gives a
birational flag-quotient interpretation. R-07 credits this numerical
identity and proves explicit centered-polynomial witnesses and their
orbit-closure survival, rather than claiming a new count.

\subsection{R-B11 --- GCT IV: proved two-row rule, its exact rectangular
specialization, and
limits}\label{r-b11--gct-iv-proved-two-row-rule-its-exact-rectangular-specialization-and-limits}

Blasiak, J., Mulmuley, K. D., \& Sohoni, M. (2015). \emph{Geometric
complexity theory IV: Nonstandard quantum group for the Kronecker
problem}. Memoirs of the American Mathematical Society, 235(1109).
\url{https://doi.org/10.1090/memo/1109}

The supplied file is the later three-author
\texttt{Apr13KroneckerGCT4.tex}, not the original two-author March 2007
version cited by the older GCT VII/VIII manuscripts. Matching author
PDF: \url{https://websites.umich.edu/~jblasiak/Apr13KroneckerGCT4.pdf} .
Page and theorem numbers below are for that Apr13 manuscript.

\begin{itemize}
\tightlist
\item
  Exact attached file:
  \texttt{sources/LIT\_Mulmuley\_GCT\_IV\_nonstandard\_quantum\_group\_Kronecker/Apr13KroneckerGCT4.tex}.
\item
  §12.1, theorem source label
  \texttt{t\ nonstandard\ schur-weyl\ duality}, source lines 4236--4290:
  nonstandard Schur--Weyl duality proved for general dimensions. This
  established structural theorem should be distinguished from the later
  two-row limitation on explicit canonical bases.
\item
  Theorem 15.21, especially (vii)--(viii), manuscript p.102, label
  \texttt{t\ main\ canonical\ basis}, source lines 5963--5978:
  highest-weight counting rule for two-row Kronecker coefficients and
  its symmetric/exterior refinement. The two-row restriction means both
  \emph{inner} GL partitions have length at most two; the outer
  partition may have four rows.
\item
  Theorem 18.2, manuscript pp.124--125, label
  \texttt{t\ consequence\ for\ Kronecker}; equation (163), label
  \texttt{e\ g\ =\ invariant-free\ g}: constructive positive two-row
  rule from five types of invariant-free components, each zero or one.
\item
  Corollary 18.4, manuscript p.127, label
  \texttt{c\ kronecker\ generating\ function\ sym\ ext}, source lines
  7736 onward; introduction equivalent equation
  \texttt{e\ intro\ Kronecker\ generating\ function}, lines 1435--1455:
  symmetric/exterior generating functions. Taking the constant
  coefficient gives the all-even/four-odd-part formula recorded as R-10.
  The numerical formula is known; our work specializes and compares it
  with the closure ring.
\item
  Example 17.8, manuscript p.122, label \texttt{ex\ not\ positive},
  source lines 7326--7351: a divided-power matrix coefficient has mixed
  Laurent signs, namely \texttt{q\^{}10-q\^{}4-q\^{}-4+q\^{}-10}. The
  source first states that no global sign readjustment makes every
  specified Chevalley matrix coefficient nonnegative. The displayed
  algebraic identity was separately recomputed in
  \texttt{check\_n2\_formulas.py}.
\item
  Conjecture 19.1, §19.1, label \texttt{cj\ canonical\ basis\ X\^{}r},
  lines 7920--7967: proposed compatible basis for the full tensor power.
  §19.2, source lines 8191--8209: naive extension can have too small a
  dimension; it has not supplied the general positive rule. Intro §1.1,
  source lines 893--903, explicitly describes the unresolved cases.
\item
  R-10 gives a direct principal-minor highest-weight proof of the
  classical identity
  \texttt{Sym\^{}d(Sym\^{}2\ C\^{}s)=direct\_sum\ S\_(2mu)\ C\^{}s}, so
  no unverified book exercise number is used for that identity.
\end{itemize}

\subsection{R-B12 --- GCT VII: proved structure versus conjectural
plethysm
reciprocity}\label{r-b12--gct-vii-proved-structure-versus-conjectural-plethysm-reciprocity}

Mulmuley, K. D. (2008a). \emph{Geometric complexity theory VII:
Nonstandard quantum group for the plethysm problem} (revised extended
abstract; arXiv:0709.0749v2; original Technical Report TR-2007-14,
2007). \url{https://arxiv.org/abs/0709.0749v2}

\begin{itemize}
\tightlist
\item
  Exact attached file:
  \texttt{sources/LIT\_Mulmuley\_GCT\_VII\_nonstandard\_quantum\_group\_plethysm/main.tex}.
\item
  Introduction, source lines 212--229, Theorem 1.2, label
  \texttt{tnonstd}: existence of the nonstandard quantization, complete
  reducibility, and quantum Peter--Weyl theorem.
\item
  Introduction, lines 257--285, Conjecture 1.3, label \texttt{cduality}:
  nonstandard duality. Lines 290--305, Conjecture 1.4, label
  \texttt{creciprocity}: nonstandard reciprocity. §6, Conjecture 6.1,
  label \texttt{crecirefined}: refined specialization statement.
\item
  Introduction, lines 309--335: conditional positive decomposition into
  \texttt{m\_lambda\^{}alpha\ n\_pi\^{}alpha} and the two subproblems
  P1/P2. Lines 336--366 explicitly state that algorithm correctness and
  the required canonical-basis properties are not proved in this source.
\item
  Abstract and introduction distinguish this proposed positive plethysm
  formula from already established existence and complete reducibility.
  Latest attached arXiv version is v2, September 1, 2008; the first
  version is September 5, 2007.
\end{itemize}

\subsection{R-B13 --- GCT VIII: precise conjectural positivity and
finite-field RH
language}\label{r-b13--gct-viii-precise-conjectural-positivity-and-finite-field-rh-language}

Mulmuley, K. D. (2008b). \emph{Geometric complexity theory VIII: On
canonical bases for the nonstandard quantum groups} (revised extended
abstract; arXiv:0709.0751v2; original Technical Report TR-2007-15,
2007). \url{https://arxiv.org/abs/0709.0751v2}

\begin{itemize}
\tightlist
\item
  Exact attached file:
  \texttt{sources/LIT\_Mulmuley\_GCT\_VIII\_canonical\_bases\_nonstandard\_quantum\_groups/main.tex}.
  Primary HTML for this version:
  \url{https://arxiv.org/html/0709.0751v2} .
\item
  Abstract, lines 144--159, and Introduction, lines 195--218: algorithms
  are explicitly called conjecturally correct; required properties are
  also conjectural.
\item
  §2.2, Conjecture 2.2, label \texttt{clocalcrystalb}, lines 634--655:
  local crystal basis existence.
\item
  §2.3, Conjecture 2.5, label \texttt{cuniquemindegleft}, lines
  947--952: unique minimum degree complexity required for globalization.
\item
  §2.4.1, Conjecture 2.6, label \texttt{ccelldecompleft}, lines
  1041--1046: cellular decomposition compatible with the restricted
  standard quantum group.
\item
  §2.4.2, Conjecture 2.7, label \texttt{cposkroneckerleft}, lines
  1073--1090: Kronecker positivity, with the precisely stated sign and
  \texttt{(q-q\^{}-1)\^{}a} factors allowed for structure constants. The
  following paragraph defines positive coefficients to mean nonnegative
  \textbf{rational} coefficients.
\item
  §2.4.3, Conjectures 2.8--2.9, labels \texttt{cnonstdposleft0} and
  \texttt{cnonstdposleft1}, lines 1151--1214: weaker general saturation
  and nonstandard almost-positivity claims; not an established
  nonnegative integer tableau count.
\item
  §3.3, Conjectures 3.1--3.2, labels \texttt{cunimindegright},
  \texttt{cunimindegright2}; §3.4.1, Conjecture 3.3, label
  \texttt{ccelldecompright}; §3.5, Conjecture 3.4, label
  \texttt{cposkroneckerright}; §3.5.1, Conjectures 3.5--3.6, labels
  \texttt{cnonstdposright}, \texttt{cnonstdposright1}; §3.5.2,
  Conjecture 3.7, label \texttt{coppquasicell}.
\item
  §6, label \texttt{scomplexity}, lines 2810--2854: efficient
  combinatorial labels and crystal operations needed for a \#P rule are
  additional proposed properties, with details deferred to a full
  version.
\item
  Introduction, lines 238--247, and §2.4.2, lines 1118--1126:
  perverse-sheaf theory and the finite-field Riemann-hypothesis
  machinery are invoked as a prospective extension needed for nonflat,
  noncommutative varieties. These passages supply no implication from
  the classical analytic RH or Dirichlet GRH to the stated algebraic
  conjectures.
\item
  Exact theorem/conjecture numbers were checked against the shared TeX
  theorem counter; principal positivity and complexity locators were
  also matched to the primary HTML. No full independent proof of these
  background quantum-group constructions is claimed.
\end{itemize}

\subsection{R-B14 --- Cubic surfaces and the small-degree four-row
exclusion}\label{r-b14--cubic-surfaces-and-the-small-degree-four-row-exclusion}

Beauville, A. (2000). Determinantal hypersurfaces. \emph{Michigan
Mathematical Journal, 48}, 39--64.
\url{https://math.univ-cotedazur.fr/u/beauvill/pubs/det.pdf}

\begin{itemize}
\tightlist
\item
  Corollary 6.4, printed pp.52--53 (PDF pages 14--15): over an
  \textbf{algebraically closed} field, every \textbf{smooth cubic
  surface} admits a \texttt{3×3} linear determinantal equation. Its
  additional count of 72 representations is not needed for R-11.
\item
  Corollary 6.5, printed p.53, states the extra
  twisted-cubic/birationality/six-lines conditions over a field that
  need not be algebraically closed. R-11 works over C and does not
  assert the unrestricted statement over other fields.
\item
  §1.9, printed p.42, cross-references the plane-curve and cubic-surface
  statements and explains the dimensional boundary for generic
  hypersurfaces.
\item
  R-11 uses Corollary 6.4 only for smooth quaternary cubics, then proves
  density, block-diagonal padding, linear-substitution closure, and
  highest-weight ideal containment. It does not claim that every
  singular cubic itself is a linear determinant; membership in the
  closure is sufficient.
\end{itemize}

\subsection{Web retrieval references retained for root citation
verification}\label{web-retrieval-references-retained-for-root-citation-verification}

The root must reopen a subagent\textquotesingle s web source before
using a web citation. The source URLs above are sufficient; retrieval
ids are recorded only for internal provenance.

\begin{itemize}
\tightlist
\item
  Rains PDF: \texttt{turn671720168view0}.
\item
  Collins--Śniady PDF: \texttt{turn671720168view1}, targeted formula
  view \texttt{turn506510132view0}.
\item
  Ikenmeyer--Panova PDF: \texttt{turn629217195view0}, stability passage
  \texttt{turn124126028view0}.
\item
  Bürgisser--Ikenmeyer--Panova PDF: \texttt{turn629217195view1}.
\item
  Manivel PDF: \texttt{turn407184572view0}, Corollary 2 passage
  \texttt{turn209001524view1}.
\item
  CFKRS v3 PDF: \texttt{turn209001524view0}.
\end{itemize}

All cited mathematical assertions are either in the supplied sources or
these primary papers. Search snippets from secondary sites were used
only to locate the primary sources.
\section{Isolation sources and exact
locators}\label{isolation-sources-and-exact-locators}

All source paths below refer to the supplied archive extracted once to
\texttt{sources/}. Existing results are not claimed as new. The dates in
bibliographic entries follow publication/version dates, not regenerated
PDF title-page timestamps.

\subsection{I-B01 --- new graph input}\label{i-b01--new-graph-input}

Hoza, W. M. (2026). \emph{NL = UL} (ECCC TR26-323). Supplied file
\texttt{ECCC\_TR26-323\_Hoza\_NL\_equals\_UL.pdf}.

\begin{itemize}
\tightlist
\item
  Definition 1.2, printed p. 2: internal-vertex path weights; uniqueness
  for all reachable vertex pairs.
\item
  Theorem 1.3, printed pp. 2--3: one seed per fixed simple graph; seed
  length O(log(WN)); logarithmic-space computability; polynomial integer
  magnitudes.
\item
  Formula (1), printed p. 3: rho\_x(a,i)=sum\_j 6\^{}(nu\_2(i)+ell-j)
  C(a)\_j x\_j.
\item
  Theorem 2.4, printed p. 5: failure probability at most 2\^{}(-d)
  W\^{}4 N\^{}2.
\item
  Used as the supplied theorem. This track did not claim a formal or
  full independent verification of the theorem.
\end{itemize}

\subsection{I-B02 --- basis isolation}\label{i-b02--basis-isolation}

Agrawal, M., Gurjar, R., Korwar, A., \& Saxena, N. (2014).
\emph{Hitting-sets for ROABP and sum of set-multilinear circuits}
{[}Preprint{]}. arXiv:1406.7535v1. \url{https://arxiv.org/abs/1406.7535}

Supplied source:
\texttt{LIT\_AGKS\_hitting\_sets\_ROABP\_basis\_isolation/deltaDistanceSICOMP.tex}.

\begin{itemize}
\tightlist
\item
  Theorem 1: source label \texttt{thm:ROABPhs}; source lines 241--246;
  PDF p. 2.
\item
  Definition 5: basis isolating weight assignment; source lines
  902--919; PDF p. 7.
\item
  Lemma 6: source label \texttt{lem:basisIsolation}, lines 927--988; PDF
  pp. 7--8.
\item
  Observation 7: \texttt{obs:smallerWeight}, lines 1005--1028; PDF p. 8.
\item
  Lemma 8: \texttt{lem:weightFunction}, lines 1040--1055 and proof
  thereafter; PDF pp. 9--11.
\item
  Additional I-17 context: Theorem 3 (\texttt{thm:invROABPHS}), printed
  pp. 5 and 19; Lemma 27 (\texttt{lem:bSConc}), printed p. 22, source
  lines 2354--2400, invertible-constant-term matrix concentration.
  Primary PDF refs \texttt{turn317946964view0},
  \texttt{turn814896166view0}, \texttt{turn814896166view1},
  \texttt{turn814896166view2}.
\item
  Primary PDF inspected: \url{https://arxiv.org/pdf/1406.7535} ; web
  reference \texttt{turn166538694view0}; targeted findings
  \texttt{turn944129952view0}, \texttt{turn944129952view1}.
\end{itemize}

\subsection{I-B03 --- matrix invariants}\label{i-b03--matrix-invariants}

Forbes, M. A., \& Shpilka, A. (2013). \emph{Explicit Noether
normalization for simultaneous conjugation via polynomial identity
testing} {[}Preprint{]}. arXiv:1303.0084v2.
\url{https://arxiv.org/abs/1303.0084}

Supplied source:
\texttt{LIT\_Forbes\_Shpilka\_explicit\_Noether\_normalization\_matrix\_invariants/pit-nn.tex}.

\begin{itemize}
\tightlist
\item
  Theorem 1.3 (\texttt{generatinginvariants}): trace generators of words
  of length at most m² in characteristic zero.
\item
  Theorem 1.13 (\texttt{thm:integral\ subring}) and Corollary 1.14:
  quasipolynomial separating invariants/integral subring.
\item
  Theorem 1.15 (\texttt{thm:orbit\ intersection}): deterministic NC
  orbit-closure intersection already follows from white-box PIT.
\item
  Construction 3.1 (\texttt{roabpinv}), source lines 510--516: M(x),
  trace-generating polynomial.
\item
  Lemma 3.2 (\texttt{diffofinv}), lines 520--526: difference has width
  2m².
\item
  Lemma 3.4 (\texttt{roabpcoeff}), lines 540--567; Corollary 3.6
  (\texttt{reducetoroabp}), lines 581--587.
\item
  Construction 3.7 (\texttt{const:hittingset}), lines 591--595.
\item
  Theorem 3.8 (\texttt{thm:sepinv}), lines 599--613:
  m²\textbar H\textbar{} homogeneous separating invariants.
\item
  Theorem 3.9: specialization of Derksen--Kemper Theorem 2.3.12,
  homogeneous separating invariants give integrality over algebraically
  closed characteristic zero.
\item
  Corollary 3.10: integral-subring consequence for the same set.
\item
  Primary PDF inspected: \url{https://arxiv.org/pdf/1303.0084} ; web
  reference \texttt{turn166538694view1}; exact-number findings
  \texttt{turn944129952view2}, \texttt{turn944129952view4}.
\end{itemize}

\subsection{I-B04 --- revised GCT V
numbering}\label{i-b04--revised-gct-v-numbering}

Mulmuley, K. D. (2017). Geometric complexity theory V: Efficient
algorithms for Noether normalization. \emph{Journal of the American
Mathematical Society, 30}, 225--309.
\url{https://doi.org/10.1090/jams/864}

Supplied source:
\texttt{LIT\_Mulmuley\_GCT\_V\_Noether\_normalization/main.tex},
matching arXiv:1209.5993v5 (26 May 2016).

\begin{itemize}
\tightlist
\item
  Theorem 5.14, source label \texttt{tequisdit}, lines 3094--3110; PDF
  pp. 45--46: strengthened general PIT versus strict e.s.o.p. for
  universal-circuit orbit closures. This is stronger than ordinary ROABP
  PIT and has an infinitesimal/border requirement.
\item
  Section 7.5 and Lemma 7.9 (\texttt{lforbes}): replace symbolic
  determinant testing by ROABP testing for matrix invariants.
\item
  Theorem 7.1 (\texttt{tmatrixstrongform}), Theorem 7.2
  (\texttt{tintroorbitclosureNew}) and Theorem 7.7
  (\texttt{tsepderksen}): exact scope of matrix normalization and
  already-known orbit-closure algorithms.
\item
  Primary PDF inspected: \url{https://arxiv.org/pdf/1209.5993} ;
  \texttt{turn477504142view0}; theorem-number findings
  \texttt{turn871722820view0}, \texttt{turn871722820view1}.
\item
  A reference in the 2013 Forbes--Shpilka paper to an earlier GCT V
  Theorem 1.6 must NOT be copied as the current theorem number. In the
  supplied revised source Theorem 1.6 is a different theorem about
  constant-dimensional groups.
\end{itemize}

\subsection{Independent source check for root finding
G-07}\label{independent-source-check-for-root-finding-g-07}

Leal, M., Lozano Huerta, C., \& Vite, M. (2024). The Noether--Lefschetz
locus of surfaces in P³ formed by determinantal surfaces.
\emph{Mathematische Nachrichten, 297}(12), 4671--4688.
\url{https://doi.org/10.1002/mana.202400132}

Primary preprint: \url{https://arxiv.org/pdf/2303.09028} , version 3 (1
November 2024), inspected in this independent review.

\begin{itemize}
\tightlist
\item
  Definition 2, PDF p. 2: closed determinantal locus det(a,b).
\item
  Proposition 1.1, PDF p. 4: irreducibility and nonempty smooth locus.
\item
  Theorem 2, PDF pp. 3 and 15: prime divisor F1 degree 320112.
\item
  Table 2, PDF p. 13: F1 has a=(5,5,5,5), b=(6,6,6,6), hence linear
  entries.
\item
  Section 4.1, PDF p. 15 immediately before Theorem 2:
  P\_(20,2)=(1/2)D\_(23,14).
\item
  Primary web reference \texttt{turn280289537view0}; exact locators
  \texttt{turn520265669view0}, \texttt{turn520265669view2},
  \texttt{turn520265669view3}.
\end{itemize}

\subsection{I-B05 --- arithmetic support and logarithmic-space
counting}\label{i-b05--arithmetic-support-and-logarithmic-space-counting}

Allender, E., \& Ogihara, M. (1996). Relationships among PL, \#L, and
the determinant. \emph{RAIRO--Theoretical Informatics and Applications,
30}(1), 1--21. \url{https://doi.org/10.1051/ita/1996300100011}

Primary scan inspected:
\url{https://www.numdam.org/item/ITA_1996__30_1_1_0.pdf}

\begin{itemize}
\tightlist
\item
  Section 2, printed p. 4: halting logarithmic-space machines and their
  accepting-minus-rejecting gap.
\item
  Theorem 1, printed p. 2: GapL and logarithmic-space functional
  reductions to integer determinant; following paragraph explicitly
  passes through iterated multiplication of \{0,1,-1\} matrices.
\item
  Section 5, Definition 2, printed p. 11: C\_=L defined by gap zero.
\item
  Theorem 14, printed p. 12: singular integer matrices are
  C\_=L-complete.
\item
  I-12 proves its particular ROABP formulation directly. It does not
  assert that coC\_=L is different from NL or that arbitrary
  polynomial-time algebraic support manipulation is impossible.
\item
  Primary web references: \texttt{turn392175736view1},
  \texttt{turn776322203view2}, \texttt{turn647216830view0}.
\end{itemize}

\subsection{Independent source check for
P-08}\label{independent-source-check-for-p-08}

Lagarias, J. C., \& Odlyzko, A. M. (1977). Effective versions of the
Chebotarev density theorem. In A. Fröhlich (Ed.), \emph{Algebraic number
fields} (pp. 409--464). Academic Press.

Primary-text TeXromancers transcription inspected:
\url{https://aareyanmanzoor.github.io/assets/articles/lagarias-odlyzko.pdf}

\begin{itemize}
\tightlist
\item
  Theorem 1.1 and Corollary 1.2, transcription p. 4; equation (1.4): GRH
  for the Dedekind zeta function of L alone, an effective absolute
  constant, and an unramified prime ideal for every conjugacy class with
  norm bounded by c₂(log Δ\_L)²(log log Δ\_L)⁴.
\item
  The parenthetical immediately following (1.4) handles L=Q by p=(2).
\item
  The source\textquotesingle s full Chebotarev proof was not
  independently checked. P-08\textquotesingle s entire reduction proof
  was independently read and checked; see
  \texttt{independent\_review\_purity.md}.
\item
  Primary web reference: \texttt{turn392175736view0}.
\end{itemize}
\section{Purity and zero-statistics source
ledger}\label{purity-and-zero-statistics-source-ledger}

Run: CUE\_GCT-20261121-purity. Researcher: GPT-6 Astra Pro
(\texttt{/root/purity\_zero\_status}); configured reasoning effort is
not exposed. Sources below are background inputs or exact
statement/hypothesis checks, not claimed independent verifications of
entire literature proofs. Deductions P-01--P-11 appear in
\texttt{findings.tex}; their identifiers designate development in this
run and do not assert historical priority. Status deductions P-Z01/P-Z02
appear in \texttt{zero\_status.tex}.

\subsection{Background IDs and exact
locators}\label{background-ids-and-exact-locators}

\subsubsection{P-B01 --- Quadratic-character small
bias}\label{p-b01--quadratic-character-small-bias}

Alon, N., Goldreich, O., Håstad, J., \& Peralta, R. (1992). Simple
constructions of almost k-wise independent random variables.
\emph{Random Structures \& Algorithms, 3}(3), 289--304.
\url{https://doi.org/10.1002/rsa.3240030308}

Primary author copy: \url{https://johanhastad.se/almostindep.pdf}
(19-page copy dated September 25, 1997). Exact locator: §4, Theorem 1
(Weil character bound), Construction 2, Proposition 2 and its proof,
printed pp. 9--10 / PDF pages 9--10. P-02 uses only the monic squarefree
case, and changes the arbitrary zero-character convention to +1; the
same n/p error applies. Publication metadata checked against publisher
and author page
\url{https://www.wisdom.weizmann.ac.il/~oded/p_aghp.html}. Web refs:
turn677747127view2, turn282302316view1, turn318993354search2.

\subsubsection{P-B02 --- Curve point
bound}\label{p-b02--curve-point-bound}

Deligne, P. (1974). La conjecture de Weil. I. \emph{Publications
Mathématiques de l'IHÉS, 43}, 273--307.
\url{https://doi.org/10.1007/BF02684373}

Primary scan:
\url{https://www.numdam.org/item/PMIHES_1974__43__273_0.pdf}. Exact
locator: formula (1.5.1), p. 275, and Theorem (1.6), p. 276. Applying
the trace formula to a smooth projective geometrically connected curve,
with H⁰ of dimension 1, H¹ of dimension 2g, and H²=Q\_l(−1), yields
\textbar\#C(F\_q)−(q+1)\textbar≤2g√q. This recovers the classical curve
theorem; no historical-priority claim is made for Deligne in dimension
one. Web refs: turn677747127view0, turn282302316view0. The original
theorem statement and trace formula were read, not the entire proof of
Weil I.

\subsubsection{P-B03 --- Trace formula and
weights}\label{p-b03--trace-formula-and-weights}

Deligne, P. (1977). \emph{Cohomologie étale (SGA 4½)} (Lecture Notes in
Mathematics, Vol. 569). Springer. Exact locator: Rapport sur la formule
des traces, Theorem 3.2 (trace formula); Theorem 3.1 and Proposition 3.3
give the determinant/L-function formulation. French primary-text
transcription inspected:
\url{https://grothendiecksga.com/read/sga4.5/fr/II_3.html}. Web ref:
turn273425817view2.

Deligne, P. (1980). La conjecture de Weil. II. \emph{Publications
Mathématiques de l'IHÉS, 52}, 137--252.
\url{https://doi.org/10.1007/BF02684780}. Exact locators: Theorem 3.3.1
and §3.3.10 for the compact-support weight bound; Theorem 3.5.3 for
equidistribution. The large original PDF was not successfully opened
through this tool; these exact locators were verified in the
author-hosted Katz--Sarnak text below, §9.0.15, which states the needed
result and explicitly cites both 3.3.1 and 3.3.10. No full-proof check
of Weil II is claimed.

Katz, N. M., \& Sarnak, P. (1999). \emph{Random matrices, Frobenius
eigenvalues, and monodromy} (American Mathematical Society Colloquium
Publications, Vol. 45). American Mathematical Society. Author-hosted
book: \url{https://web.math.princeton.edu/~nmk/RMFEM.pdf}. Exact
locators: §§9.0.13--9.0.16, printed pp. 269--270 (duality, affine
vanishing, trace formula, weights, geometric coinvariants); Theorem
9.2.6 (equidistribution under the stated monodromy hypotheses);
§§9.1.5--9.1.7 for smooth proper cohomology and purity. Web refs:
turn387134677view1, turn905674939view2, turn991829417view1. These
sections/theorem were read; no claim of reading the entire book.

\subsubsection{P-B04 --- Effective Chebotarev under Dedekind
GRH}\label{p-b04--effective-chebotarev-under-dedekind-grh}

Lagarias, J. C., \& Odlyzko, A. M. (1977). Effective versions of the
Chebotarev density theorem. In A. Fröhlich (Ed.), \emph{Algebraic number
fields} (pp. 409--464). Academic Press.

Original paper in a later TeXromancers transcription:
\url{https://aareyanmanzoor.github.io/assets/articles/lagarias-odlyzko.pdf}.
Exact locator: Theorem 1.1, Corollary 1.2, formula (1.4), transcription
p. 4. Corollary 1.2 assumes GRH for the Dedekind zeta function of L and
provides an unramified prime in each conjugacy class with norm at most
c₂(log Δ\_L)²(log log Δ\_L)⁴, with an effective absolute constant. L=Q
is explicitly handled by prime 2. Bibliographic identity checked at
author publication list
\url{https://websites.umich.edu/lagarias/PUBLICATIONS/zeta.html}. Web
refs: turn677747127view1, turn858301579view1. P-08 is a new conditional
reduction using this background bound; it is not a claim that a suitable
common arithmetic encoding exists for ROABPs.

\subsubsection{P-B05 --- Montgomery supported pair
correlation}\label{p-b05--montgomery-supported-pair-correlation}

Montgomery, H. L. (1973). The pair correlation of zeros of the zeta
function. In \emph{Analytic number theory} (Proceedings of Symposia in
Pure Mathematics, Vol. 24, pp. 181--193). American Mathematical Society.

Author-hosted scan:
\url{https://public.websites.umich.edu/hlm/paircor1.pdf}. Exact locator:
§1, theorem and equations (1)--(4), pp. 181--182. RH is assumed at the
opening. The weight is 4/(4+u²); normalization is T log T/(2π); the
theorem gives the form factor uniformly on compact subintervals of
\textbar α\textbar\textless1. P-Z01 is the Fourier consequence for
compactly supported transforms in (−1,1), preserving the original weight
and ordered-index diagonal convention. The full-support conjecture is
distinguished in equations (11)--(12), pp. 183--184. Web ref:
turn991829417view2. The needed theorem, normalization, and Fourier
conversion were checked, not the complete 13-page proof.

\subsubsection{P-B06 --- Hughes--Rudnick}\label{p-b06--hughesrudnick}

Hughes, C. P., \& Rudnick, Z. (2003). Linear statistics of low-lying
zeros of L-functions. \emph{The Quarterly Journal of Mathematics,
54}(3), 309--333. \url{https://doi.org/10.1093/qmath/hag021}

Primary supplied source:
\texttt{sources/LIT\_Hughes\_Rudnick\_low-lying\_zeros\_linear\_statistics/main.tex}.
Exact locators: Introduction (complex ordinates allowed, no GRH assumed,
source lines 94--150); Definition 2.1 (admissible f, line 340); Theorem
3.1, label \texttt{thm:expectation} (line 383, support {[}−2,2{]});
Theorem 4.1, label \texttt{thm:Wf\_variance} (line 505, support
{[}−1,1{]}); Theorem 5.1, label \texttt{thm-L:Gaussian\ moments} (line
621, m\textless2/α); Theorem 8.1, label \texttt{thm:inf\_less\_1/4R}
(line 1349); Corollary 8.2, label \texttt{thm:inf\_less\_1/4} (line
1454); Theorem 8.3, label \texttt{thm:pos\_prop\_small\_zeros} (line
1484, β≥0.633). Definition x\_\{χ,j\}=γ\_\{χ,j\} log q/(2π); x\_\{χ,1\}
is the first nonnegative ordinate. Mean/variance/restricted moments are
already unconditional analytic statements; §8 explicitly assumes
Dirichlet GRH. Publication year verified as 2003, not 2004, using the
primary OUP record and author publication list. Web refs:
turn669156525search0, turn669156525search17.

\subsubsection{P-B07 --- Rudnick--Sarnak}\label{p-b07--rudnicksarnak}

Rudnick, Z., \& Sarnak, P. (1996). Zeros of principal L-functions and
random matrix theory. \emph{Duke Mathematical Journal, 81}(2), 269--322.

Primary paper copy:
\url{https://empslocal.ex.ac.uk/people/staff/mrwatkin/zeta/rudnick-sarnak.pdf}.
Exact locators: Theorems 1.1--1.2, coefficient Hypothesis (1.7), and
support condition Σ\textbar ξ\_j\textbar\textless2/m. Theorem 1.1
permits complex ordinates in its smoothed analytic statistic; Theorem
1.2 adds target RH for the real-zero/unsmoothed form. The coefficient
condition is required in the stated treatment for degree m\textgreater3.
The m=1 Dirichlet/zeta instances are the instances discharged by
stipulated Alpöge. Web ref: turn696136632view1. Statement/hypothesis
audit, not full-paper proof verification.

\subsubsection{P-B08 ---
Iwaniec--Luo--Sarnak}\label{p-b08--iwaniecluosarnak}

Iwaniec, H., Luo, W., \& Sarnak, P. (2000). Low lying zeros of families
of L-functions. \emph{Publications Mathématiques de l'IHÉS, 91},
55--131. \url{https://doi.org/10.1007/BF02698741}

Primary scan: \url{https://www.numdam.org/item/10.1007/BF02698741.pdf}.
Exact locators: §1 standing RH convention p. 56; Theorems 1.1--1.3 for
cuspidal GL₂ target families; Theorems 1.4--1.5 for symmetric-square GL₃
target families; Theorem 1.9 and Theorem 1.10/Hypothesis S (both p. 66);
§4, pp. 88--89, on auxiliary symmetric-square/Rankin--Selberg individual
prime sums and their averaged replacements. The proof does not justify
declaring all the auxiliary individual RH assumptions irreducibly
necessary: the authors explicitly remove some by averaging in the
claimed ranges. Conversely, the target cuspidal and symmetric-square RH
is not Dirichlet or one-dimensional Hecke-character GRH. Web refs:
turn167264000view0, turn651032895view0, turn651032895view1,
turn651032895view2, turn114835679view1, turn890935485view0. Statement
and dependency audit only.

\subsubsection{P-B09 --- Conditional Hecke-family
example}\label{p-b09--conditional-hecke-family-example}

Gao, P., \& Zhao, L. (2020). One-level density of low-lying zeros of
quadratic and quartic Hecke L-functions. \emph{Canadian Journal of
Mathematics, 72}(2), 427--454.
\url{https://doi.org/10.4153/S0008414X1900021X}

Primary preprint: \url{https://arxiv.org/pdf/1708.01701} (v3, March 22,
2019). Exact locators: Theorem 1.1, p. 3, quadratic Q(i) family with
support (−2,2), symplectic one-level density; Theorem 1.3, p. 3, quartic
family with support (−20/19,20/19), unitary density; Lemma 2.8, p. 9,
auxiliary nonprincipal Hecke-character prime sum under GRH. These
dependencies are covered conditionally by the user-stated Masullo Hecke
GRH, not by Dirichlet GRH alone. Web refs: turn696136632view2,
turn858301579view3, turn394727321view0.

\subsubsection{P-B10 --- Function-field moment
limitations}\label{p-b10--function-field-moment-limitations}

Sawin, W. (2020). A representation theory approach to integral moments
of L-functions over function fields. \emph{Algebra \& Number Theory,
14}(4), 867--906. \url{https://doi.org/10.2140/ant.2020.14.867}

Primary supplied source:
\texttt{sources/LIT\_representation\_theory\_approach\_to\_moments\_function\_fields/main.tex};
arXiv \url{https://arxiv.org/abs/1810.01303}. Exact locators: Hypothesis
1.3, label \texttt{hypo}, line 77; Theorem 1.4, label \texttt{main},
line 85; Remark 1.9(1),(5), lines 204--224. The theorem assumes
H(n,r,tilde r,w), n\textgreater2max(r,tilde r)+1, and characteristic ≠2
if n=4 or 5. Its error includes q\^{}((w−n)/2) C\_\{r,tilde r\}\^{}n
n\^{}(r+tilde r). Remark 1.9(5) obtains w=n−1 from duality, hence
O(q\^{}(−1/2)) with other parameters fixed. Remark 1.9(1) retains
stronger cohomology and uniform Betti-number gates for general stronger
degree-aspect/ratio/correlation conclusions. The paper also proves an
unconditional unbalanced case in Corollary 1.8 (source line 156,
repeated later at label \texttt{combined-result}, line 941), with r\_-=1
and error
sqrt(q)\emph{(q\^{}\{−(p−2r)/(2pr)\}}(r+2)\^{}\{r+1\})\^{}n*n\^{}\{r+1\},
under the stated n/characteristic conditions. This restricted
degree-aspect result must not be omitted. Publisher PDF metadata:
\url{https://msp.org/ant/2020/14-4/ant-v14-n4-p03-s.pdf}; web refs
turn404601268search12, turn404601268search0. The original source is
supplied, so no current universal claim about all later function-field
moment literature is inferred from its 2020 discussion.

\subsubsection{P-B11 --- CFKRS and CUE
distinction}\label{p-b11--cfkrs-and-cue-distinction}

Conrey, J. B., Farmer, D. W., Keating, J. P., Rubinstein, M. O., \&
Snaith, N. C. (2005). Integral moments of L-functions. \emph{Proceedings
of the London Mathematical Society, 91}(1), 33--104.
\url{https://doi.org/10.1112/S0024611504015175}

Primary supplied source:
\texttt{sources/LIT\_CFKRS\_integral\_moments\_of\_L-functions/integralmoments2v.tex}.
Exact locators: §1.5, \texttt{sec:mainconjectures}, from source line
883; Conjecture 1.5.1, \texttt{thm:zeta2kconjecture}, line 919; exact
finite-N CUE Theorem 1.5.2, \texttt{thm:UN2k}, line 968. The recipe
itself is §2.1. The rectangular/CUE agent owns detailed numbering and
the exact matrix-integral derivation. The arithmetic conjecture is not
obtained merely by discharging RH/GRH in unrelated supported
zero-statistics theorems.

\subsubsection{P-B12 --- Quantum-chaos sources; memory
qualification}\label{p-b12--quantum-chaos-sources-memory-qualification}

Bohigas, O., Giannoni, M. J., \& Schmit, C. (1984). Characterization of
chaotic quantum spectra and universality of level fluctuation laws.
\emph{Physical Review Letters, 52}, 1--4.
\url{https://doi.org/10.1103/PhysRevLett.52.1}

Berry, M. V., \& Keating, J. P. (1999). The Riemann zeros and eigenvalue
asymptotics. \emph{SIAM Review, 41}(2), 236--266.
\url{https://doi.org/10.1137/S0036144598347497}

Berry, M. V., \& Keating, J. P. (1999). H=xp and the Riemann zeros. In
J. P. Keating, D. E. Khmelnitskii, \& I. V. Lerner (Eds.),
\emph{Supersymmetry and trace formulae: Chaos and disorder} (pp.
355--367). Plenum Press. Author-institution bibliographic record:
\url{https://research-information.bris.ac.uk/en/publications/ih-xpi-and-the-riemann-zeros/}.

The symmetry-class connection was recalled from memory and is explicitly
marked that way in the paper. Publication locators and the SIAM
abstract\textquotesingle s H\_cl=XP speculation were verified against
primary publication/author records (turn669156525search2,
turn669156525search4; APS DOI above). No uninspected chapter or equation
number is claimed. The stipulated statements do not construct a
specified differential or pseudodifferential Hamiltonian with the
desired arithmetic spectral interpretation. The paper explicitly
distinguishes this from the elementary abstract diagonal self-adjoint
operator obtained by inserting the already-given real zero ordinates
into its definition.

\subsubsection{P-B13 --- User-stipulated inputs and limits of the
attachments}\label{p-b13--user-stipulated-inputs-and-limits-of-the-attachments}

Alpöge: \texttt{sources/ALPOGE\_Main.lean} and
\texttt{sources/ALPOGE\_Statements.lean}. Read the actual exported
declarations \texttt{riemannHypothesis},
\texttt{riemannZeta\_eq\_zero\_re},
\texttt{LFunction\_ne\_zero\_of\_half\_lt\_re},
\texttt{LFunction\_eq\_zero\_re}, and
\texttt{LFunction\_eq\_zero\_re\_of\_isPrimitive}, and their statement
predicates. These are used as stipulated theorem input; imports were not
rebuilt, and this ledger is not a truth audit.

Masullo: \texttt{sources/MASULLO\_STATEMENTS\_INDEX.md} only. The prompt
attributes Hecke GRH to §1.2; this section\textquotesingle s numbered
statement/full proof is not present in the supplied index. Theorem
13.3.1.1, indexed PDF p. 522, is a trace-formula claim; its following
malformed row contains editorial placeholders. Conjecture 14.1.2.1, p.
529, concerns general automorphic GL\_n GRH/GSH and is explicitly a
conjecture. Proposition 14.1.2.3, p. 531, assumes Conjectures 14.1.2.1
and 14.1.2.2. These entries do not validate a general automorphic
theorem. All uses of M remain conditional on the user-stated Hecke
claim, with no invented theorem number or proof check.

\subsubsection{P-B14 --- Explicit conductor
budget}\label{p-b14--explicit-conductor-budget}

Grothendieck, A. (1977). Formule d'Euler--Poincaré en cohomologie étale
(I. Bucur, rédaction). In \emph{Cohomologie l-adique et fonctions L (SGA
5)} (Lecture Notes in Mathematics, Vol. 589), Exposé X. Springer. Exact
locators: Theorem 7.1, Corollary 7.12, formula (7.13). Primary-text
English transcription inspected:
\url{https://grothendiecksga.com/read/sga5/en/X-7.html}
(turn343029532view0); its French counterpart did not load in this tool.
The l-adic formula used in P-06 is explicitly stated in the
author-hosted Katz--Sarnak book, proof of Lemma 10.1.12, printed p. 296
(turn387134677view0). That formula plus h\_c\^{}0=0 and h\_c\^{}2=δ\_ij
gives the exact conductor expression; the whole SGA proof is not claimed
independently verified.

\subsubsection{P-B15 --- Rank-one measurements and energy
distortion}\label{p-b15--rank-one-measurements-and-energy-distortion}

Cai, T. T., \& Zhang, A. (2015). ROP: Matrix recovery via rank-one
projections. \emph{The Annals of Statistics, 43}(1), 102--138.
\url{https://doi.org/10.1214/14-AOS1267}

Primary author-hosted article:
\url{https://tony-cai.com/assets/pdf/papers/ROP.pdf}. Exact locators:
Definition 2.1, Theorems 2.1--2.2, p. 107, for the ℓ₁ measurement
condition; §2.2, pp. 108--110; Lemma 2.1, p. 109, for failure of
constant-distortion ℓ₂ restricted isometry with subambient sampling in
the Gaussian rank-one matrix model. The lemma\textquotesingle s proof is
assigned by the paper to supplementary material, which was not read.
This is a relevant predecessor, not an attribution of the exact weighted
tensor statement P-09 to the paper. P-09 has its own complete
deterministic trace proof. Web refs: turn851314685view0,
turn170484496view1, turn637901333view2.

\subsubsection{P-B16 --- Width-one kernels as completely entangled
subspaces}\label{p-b16--width-one-kernels-as-completely-entangled-subspaces}

Parthasarathy, K. R. (2004). On the maximal dimension of a completely
entangled subspace for finite level quantum systems. \emph{Proceedings
of the Indian Academy of Sciences---Mathematical Sciences, 114}(4),
365--374. \url{https://doi.org/10.1007/BF02829441}

Primary author preprint: \url{https://arxiv.org/pdf/quant-ph/0405077}.
Exact locators: Proposition 1.3, Proposition 1.4, and Theorem 1.5,
preprint pp. 4--6. Read both complete proposition proofs and the theorem
deduction. The Vandermonde product-vector construction is the diagonal
polynomial-evaluation construction in P-10; the optimal codimension is
Σ\_i(d\_i−1)+1. The elementary arbitrary-node annihilator and the
tensor-train kernel/conditioning formulation are written out in P-10.
Publication metadata corroborated by
\url{https://repository.ias.ac.in/36155/}. Web refs: turn637901333view0,
turn698733368view0, turn444985646search1. The later sections on
perfectly entangled subspaces were not reviewed.

\subsubsection{P-B17 --- Tensor-train ranks and closed coefficient
classes}\label{p-b17--tensor-train-ranks-and-closed-coefficient-classes}

Oseledets, I. V. (2011). Tensor-train decomposition. \emph{SIAM Journal
on Scientific Computing, 33}(5), 2295--2317.
\url{https://doi.org/10.1137/090752286}

Primary published article copy:
\url{https://users.math.msu.edu/users/iwenmark/Teaching/CMSE890/TENSOR_oseledets2011.pdf}.
Exact locators: §2, definition and equation (2.1); Theorem 2.1 and
complete rank-factorization proof, pp. 2298--2299. These were read.
Corollary 2.4, p. 2300, also records closedness/existence for bounded TT
ranks; its proof was read. P-10 uses only exact rank factorization,
homogeneous minor equations, compactness, and the elementary parameter
upper bound. No numerical TT-SVD algorithm was implemented or certified.
Web refs: turn637901333view1, turn698733368view1, turn698733368view2;
publisher metadata turn444985646search2.

\subsection{Scope of citation checks}\label{scope-of-citation-checks}

The bibliography distinguishes primary text, author-hosted copies,
publisher metadata, later transcriptions, supplied source files, and
memory-based explanations. Numerical theorem labels derived from
supplied LaTeX use the source\textquotesingle s shared theorem counters;
integration should preserve the source label as a second locator. This
ledger does not certify source completeness, detect unseen subsequent
literature, or claim formal verification.
\section{Integral geometry and archive source
register}\label{integral-geometry-and-archive-source-register}

These are framework and source-dependency checks. They do not claim full
independent verification of every cited literature proof. New finite
constructions and transfers are proved under their finding IDs in the
manuscript.

\subsection{G-B01 --- Grothendieck\textquotesingle s supplied cote 133
transcriptions}\label{g-b01--grothendiecks-supplied-cote-133-transcriptions}

Grothendieck, A. (n.d.). \emph{Cote 133, scans 006, 008, 009, 019, and
028} {[}Supplied LaTeX transcriptions{]}.

Exact files: ARCHIVE\_133\_scan006.tex, ARCHIVE\_133\_scan008.tex,
ARCHIVE\_133\_scan009.tex, ARCHIVE\_133\_scan019.tex,
ARCHIVE\_133\_scan028.tex. All five transcriptions were read. No
original image is supplied. See TRANSCRIPTION\_NOTES.tsv for file-line
locators and separate mathematical and paleographic confidence. The
accompanying SA17\_STUDY\_cote133\_lambda.md is an earlier research
record, not an independently reverified image reading.

\subsection{G-B02 --- SGA 3}\label{g-b02--sga-3}

Demazure, M., \& Grothendieck, A. (Eds.). (1970). \emph{Schémas en
groupes (SGA 3)} (Lecture Notes in Mathematics, Vols. 151--153).
Springer.

Exact locators: Exposé I, Propositions 4.7.2--4.7.4 (comodules,
diagonalizable groups and gradings, splitting), Theorem 5.3.3
(diagonalizable-group cohomology over an affine base). Original Exposé I
inspected at
\url{https://webusers.imj-prg.fr/~patrick.polo/SGA3/origExp1.pdf} .
Exposé VI\_B, Proposition 4.1 (fiber dimension for group schemes),
electronic primary-text transcription
\url{https://grothendiecksga.com/read/sga3/fr/VIB-4.html} . It
explicitly cites EGA IV\_3, Theorem 13.1.3. The entire EGA proof was not
reread here.

\subsection{G-B03 --- SGA 6}\label{g-b03--sga-6}

Berthelot, P., Grothendieck, A., \& Illusie, L. (Eds.). (1971).
\emph{Théorie des intersections et théorème de Riemann--Roch (SGA 6)}
(Lecture Notes in Mathematics, Vol. 225). Springer.
\url{https://doi.org/10.1007/BFb0066283}

Exact locators: Exposé V, Definition 2.1, Example 2.2(b), Lemma 2.2.1,
Definition 2.4 and equation (2.4.2), Proposition 3.2, §3.10 equation
(3.10.1), §7.1 equation (7.1.1), Proposition 7.5; Exposé VI, Corollary
3.2. The V passages were read in
\url{https://grothendiecksga.com/read/sga6/fr/V-2.html} ,
\url{https://grothendiecksga.com/read/sga6/fr/V-3.html} and
\url{https://grothendiecksga.com/read/sga6/fr/V-7.html} . The VI 3.2
locator was located/corroborated, not claimed as a complete proof
reading. A separate 24-page PDF found during retrieval was an incomplete
front-matter excerpt and was not treated as the full volume.

\subsection{G-B04 --- Integral representation
rings}\label{g-b04--integral-representation-rings}

Serre, J.-P. (1968). Groupe de Grothendieck des schémas en groupes
réductifs déployés. \emph{Publications Mathématiques de
l\textquotesingle IHÉS, 34}, 37--52.
\url{https://doi.org/10.1007/BF02684589}

Exact locators: Theorem 5, §3.7; §3.8 for GL\_n. Author-hosted original
scan inspected:
\url{https://www.college-de-france.fr/media/jean-pierre-serre/UPL8576975072934521334_Serre_Grpes_de_Gr.pdf}
. The theorem concerns Grothendieck groups for split reductive group
schemes; it is not a claim that characteristic-zero simple
decompositions remain semisimple in positive characteristic, or a
solution of the distinct finite-S\_n lattice question.

\subsection{G-B05 --- Hilbert and Quot
schemes}\label{g-b05--hilbert-and-quot-schemes}

Grothendieck, A. (1961). Techniques de construction et théorèmes
d\textquotesingle existence en géométrie algébrique IV: Les schémas de
Hilbert. \emph{Séminaire Bourbaki, 6}, Exposé 221, 249--276.
\url{https://www.numdam.org/article/SB_1960-1961__6__249_0.pdf}

Exact locators: Theorem 3.1 (Quot representability), Theorem 3.2 (fixed
Hilbert polynomial), Proposition 3.8 (Grassmannian immersion). Relevant
original pages inspected. These are parameter-space theorems, not
efficient algorithms for the present rank computations.

\subsection{G-B06 --- EGA graded modules and
dimension}\label{g-b06--ega-graded-modules-and-dimension}

Grothendieck, A., \& Dieudonné, J. (1961). Éléments de géométrie
algébrique II: Étude globale élémentaire de quelques classes de
morphismes. \emph{Publications Mathématiques de l\textquotesingle IHÉS,
8}, 5--222.

Exact locator: §§2.4--2.5, Proposition 2.5.1 for the sheaf associated to
a graded module. Electronic translation inspected at
\url{https://jcreinhold.github.io/ega/ii/02-02-homogeneous-prime-spectra.html}
. The original large PDF did not load in the web tool. EGA IV\_3,
Theorem 13.1.3 is used through its explicit citation and application in
SGA 3 VI\_B Proposition 4.1, with this check scope recorded.

\subsection{G-B07 --- Determinantal plane curves and cubic
surfaces}\label{g-b07--determinantal-plane-curves-and-cubic-surfaces}

Beauville, A. (2000). Determinantal hypersurfaces. \emph{Michigan
Mathematical Journal, 48}, 39--64.
\url{https://doi.org/10.1307/mmj/1030132717}

Author PDF: \url{https://math.univ-cotedazur.fr/u/beauvill/pubs/det.pdf}
. Exact locators: Proposition 3.1(a), §3.2, Proposition 3.3 and Remark
3.4 for plane curves and their line bundles; Corollary 6.4 for smooth
cubic surfaces; Corollary 1.12 for the relevant linear resolution. These
statements were inspected. G05 supplies the required
density/highest-weight deduction; R11 supplies the low-degree padding
deduction.

\subsection{G-B08 --- The reduced determinantal-quartic
divisor}\label{g-b08--the-reduced-determinantal-quartic-divisor}

Leal, M., Lozano Huerta, C., \& Vite, M. (2024). Noether--Lefschetz loci
of determinantal surfaces. \emph{Mathematische Nachrichten, 297}(12),
4671--4688. \url{https://doi.org/10.1002/mana.202400132}

Primary version inspected: arXiv:2303.09028v3,
\url{https://arxiv.org/pdf/2303.09028} . Exact locators: Theorem 2;
Table 2, component F\_1 with a=(5,5,5,5), b=(6,6,6,6), degree 320112;
§4.1 immediately preceding Theorem 2, P\_(20,2)=(1/2)D\_(23,14).
Original table and theorem were visually checked. The attached Landsberg
survey\textquotesingle s section zhanglee/theorem zlthm gives 640224;
the reduced divisor degree required for a principal equation is 320112.
G07 records this cycle-multiplicity correction, separately from
manuscript transcription.

\subsection{G-B09 --- Known optimal order for uniform determinantal
representations}\label{g-b09--known-optimal-order-for-uniform-determinantal-representations}

Boralevi, A., van Doornmalen, J., Draisma, J., Hochstenbach, M. E., \&
Plestenjak, B. (2017). Uniform determinantal representations. \emph{SIAM
Journal on Applied Algebra and Geometry, 1}(1), 415--441.
\url{https://doi.org/10.1137/16M1085656}

Primary author copy:
\url{https://hochsten.win.tue.nl/pdf/siag-unidetrep.pdf} . Exact
locators: Theorem 1.3, printed p.417, and its construction in §4,
printed pp.425--427; Proposition 4.5; Remark 4.8. The asymptotic order
d\^{}(v/2) for v fixed affine variables is a known theorem. G09 gives a
displayed finite shared-chain bound and the low-row multiplicity
consequence, without claiming the asymptotic exponent as new. Both
independent reviewers checked the stated scope and the
construction\textquotesingle s coordinate-chart qualification.

\subsection{G-B10 --- Supplied analytic
inputs}\label{g-b10--supplied-analytic-inputs}

Alpöge, L. (2026, November 20). \emph{Remark} {[}User-stipulated theorem
and supplied note{]}. \url{https://alpo.ge/remark.pdf}

Exact locators: steps 2--5 p.3; steps 7,8,13,17 p.4; exported theorem
names and predicates in ALPOGE\_Main.lean and ALPOGE\_Statements.lean.
The source files and four-page note were read; the imported Lean archive
was not supplied or rebuilt.

Masullo, A. M. (2026, November 20). \emph{The music of the primes}
{[}Draft; supplied statement index only{]}.

The general Hecke-GRH claim is the prompt\textquotesingle s §1.2
attribution; its numbered statement is absent from the supplied index.
Theorem13.3.1.1 p.522 and Conjecture14.1.2.1 p.529 were inspected in the
index. Every dependency on the general Hecke claim is conditional.

\end{document}
